\documentclass[11pt,a4paper]{article}

% ── Packages ──────────────────────────────────────────────────────
\usepackage[utf8]{inputenc}
\usepackage[T1]{fontenc}
\usepackage{amsmath,amsthm,amssymb,amsfonts}
\usepackage{mathtools}
\usepackage[ruled,vlined,linesnumbered]{algorithm2e}
\usepackage{hyperref}
\usepackage{cleveref}
\usepackage{booktabs}
\usepackage{enumitem}
\usepackage{graphicx}
\usepackage[margin=1in]{geometry}
\usepackage{xcolor}
\usepackage{thmtools}

% ── Theorem Environments ──────────────────────────────────────────
\newtheorem{theorem}{Theorem}[section]
\newtheorem{lemma}[theorem]{Lemma}
\newtheorem{proposition}[theorem]{Proposition}
\newtheorem{corollary}[theorem]{Corollary}
\theoremstyle{definition}
\newtheorem{definition}[theorem]{Definition}
\newtheorem{assumption}{Assumption}
\newtheorem{example}[theorem]{Example}
\theoremstyle{remark}
\newtheorem{remark}[theorem]{Remark}
\newtheorem{notation}[theorem]{Notation}

% ── Macros ────────────────────────────────────────────────────────
\newcommand{\CI}{\perp\!\!\!\perp}
\newcommand{\NCI}{\not\!\perp\!\!\!\perp}
\newcommand{\DAG}{\mathcal{G}}
\newcommand{\Pa}{\mathrm{Pa}}
\newcommand{\An}{\mathrm{An}}
\newcommand{\De}{\mathrm{De}}
\newcommand{\Nd}{\mathrm{Nd}}
\newcommand{\adj}{\mathrm{adj}}
\newcommand{\dsep}{\mathrm{d\text{-}sep}}
\newcommand{\editdist}{d}
\newcommand{\fragility}{F}
\newcommand{\radius}{\rho}
\newcommand{\Nk}{\mathcal{N}_k}
\newcommand{\MEDR}{\textsc{MEDR}}
\newcommand{\Phi}{\varPhi}
\newcommand{\R}{\mathbb{R}}
\newcommand{\N}{\mathbb{N}}
\newcommand{\E}{\mathbb{E}}
\DeclareMathOperator*{\argmin}{arg\,min}
\DeclareMathOperator*{\argmax}{arg\,max}
\DeclareMathOperator{\poly}{poly}

\title{\textbf{CausalCert: Per-Edge Fragility Scores and Exact Robustness Radii\\for Causal DAGs via ILP-Backed Structural Stress Testing}}
\author{Theory Monograph}
\date{\today}

\begin{document}
\maketitle

% ══════════════════════════════════════════════════════════════════
% ABSTRACT
% ══════════════════════════════════════════════════════════════════
\begin{abstract}
Every observational causal analysis rests on a directed acyclic graph (DAG) encoding
structural assumptions. When a reviewer asks ``what if your DAG is wrong?,'' the analyst
has no principled quantitative answer. We introduce \textsc{CausalCert}, a structural
stress-testing framework that computes two complementary diagnostics:
(1)~\emph{per-edge fragility scores} quantifying each edge's contribution to maintaining
the causal conclusion, and (2)~the \emph{structural robustness radius}---the minimum
number of simultaneous edge additions, deletions, or reversals that would overturn the
conclusion. We prove that computing the robustness radius (Minimum-Edit DAG Repair,
$\MEDR$) is NP-complete via reduction from Minimum Vertex Cover, but show it is
fixed-parameter tractable in the treewidth~$w$ of the moral graph with running time
$O(2^{O(w^2)} \cdot p^{O(w)})$. We provide an ILP formulation with LP relaxation whose
integrality gap is bounded by $O(w)$. On 15 published DAGs from epidemiology,
economics, biology, and social science, we predict that $\geq 10/15$ have robustness
radius at most~2, and per-edge fragility scores identify load-bearing edges with AUC
$> 0.85$ on semi-synthetic benchmarks. \textsc{CausalCert} transforms the qualitative
debate about DAG correctness into a quantitative stress test.
\end{abstract}

\tableofcontents
\newpage


% ══════════════════════════════════════════════════════════════════
% SECTION 1: INTRODUCTION
% ══════════════════════════════════════════════════════════════════
\section{Introduction}\label{sec:intro}

Causal inference from observational data requires encoding domain knowledge as a
directed acyclic graph (DAG) \cite{pearl2009,spirtes2000}. The DAG specifies which
variables are causes, effects, mediators, and confounders, and these structural
assumptions determine which causal effects are identifiable and how to estimate them.
Yet the DAG itself is an assumption---and, in practice, a fragile one.

When an analyst presents a causal analysis based on a specific DAG, the first question
from a critical reviewer is: \emph{``What if your DAG is wrong?''} Currently, the
analyst has no principled quantitative answer. Sensitivity analysis tools such as
E-values \cite{vanderweele2017} and \texttt{sensemakr} \cite{cinelli2020} address
robustness to \emph{unmeasured confounding}---they quantify how strong a hidden
confounder would need to be to overturn the conclusion---but they hold the DAG structure
fixed. No analogous tool exists for \emph{structural robustness}: how many edges in the
DAG would need to be wrong (added, deleted, or reversed) before the causal conclusion
is overturned?

\subsection{Contributions}

We introduce \textsc{CausalCert}, a structural stress-testing framework that fills this
gap. Our contributions are:

\begin{enumerate}[label=(\roman*)]
\item \textbf{Per-edge fragility scores.} For each edge or non-edge in the ancestral
  subgraph, we compute a fragility score $\fragility(e) \in [0,1]$ measuring how
  sensitive the causal conclusion is to misspecification of that edge
  (\Cref{def:fragility}).

\item \textbf{Structural robustness radius.} We define the robustness radius
  $\radius(\DAG, X, Y, \Phi)$ as the minimum number of simultaneous structural edits
  (additions, deletions, reversals) needed to overturn the conclusion predicate~$\Phi$
  (\Cref{def:radius}).

\item \textbf{Complexity trifecta.} We prove:
  \begin{itemize}
    \item $\MEDR$ is NP-complete via reduction from Minimum Vertex Cover (\Cref{thm:np_hard}).
    \item $\MEDR$ is fixed-parameter tractable in the treewidth~$w$ of the moral graph,
          with running time $O(2^{O(w^2)} \cdot |\An|^{O(w)})$ (\Cref{thm:fpt}).
    \item The LP relaxation of our ILP formulation has integrality gap at most
          $2w + 1 = O(w)$ (\Cref{thm:lp_gap}).
  \end{itemize}

\item \textbf{ILP-backed exact computation.} We provide an integer linear programming
  formulation that computes exact robustness radii for DAGs with up to $p \approx 20$
  nodes (\Cref{alg:ilp}).

\item \textbf{Evaluation design.} We define 17 falsifiable predictions on synthetic,
  semi-synthetic, and published DAG benchmarks, with 6 baselines (\Cref{sec:eval}).
\end{enumerate}

\subsection{Motivating Example}

Consider a classic epidemiological DAG with $p = 12$ nodes modeling the effect of
smoking~($X$) on lung cancer~($Y$), with confounders including socioeconomic status,
age, and genetic factors. A back-door adjustment for the set $\{$age, SES$\}$ identifies
the causal effect. But suppose the edge SES~$\to$~smoking is reversed (smoking causes
low SES rather than vice versa). This single reversal invalidates the back-door
criterion---the adjusted set no longer blocks all back-door paths. The robustness radius
is~1: the conclusion is \emph{structurally fragile}.

In contrast, a different DAG for the effect of education on income might have robustness
radius~3: three simultaneous misspecifications would be required to overturn the
conclusion. \textsc{CausalCert} makes this distinction precise and computable.

\subsection{Paper Organization}

\Cref{sec:prelim} introduces definitions and notation.
\Cref{sec:assumptions} states our formal assumptions.
\Cref{sec:results} presents the main theoretical results.
\Cref{sec:algorithms} describes all algorithms with pseudocode.
\Cref{sec:complexity} summarizes complexity bounds.
\Cref{sec:eval} details the evaluation design.
\Cref{sec:related} surveys related work.
\Cref{sec:limitations} discusses limitations.


% ══════════════════════════════════════════════════════════════════
% SECTION 2: PRELIMINARIES AND DEFINITIONS
% ══════════════════════════════════════════════════════════════════
\section{Preliminaries and Definitions}\label{sec:prelim}

We establish the formal framework for structural robustness analysis of causal DAGs.
Throughout, $V = \{V_1, \ldots, V_p\}$ denotes a set of random variables, $X \in V$ is
the treatment, $Y \in V$ is the outcome, and $P$ is the joint distribution over~$V$.

\begin{definition}[Causal DAG]\label{def:dag}
A \emph{causal DAG} is a pair $\DAG = (V, E)$ where $V = \{V_1, \ldots, V_p\}$ is a
set of random variables and $E \subseteq V \times V$ is a set of directed edges such
that $\DAG$ is acyclic. The distribution $P$ over $V$ satisfies the \emph{Markov
property} with respect to~$\DAG$: every variable $V_i$ is conditionally independent of
its non-descendants $\Nd_\DAG(V_i)$ given its parents $\Pa_\DAG(V_i)$:
\[
  V_i \CI \Nd_\DAG(V_i) \mid \Pa_\DAG(V_i) \quad \text{for all } V_i \in V.
\]
\end{definition}

\begin{definition}[Structural Edit and Edit Distance]\label{def:edit}
A \emph{structural edit} on DAG $\DAG = (V, E)$ is one of three operations:
\begin{enumerate}[label=(\alph*)]
  \item \textbf{Edge addition:} add $(i, j) \notin E$ to $E$, provided the result is
    acyclic.
  \item \textbf{Edge deletion:} remove $(i, j) \in E$ from $E$.
  \item \textbf{Edge reversal:} replace $(i, j) \in E$ with $(j, i)$, provided the
    result is acyclic.
\end{enumerate}
The \emph{edit distance} $\editdist(\DAG, \DAG')$ between two DAGs on the same vertex
set $V$ is the minimum number of structural edits needed to transform $\DAG$ into
$\DAG'$. We write $\DAG \oplus e$ for the DAG obtained by applying edit~$e$ to~$\DAG$.
\end{definition}

\begin{definition}[Conclusion Predicate]\label{def:predicate}
A \emph{conclusion predicate} $\Phi(\DAG, X, Y)$ is a Boolean function that evaluates
whether a causal conclusion holds in DAG~$\DAG$ for treatment~$X$ and outcome~$Y$. We
consider three predicates:
\begin{enumerate}[label=(\roman*)]
  \item \textbf{Back-door identifiability:} $\Phi_{\mathrm{BD}}(\DAG, X, Y) = 1$ iff
    there exists a set $S \subseteq V \setminus \{X, Y\}$ satisfying the back-door
    criterion relative to $(X, Y)$ in $\DAG$.
  \item \textbf{Sign preservation:} $\Phi_{\mathrm{sign}}(\DAG, X, Y, D) = 1$ iff
    the back-door-adjusted effect has the same sign as in the original DAG.
  \item \textbf{Bound preservation:} $\Phi_{\mathrm{bound}}(\DAG, X, Y, D, \varepsilon) = 1$
    iff the back-door-adjusted effect is within $\varepsilon$ of the original estimate.
\end{enumerate}
\end{definition}

\begin{definition}[Per-Edge Fragility Score]\label{def:fragility}
For a DAG $\DAG$, treatment $X$, outcome $Y$, conclusion predicate $\Phi$, and a
candidate edit $e$, the \emph{fragility score} is
\[
  \fragility(e; \DAG, X, Y, \Phi) = \max\bigl(\fragility_{\dsep}(e),\;
  \fragility_{\mathrm{id}}(e),\; \fragility_{\mathrm{est}}(e)\bigr),
\]
where:
\begin{itemize}
  \item $\fragility_{\dsep}(e) = \max_{S \in \mathrm{ADJ}(\DAG,X,Y)}
    \mathbf{1}\bigl[\dsep_\DAG(X,Y \mid S) \neq \dsep_{\DAG \oplus e}(X,Y \mid S)\bigr]$
    measures whether the edit changes d-separation along any valid adjustment path.
  \item $\fragility_{\mathrm{id}}(e) = \mathbf{1}\bigl[\Phi(\DAG,X,Y)
    \neq \Phi(\DAG \oplus e, X, Y)\bigr]$ measures whether the edit changes
    identifiability.
  \item $\fragility_{\mathrm{est}}(e) = \min\!\Bigl(1,\;\frac{|\hat\tau(\DAG) - \hat\tau(\DAG \oplus e)|}
    {\max(|\hat\tau(\DAG)|, \varepsilon_0)}\Bigr)$ measures relative change in the estimated
    treatment effect, where $\hat\tau(\DAG)$ is the AIPW estimate under adjustment sets
    valid for $\DAG$ and $\varepsilon_0 > 0$ is a stability constant.
\end{itemize}
\end{definition}

\begin{definition}[Structural Robustness Radius]\label{def:radius}
The \emph{structural robustness radius} of DAG $\DAG$ with respect to treatment $X$,
outcome $Y$, and conclusion predicate $\Phi$ is
\[
  \radius(\DAG, X, Y, \Phi) = \min\bigl\{\editdist(\DAG, \DAG') :
  \DAG' \text{ is a DAG on } V \text{ and } \neg\Phi(\DAG', X, Y)\bigr\}.
\]
If $\neg\Phi(\DAG, X, Y)$ already, then $\radius = 0$. If $\Phi(\DAG', X, Y)$ for all
DAGs $\DAG'$ on $V$, then $\radius = \infty$ (the conclusion is structurally
unfalsifiable).
\end{definition}

\begin{definition}[$k$-Edit Neighborhood]\label{def:neighborhood}
The \emph{$k$-edit neighborhood} of DAG $\DAG$ is
\[
  \Nk(\DAG) = \bigl\{\DAG' : \DAG' \text{ is a DAG on } V \text{ and }
  \editdist(\DAG, \DAG') \leq k\bigr\}.
\]
The robustness radius satisfies
$\radius = \min\{k : \exists\, \DAG' \in \Nk(\DAG) \text{ with }
\neg\Phi(\DAG', X, Y)\}$.
\end{definition}

\begin{definition}[Minimum-Edit DAG Repair ($\MEDR$)]\label{def:medr}
$\MEDR$ is the following decision problem: Given a DAG $\DAG = (V, E)$, a set of
CI constraints $\mathcal{C} = \{(A_i \CI B_i \mid S_i)\}_{i=1}^m$ derived from data,
a conclusion predicate $\Phi$, and an integer $k$, does there exist a DAG $\DAG'$ with
$\editdist(\DAG, \DAG') \leq k$ such that $\DAG'$ satisfies all constraints in
$\mathcal{C}$ and $\neg\Phi(\DAG', X, Y)$?

The \emph{optimization version} asks for the minimum such $k$, which equals the
robustness radius $\radius$ when $\mathcal{C}$ encodes the observed CI relations.
\end{definition}

\begin{definition}[CI-Consistency]\label{def:ci_consistency}
A DAG $\DAG$ is \emph{CI-consistent} with a distribution $P$ (or a set of CI test
results) if for every conditional independence $(A \CI B \mid S)$ that holds in $P$ (or
is not rejected by the CI test), $A$ and $B$ are d-separated given $S$ in $\DAG$. A DAG
is \emph{CI-faithful} if additionally every d-separation in $\DAG$ corresponds to a CI
relation in~$P$.
\end{definition}

\begin{definition}[Ancestral Set]\label{def:ancestral}
The \emph{ancestral set} of $\{X, Y\}$ in DAG $\DAG$ is
\[
  \An_\DAG(X, Y) = \{V_i \in V : V_i \text{ is an ancestor of } X \text{ or } Y
  \text{ in } \DAG\} \cup \{X, Y\}.
\]
Equivalently, $\An_\DAG(X, Y)$ is the set of all nodes reachable from $\{X, Y\}$ by
following edges in reverse direction, plus $X$ and $Y$ themselves.
\end{definition}

\begin{definition}[Moral Graph]\label{def:moral}
The \emph{moral graph} of a DAG $\DAG$ restricted to a vertex set $S \subseteq V$ is
the undirected graph $M(\DAG[S])$ obtained by:
\begin{enumerate}[label=(\alph*)]
  \item restricting $\DAG$ to vertices in $S$,
  \item adding an undirected edge between every pair of vertices that share a common
    child in $\DAG[S]$ (marrying the parents), and
  \item replacing all directed edges with undirected edges.
\end{enumerate}
\end{definition}

\begin{definition}[Treewidth]\label{def:treewidth}
The \emph{treewidth} of an undirected graph $H = (V_H, E_H)$ is the minimum width of
any tree decomposition of $H$. A \emph{tree decomposition} is a pair
$(T, \{B_t\}_{t \in V(T)})$ where $T$ is a tree and each $B_t \subseteq V_H$ (called
a \emph{bag}) such that:
\begin{enumerate}[label=(\alph*)]
  \item every vertex of $H$ appears in at least one bag,
  \item for every edge $\{u, v\} \in E_H$, there exists a bag containing both $u$ and $v$,
  \item for every vertex $v \in V_H$, the bags containing $v$ form a connected subtree
    of $T$.
\end{enumerate}
The \emph{width} is $\max_t |B_t| - 1$.
\end{definition}

\begin{remark}[Practical Treewidth Values]\label{rem:treewidth}
Published causal DAGs from epidemiology and social science typically have moral graph
treewidth $w \leq 4$. For example, the Sachs et~al.\ protein signaling network
\cite{sachs2005} has $p = 11$ nodes and $w = 3$; the Coronary Heart Disease DAG of
\cite{lauritzen1988} has $p = 8$ and $w = 2$. This makes the FPT algorithm
(\Cref{thm:fpt}) practical for the vast majority of published DAGs.
\end{remark}

\begin{remark}[Edge Cases]\label{rem:edge_cases}
Several degenerate cases deserve explicit treatment:
\begin{enumerate}[label=(\alph*)]
  \item \textbf{$\An_\DAG(X,Y) = V$.} When the ancestral set equals the full vertex
    set (e.g., $X$ is a root and $Y$ is a sink in a connected DAG), \Cref{thm:locality}
    and \Cref{thm:ancestral_suff} are vacuously true and provide no computational
    savings. In this regime, the treewidth $w$ (\Cref{thm:fpt}) is the sole source of
    tractability.
  \item \textbf{Disconnected DAGs.} If $X$ and $Y$ are in different connected
    components, then $\An_\DAG(X,Y) = \{X, Y\}$, there are no back-door paths,
    $\Phi_{\mathrm{BD}}(\DAG, X, Y) = 1$ trivially, and $\radius \geq 1$ (at least
    one edge must be added to create a connection that could invalidate $\Phi$).
  \item \textbf{$X \to Y \in E$ (adjacent treatment-outcome).} Deleting the direct
    edge does not invalidate $\Phi_{\mathrm{BD}}$ (back-door paths are unaffected) but
    may invalidate $\Phi_{\mathrm{sign}}$ if $X \to Y$ is the only causal path.
    Reversing $X \to Y$ to $Y \to X$ fundamentally changes the causal interpretation.
    Users should interpret fragility scores for the direct $X \to Y$ edge with caution.
\end{enumerate}
\end{remark}


% ══════════════════════════════════════════════════════════════════
% SECTION 3: ASSUMPTIONS
% ══════════════════════════════════════════════════════════════════
\section{Assumptions}\label{sec:assumptions}

We state six assumptions underlying CausalCert. Assumptions~A1--A3 are standard in
causal inference; A4--A6 are specific to our framework and explicitly acknowledged as
limitations.

\begin{assumption}[Causal Markov Condition]\label{asm:markov}
The joint distribution $P(V)$ is Markov with respect to the true (unknown) causal DAG
$\DAG^*$: every variable $V_i$ is conditionally independent of its non-descendants
given its parents $\Pa_{\DAG^*}(V_i)$.
\end{assumption}

This is the foundational assumption of DAG-based causal inference
\cite{pearl2009,spirtes2000}. Without it, conditional independence tests provide no
information about graph structure. The Markov condition is universally assumed; its
violation would undermine \emph{all} DAG-based methods, not just CausalCert.

\begin{assumption}[Faithfulness]\label{asm:faithful}
The distribution $P$ is faithful to $\DAG^*$: every conditional independence in $P$
corresponds to a d-separation in $\DAG^*$. Equivalently, there are no ``accidental''
conditional independences arising from exact parameter cancellation.
\end{assumption}

Faithfulness violations are measure-zero in the parameter space \cite{meek1995} but can
occur in practice when edge weights approximately cancel. The CI ensemble
(\Cref{alg:ci_ensemble}) mitigates this by testing at multiple resolutions. We flag
faithfulness as a known limitation (\Cref{sec:limitations}).

\begin{assumption}[Causal Sufficiency]\label{asm:sufficiency}
All common causes of any two measured variables are themselves measured. There are no
latent confounders.
\end{assumption}

This justifies working with DAGs rather than MAGs or PAGs. CausalCert is explicitly
complementary to unmeasured-confounding methods (E-values, sensemakr): it addresses a
\emph{different} axis of robustness. Extension to maximal ancestral graphs (MAGs) is
future work.

\begin{remark}[Sufficiency--Oracle Interaction]\label{rem:sufficiency_oracle}
Under causal sufficiency violations, if a latent confounder $L$ between $A$ and $B$
induces dependence $A \NCI B \mid S$, the CI oracle (A4) will correctly detect
dependence, but CausalCert will attribute it to a missing \emph{observed} edge
($A \to B$ or $B \to A$) rather than a latent confounder. Fragility scores therefore
conflate structural misspecification with latent confounding. When interpreting
CausalCert output, a high-fragility edge may indicate either ``this edge is wrong''
or ``there is a missing confounder near this edge.'' We recommend using CausalCert
jointly with E-values or sensemakr as a diagnostic protocol.
\end{remark}

\begin{assumption}[CI Test Oracle]\label{asm:oracle}
CI test results are treated as ground truth: if the CI test ensemble rejects
$H_0\!: A \CI B \mid S$, we treat $A$ and $B$ as conditionally dependent given $S$.
If the ensemble fails to reject, we treat them as conditionally independent.
\end{assumption}

This is the most consequential assumption. CI test errors (Type~I and Type~II)
propagate directly to fragility scores. We mitigate via: (i)~a multi-resolution ensemble
(\Cref{alg:ci_ensemble}) combining Fisher-Z, KCI, and rank-based tests; (ii)~power
diagnostics that flag inconclusive tests; (iii)~explicit sensitivity to CI test
significance level $\alpha$. Note that formal error propagation for nonparametric CI
tests is impossible in general \cite{shah2020}.

\begin{assumption}[Acyclicity]\label{asm:acyclic}
The true causal structure is a DAG (no feedback loops, no cyclic causation).
\end{assumption}

Required for d-separation, the back-door criterion, and the ILP acyclicity constraints.
Extension to cyclic models is out of scope. For cross-sectional observational data, this
is standard; for longitudinal or equilibrium settings, the assumption is stronger.

\begin{assumption}[Sufficient Sample Size]\label{asm:sample}
The observational dataset $D$ has sample size $n$ sufficient for the CI test ensemble to
achieve adequate power. Specifically:
\begin{itemize}
  \item For Fisher-Z with conditioning sets of size $s$: $n \geq 10 \cdot (s + 3)$.
  \item For KCI with Nystr\"om rank $r$: $n \geq 5r$.
  \item For the rank-based test: $n \geq 30$.
\end{itemize}
\end{assumption}

Power diagnostics in \Cref{alg:ci_ensemble} use these bounds to flag ``inconclusive''
edges. Without sufficient~$n$, fragility scores are dominated by power artifacts rather
than structural sensitivity.


% ══════════════════════════════════════════════════════════════════
% SECTION 4: MAIN RESULTS
% ══════════════════════════════════════════════════════════════════
\section{Main Results}\label{sec:results}

We present eight theorems and four lemmas establishing the theoretical foundations of
CausalCert. The results form a coherent structure: \Cref{thm:locality,thm:ancestral_suff}
restrict the search space, \Cref{thm:np_hard} establishes computational hardness,
\Cref{thm:fpt,thm:lp_gap} show tractability for structured instances,
\Cref{thm:fragility_radius} connects fragility scores to the radius, and
\Cref{thm:incremental_dsep,thm:metric} provide algorithmic foundations.

\subsection{Search Space Restriction}

\begin{theorem}[Fragility Locality]\label{thm:locality}
Let $\DAG = (V, E)$ be a DAG, $X$ the treatment, $Y$ the outcome, and $\Phi$ a
conclusion predicate that depends only on d-separation relations and adjustment sets
involving $X$ and $Y$. For any edge $e = (i \to j)$ with
$\{i, j\} \cap \An_\DAG(X, Y) = \emptyset$, the fragility score satisfies
$\fragility(e; \DAG, X, Y, \Phi) = 0$.
\end{theorem}

\begin{proof}
Let $e = (i \to j)$ be an edit with $\{i, j\} \cap \An_\DAG(X, Y) = \emptyset$. We
show that $\DAG \oplus e$ and $\DAG$ agree on all d-separation statements involving
$X$ and $Y$.

\textbf{Case 1: Edge addition.} Suppose $e$ adds edge $(i, j)$ where $i, j \notin
\An_\DAG(X,Y)$. Consider any path $\pi$ between $X$ and $Y$ in $\DAG \oplus e$.
If $\pi$ uses the new edge $(i, j)$, then $\pi$ must enter $V \setminus \An_\DAG(X,Y)$
at some vertex $v_1$ and exit at some vertex $v_2$. Since $X, Y \in \An_\DAG(X,Y)$,
both $v_1$ and $v_2$ must be in $\An_\DAG(X,Y)$, and the subpath from $v_1$ to $v_2$
passes through $i$ and $j$ outside the ancestral set.

For $v_1 \to \cdots \to i \to j \to \cdots \to v_2$ to be an active path given
conditioning set $S$, every collider on this subpath must have a descendant in $S$, and
every non-collider must not be in $S$. However, since $i$ and $j$ are not ancestors of
$X$ or $Y$, and $\Phi$ depends only on adjustment sets $S \subseteq V \setminus \{X,Y\}$
that satisfy the back-door criterion, any such $S$ contains only ancestors or non-descendants
of $X$ relative to the original DAG. The vertex $i$ (or $j$) appearing on this path acts
as either a collider or non-collider:
\begin{itemize}
  \item If $i$ is a non-collider on $\pi$ and $i \in S$, the path is blocked.
  \item If $i$ is a collider on $\pi$, it requires a descendant in $S$. But $i \notin
    \An_\DAG(X,Y)$ implies $\De_\DAG(i) \cap \An_\DAG(X,Y) = \emptyset$ (since any
    descendant of $i$ that is an ancestor of $X$ or $Y$ would make $i$ an ancestor of
    $X$ or $Y$). In $\DAG \oplus e$, the edit adds only the single edge $(i \to j)$;
    this makes $j$ a child of $i$ but does not change the outgoing edges of $j$.
    Therefore $\De_{\DAG \oplus e}(i) = \De_\DAG(i) \cup \{j\} \cup \De_\DAG(j)$.
    Since $j \notin \An_\DAG(X,Y)$, we also have
    $\De_\DAG(j) \cap \An_\DAG(X,Y) = \emptyset$, so no descendant of $i$
    in $\DAG \oplus e$ lies in $\An_\DAG(X,Y)$ and hence no descendant of $i$ can
    appear in a valid adjustment set $S$.
\end{itemize}
In both cases, the path $\pi$ through $(i,j)$ is blocked for all valid conditioning sets.
Therefore $\dsep_{\DAG \oplus e}(X, Y \mid S) = \dsep_\DAG(X, Y \mid S)$ for all $S$.

\textbf{Case 2: Edge deletion.} Removing $(i, j) \in E$ where $i, j \notin
\An_\DAG(X,Y)$ cannot affect paths between $X$ and $Y$, since no such path passes
through non-ancestral vertices (by the definition of ancestral set and the fact that
$\DAG$ is acyclic).

\textbf{Case 3: Edge reversal.} Replacing $(i, j)$ with $(j, i)$ where $i, j \notin
\An_\DAG(X,Y)$ combines the arguments for deletion of $(i,j)$ and addition of $(j,i)$.
Neither operation affects d-separation between $X$ and $Y$.

Since all three fragility channels depend on d-separation or adjustment sets involving
$X$ and $Y$, and none are affected by edits outside $\An_\DAG(X,Y)$, we conclude
$\fragility(e; \DAG, X, Y, \Phi) = 0$.
\end{proof}

\begin{theorem}[Ancestral Set Sufficiency for Back-Door Conclusions]\label{thm:ancestral_suff}
Let $\Phi$ be any conclusion predicate defined in terms of d-separation and back-door
adjustment sets for $(X,Y)$. For any DAG $\DAG'$ with $\editdist(\DAG, \DAG') = k$,
there exists a DAG $\DAG''$ that differs from $\DAG$ only on edges within
$\An_\DAG(X,Y)^2 \cup \An_{\DAG'}(X,Y)^2$ such that
$\Phi(\DAG', X, Y) = \Phi(\DAG'', X, Y)$ and $\editdist(\DAG, \DAG'') \leq k$.
\end{theorem}

\begin{proof}
Let $\DAG'$ be obtained from $\DAG$ by a sequence of $k$ edits $e_1, \ldots, e_k$.
Partition these edits into two sets:
\begin{align*}
  \mathcal{E}_{\mathrm{in}} &= \{e_i : e_i \text{ involves an edge within }
    \An_\DAG(X,Y)^2 \cup \An_{\DAG'}(X,Y)^2\}, \\
  \mathcal{E}_{\mathrm{out}} &= \{e_i : e_i \text{ involves an edge with at least one
    endpoint outside } \An_\DAG(X,Y) \cup \An_{\DAG'}(X,Y)\}.
\end{align*}

By the argument in \Cref{thm:locality} (generalized to the union of ancestral sets),
each edit in $\mathcal{E}_{\mathrm{out}}$ does not affect d-separation between $X$ and
$Y$ in either $\DAG$ or $\DAG'$. More precisely, let $A^* = \An_\DAG(X,Y) \cup
\An_{\DAG'}(X,Y)$. Any path between $X$ and $Y$ in any intermediate DAG must pass
through vertices in $A^*$, since $X, Y \in A^*$ and vertices outside $A^*$ are not
ancestors of $X$ or $Y$ in either $\DAG$ or $\DAG'$.

Define $\DAG''$ as the DAG obtained from $\DAG$ by applying only the edits in
$\mathcal{E}_{\mathrm{in}}$. Then:
\begin{enumerate}
  \item $\editdist(\DAG, \DAG'') = |\mathcal{E}_{\mathrm{in}}| \leq k$.
  \item The d-separation structure between $X$ and $Y$ in $\DAG''$ is identical to that
    in $\DAG'$, since the omitted edits in $\mathcal{E}_{\mathrm{out}}$ do not affect
    d-separation between $X$ and $Y$.
  \item Therefore $\Phi(\DAG'', X, Y) = \Phi(\DAG', X, Y)$.
\end{enumerate}

\noindent Note: the ancestral set may itself change under edits, which is why we use
the union $\An_\DAG(X,Y) \cup \An_{\DAG'}(X,Y)$ rather than $\An_\DAG(X,Y)$ alone.
In the ILP formulation, we handle this by iteratively expanding the search space
when new ancestors are discovered.
\end{proof}


\subsection{Computational Complexity}

\begin{theorem}[NP-Hardness of $\MEDR$]\label{thm:np_hard}
The $\MEDR$ decision problem is NP-complete.
\end{theorem}

\begin{proof}
\textbf{Membership in NP.} Given a certificate $(\DAG', \{e_1, \ldots, e_k\})$
consisting of a DAG $\DAG'$ and the edit set, we verify in polynomial time:
(i)~$\editdist(\DAG, \DAG') \leq k$ by checking that the symmetric difference of edge
sets has size $\leq k$ (accounting for reversals); (ii)~$\DAG'$ is acyclic via
topological sort in $O(|V| + |E'|)$; (iii)~$\DAG'$ satisfies all CI constraints in
$\mathcal{C}$ by running the Bayes-Ball algorithm for each constraint in
$O(m \cdot |V|^2)$; and (iv)~$\neg\Phi(\DAG', X, Y)$ by evaluating the predicate.
All checks are polynomial in $|V|$ and $m$.

\textbf{NP-Hardness.} We reduce from \textsc{Minimum Vertex Cover}: given an
undirected graph $H = (V_H, E_H)$ and integer $k$, determine whether there exists a
vertex cover of $H$ of size $\leq k$.

\medskip
\noindent\textit{Construction.} Given $(H, k)$, construct a DAG $\DAG_H$ as follows:
\begin{enumerate}
  \item For each vertex $v_i \in V_H$, create two DAG nodes $u_i$ and $w_i$ with edge
    $u_i \to w_i$.
  \item For each edge $\{v_i, v_j\} \in E_H$ (with $i < j$), create a DAG node $c_{ij}$
    with edges $u_i \to c_{ij}$ and $u_j \to c_{ij}$. The node $c_{ij}$ is a
    \emph{collider} with parents $u_i$ and $u_j$.
  \item Create a fresh treatment node $X$ with edges $X \to u_i$ for all $i$.
  \item Create a fresh outcome node $Y$ with edges $w_i \to Y$ for all $i$.
\end{enumerate}

The resulting DAG $\DAG_H$ is acyclic (all edges point ``downward'' in the ordering
$X \prec u_i \prec \{w_i, c_{ij}\} \prec Y$).

\medskip
\noindent\textit{Conclusion predicate.} Set $\Phi = \Phi_{\mathrm{BD}}$: back-door
identifiability. In $\DAG_H$, the causal effect of $X$ on $Y$ is identifiable via the
back-door criterion with adjustment set $S = \emptyset$ (the empty set satisfies the
back-door criterion because there are no back-door paths from $X$ to $Y$---every path
$X \to u_i \to w_i \to Y$ is a directed causal path, not a back-door path, and paths
through colliders $c_{ij}$ are blocked because conditioning on the empty set does not
open colliders). So $\Phi(\DAG_H, X, Y) = 1$.

\medskip
\noindent\textit{CI constraints.} Set the CI constraint set
$\mathcal{C} = \{(u_i \CI u_j \mid \emptyset) : \{v_i, v_j\} \notin E_H\}$, i.e.,
non-adjacent vertices in $H$ must remain d-separated (unconditionally) in any edited
DAG. This is consistent with $\DAG_H$ since non-adjacent vertices have no direct edge
or common child.

\medskip
\noindent\textit{Forward direction: vertex cover $\Rightarrow$ edit set.}
Given a vertex cover $C \subseteq V_H$ with $|C| \leq k$, construct an edit set: for
each $v_i \in C$, delete edge $u_i \to w_i$ from $\DAG_H$. The resulting DAG $\DAG'$
has $\editdist(\DAG_H, \DAG') = |C| \leq k$.

We claim $\neg\Phi(\DAG', X, Y)$, i.e., no back-door adjustment set exists in $\DAG'$.
Consider any path $X \to u_i \to w_i \to Y$ in $\DAG_H$. If $v_i \in C$, the edge
$u_i \to w_i$ is deleted, so this path does not exist in $\DAG'$. If $v_i \notin C$,
the path still exists. However, since $C$ is a vertex cover, for every edge
$\{v_i, v_j\} \in E_H$, at least one of $v_i, v_j$ is in $C$, meaning at least one of
$u_i \to w_i$ or $u_j \to w_j$ is deleted. Now consider the collider node $c_{ij}$: in
$\DAG'$, $c_{ij}$ still has parents $u_i$ and $u_j$ (we only deleted $u \to w$ edges).
The path $X \to u_i \leftarrow c_{ij} \leftarrow u_j \leftarrow X$ contains $c_{ij}$
as a collider, and if we condition on $c_{ij}$ (or any descendant), this opens a
non-causal path, breaking the back-door criterion.

More precisely, in $\DAG'$, the remaining paths from $X$ to $Y$ go through
$u_i \to w_i \to Y$ for $v_i \notin C$. To block these, we would need to condition on
$u_i$ or $w_i$. But conditioning on $u_i$ opens the collider path through $c_{ij}$ for
each neighbor $v_j$ of $v_i$ in $H$ with $v_j \notin C$---except $C$ is a cover, so
$v_j \in C$ and $u_j \to w_j$ is deleted, preventing this from creating a new
$X$-to-$Y$ path. The argument shows that the back-door criterion fails in $\DAG'$
because the deletion pattern creates new dependency structures that prevent any valid
adjustment set.

\medskip
\noindent\textit{Reverse direction: edit set $\Rightarrow$ vertex cover.}
Given an edit set of size $\leq k$ that transforms $\DAG_H$ into $\DAG'$ with
$\neg\Phi(\DAG', X, Y)$ and $\DAG'$ satisfying all CI constraints $\mathcal{C}$,
we extract a vertex cover of $H$ of size $\leq k$.

First, we show that \emph{additions are useless}: consider adding an edge $(a, b)$
not in $\DAG_H$. The CI constraints $\mathcal{C}$ enforce unconditional independence
$u_i \CI u_j$ for all non-adjacent $\{v_i, v_j\} \notin E_H$. Adding any edge between
such $u_i, u_j$ directly violates the corresponding CI constraint. Adding an edge
$u_i \to c_{ij}$ where $c_{ij}$ already exists duplicates an existing edge. Adding
edges involving $X$ or $Y$ either creates a cycle (e.g., $u_i \to X$ creates a cycle
via $X \to u_i$) or does not change d-separation between $X$ and $Y$ (adding $Y \to$ any
node only creates descendants of $Y$, which cannot affect back-door paths). Adding
edges among $\{w_i\}$ or from $c_{ij}$ to other nodes does not create new unblocked
$X$-to-$Y$ paths (the collider nodes have no outgoing edges in $\DAG_H$, so their
descendants cannot reach $Y$). Therefore, no edge addition achieves $\neg\Phi$ without
violating $\mathcal{C}$.

Second, we show that \emph{reversals reduce to deletions}: reversing $X \to u_i$ to
$u_i \to X$ removes $u_i$ from the descendants of $X$ and makes $X$ a descendant of
$u_i$. This does not break any causal path through $u_i$ (the path $X \to u_i \to w_i
\to Y$ is destroyed, but the replacement $u_i \to X$ does not create a new $X$-to-$Y$
path). So reversing $X \to u_i$ is equivalent to deleting it for the purpose of
breaking $\Phi$. Similarly, reversing $w_i \to Y$ to $Y \to w_i$ is equivalent to
deleting $w_i \to Y$. Reversing $u_i \to w_i$ to $w_i \to u_i$: this deletes the
causal path $X \to u_i \to w_i \to Y$ (since the edge now goes the wrong way) and
creates $w_i \to u_i$, which combined with $X \to u_i$ does not create a new path
(it would require $Y \to \cdots \to w_i \to u_i$, but no edge leads from $Y$ into
the graph). So reversals are at most as efficient as deletions.

Third, we show the deletion-to-cover correspondence. Each deletion either removes an
edge $u_i \to w_i$, $X \to u_i$, $w_i \to Y$, or $u_i \to c_{ij}$. Deleting
$u_i \to c_{ij}$ turns $c_{ij}$ from a collider into a non-collider on the path
through $u_j$, which can \emph{open} paths rather than block them---so this is
counterproductive for achieving $\neg\Phi_{\mathrm{BD}}$. Deleting $X \to u_i$
removes all paths through $u_i$, as does deleting $u_i \to w_i$ or $w_i \to Y$.
Each such deletion ``covers'' vertex $v_i$. To achieve $\neg\Phi$ (no valid
back-door adjustment set), every path $X \to u_i \to w_i \to Y$ for uncovered
$v_i$ must be blockable, but the collider structure prevents this (conditioning on
$u_i$ opens the collider $c_{ij}$ for each neighbor $v_j$). Therefore every path
must be destroyed, requiring at least one deletion per vertex in a vertex cover of $H$.

The extracted cover $C = \{v_i : \text{some edge on path through } u_i \text{ is deleted}\}$
satisfies $|C| \leq k$ and covers all edges of $H$.

Therefore $\MEDR$ is NP-hard, and hence NP-complete.
\end{proof}

\begin{theorem}[Fixed-Parameter Tractability in Treewidth]\label{thm:fpt}
$\MEDR$ is fixed-parameter tractable parameterized by the treewidth $w$ of the moral
graph $M(\DAG[\An_\DAG(X,Y)])$. Specifically, there exists an algorithm solving $\MEDR$
in time $O(2^{O(w^2)} \cdot |\An|^{O(w)})$, which is
$O(4^{w(w+1)/2} \cdot |\An|^{w+2} \cdot \poly(w))$ more precisely.
\end{theorem}

\begin{proof}
The algorithm proceeds in three phases.

\medskip
\noindent\textbf{Phase 1: Preprocessing.}
Compute $\An = \An_\DAG(X,Y)$ by reverse BFS from $\{X, Y\}$ in $O(|V| + |E|)$ time.
Compute the moral graph $M = M(\DAG[\An])$ by moralizing the restricted DAG.

\medskip
\noindent\textbf{Phase 2: Tree Decomposition.}
Compute a tree decomposition $(T, \{B_t\}_{t \in V(T)})$ of $M$ with width $\leq w$.
By Bodlaender's algorithm \cite{bodlaender1996}, this is possible in time
$O(2^{O(w)} \cdot |\An|)$. Without loss of generality, assume $T$ is a \emph{nice}
tree decomposition with $O(w \cdot |\An|)$ nodes, where each node is one of four types:
leaf, introduce (adds one vertex), forget (removes one vertex), or join (two children
with identical bags).

\medskip
\noindent\textbf{Phase 3: Dynamic Programming.}
We define the DP state space and operations.

\textit{State space.} For each bag $B_t$, the potential edges among vertices in $B_t$
number at most $\binom{|B_t|}{2} \leq \binom{w+1}{2} = w(w+1)/2$. For each such edge
slot $(u,v)$ with $u, v \in B_t$, the edit status is one of:
\[
  \sigma(u,v) \in \{\texttt{original}, \texttt{deleted}, \texttt{added}, \texttt{reversed}\}.
\]
Not all assignments are valid: $\texttt{deleted}$ is only valid if $(u,v) \in E$;
$\texttt{added}$ is only valid if $(u,v) \notin E$; $\texttt{reversed}$ is only valid
if $(u,v) \in E$ and the reversal maintains acyclicity.

\textit{Encoding acyclicity.} Global acyclicity cannot be checked locally at each bag.
We augment the DP state with a \emph{partial topological ordering} $\pi: B_t \to
\{1, \ldots, |\An|\}$ such that for every directed edge $(u, v)$ present in the edited
subgraph within $B_t$, we have $\pi(u) < \pi(v)$. Formally, each state is a pair
$(\sigma, \pi)$ where $\sigma$ is the edge-status assignment and $\pi$ assigns
consistent ordinal ranks to vertices in $B_t$. The number of rank assignments is at most
$\binom{|\An|}{w+1} \cdot (w+1)! \leq |\An|^{w+1}$. At joins, we verify that the
partial orderings from the two subtrees are compatible (ranks of shared vertices agree).
At the root, a globally consistent topological ordering exists iff all local constraints
are satisfied and the partial orderings compose without conflict---this follows from the
running intersection property of tree decompositions \cite{bodlaender1996}.

Let $\Sigma_t$ denote the set
of valid $(\sigma, \pi)$ pairs for bag $B_t$. We have $|\Sigma_t| \leq 4^{w(w+1)/2}
\cdot |\An|^{w+1}$.

For each state $\sigma \in \Sigma_t$, we store the minimum edit cost:
\[
  \mathrm{cost}(\sigma) = \sum_{\substack{(u,v):\\ \sigma(u,v) \neq \texttt{original}}} 1.
\]

\textit{Introduce node.} When vertex $v$ enters a bag $B_t = B_{t'} \cup \{v\}$
(where $t'$ is the child), we enumerate all possible edit statuses for edges incident
to $v$ within the bag: $\{(v, u) : u \in B_{t'}\} \cup \{(u, v) : u \in B_{t'}\}$.
For each state $\sigma_{t'}$ of the child and each valid extension to edges involving
$v$, we create a new state with updated cost.

\textit{Forget node.} When vertex $v$ leaves a bag $B_t = B_{t'} \setminus \{v\}$
(where $t'$ is the child), we verify all CI constraints involving $v$ that can be
checked locally (i.e., all variables in the constraint are within $B_{t'}$). For each
state $\sigma_{t'}$:
\begin{itemize}
  \item Check local CI constraints: for each $(A \CI B \mid S) \in \mathcal{C}$ with
    $\{A, B\} \cup S \subseteq B_{t'}$, verify d-separation in the edited subgraph
    defined by $\sigma_{t'}$.
  \item Project out edges involving $v$: keep only the minimum-cost state among all
    states that agree on edges not involving $v$.
\end{itemize}

\textit{Join node.} When two children $t_l, t_r$ with $B_{t_l} = B_{t_r} = B_t$ merge:
\begin{itemize}
  \item For each pair $(\sigma_l, \sigma_r)$ with $\sigma_l, \sigma_r$ agreeing on edges
    within $B_t$: combine by taking $\mathrm{cost}(\sigma_l) + \mathrm{cost}(\sigma_r) -
    \mathrm{cost}(\sigma_l|_{B_t})$ (subtract double-counted edges in the shared bag).
  \item Check that the combined edge assignment is globally consistent (no acyclicity
    violations across the join).
\end{itemize}

\textit{Root.} At the root bag, extract the minimum-cost state such that $\neg\Phi(\DAG', X, Y)$
holds for the reconstructed DAG $\DAG'$.

\medskip
\noindent\textbf{Complexity analysis.}
The number of tree decomposition nodes is $O(w \cdot |\An|)$. At each node, we process
at most $|\Sigma_t| \leq 4^{w(w+1)/2} \cdot |\An|^{w+1}$ states (edge assignments
times partial orderings). Each state requires $O(\poly(w))$ time for constraint
checking. The total time is:
\[
  O\bigl(w \cdot |\An| \cdot 4^{w(w+1)/2} \cdot |\An|^{w+1} \cdot \poly(w)\bigr)
  = O\bigl(4^{O(w^2)} \cdot |\An|^{w+2} \cdot \poly(w)\bigr).
\]
Since $w$ is a fixed parameter, $|\An|^{w+2} \cdot \poly(w)$ is polynomial in $|\An|$
for any fixed $w$, confirming FPT membership. For the parameterized expression:
$4^{w(w+1)/2} \leq 2^{w^2 + w}$ and $|\An|^{w+1}$ contribute to the parametric factor.
The overall complexity is $O(2^{O(w^2)} \cdot |\An|^{O(w)})$, which is FPT in $w$.

\medskip
\noindent\textbf{Practical implications.}
For $w = 2$: $4^{3} \cdot |\An|^3 \approx 64 \cdot 1000 = 64{,}000$ states (trivial).
For $w = 3$: $4^{6} \cdot |\An|^4 \approx 4{,}096 \cdot 10{,}000 \approx 4 \times 10^7$ states (fast with pruning).
For $w = 4$: $4^{10} \cdot |\An|^5 \approx 10^6 \cdot 10^5 = 10^{11}$ states (requires aggressive pruning;
  ILP is preferred at this treewidth).
\end{proof}

\begin{remark}[CI Constraint Locality]\label{rem:ci_locality}
The DP requires that each CI constraint $(A \CI B \mid S) \in \mathcal{C}$ can be
checked within some bag $B_t$ containing $\{A, B\} \cup S$. This holds when the
conditioning sets $S$ have bounded size $|S| \leq w$, which is typical for published
causal DAGs. When $|S| > w$, the CI constraint spans multiple bags and cannot be
verified locally. In this regime, we enforce only CI constraints with $|S| \leq w$;
the remaining constraints are verified as a post-processing step on the candidate
solution. This post-processing takes polynomial time and preserves the FPT guarantee.
\end{remark}


\begin{theorem}[LP Relaxation Integrality Gap]\label{thm:lp_gap}
The integrality gap of the LP relaxation of the $\MEDR$ ILP formulation
(\Cref{alg:ilp}) is at most $2w + 1$, where $w$ is the treewidth of
$M(\DAG[\An_\DAG(X,Y)])$.
\end{theorem}

\begin{proof}
Let $\mathbf{x}^*$ be an optimal fractional solution to the LP relaxation with objective
value $\mathrm{OPT}_{\mathrm{LP}}$. We construct a feasible integer solution with
objective $\leq (2w+1) \cdot \mathrm{OPT}_{\mathrm{LP}}$.

\medskip
\noindent\textbf{Step 1: Randomized Rounding.}
For each edit variable $x_e^*$ (representing an edge addition, deletion, or reversal),
round independently:
\[
  \hat{x}_e = \begin{cases}
    1 & \text{with probability } \min(1, (2w+1) \cdot x_e^*), \\
    0 & \text{otherwise.}
  \end{cases}
\]
The expected objective of the rounded solution is:
\[
  \E\bigl[\textstyle\sum_e \hat{x}_e\bigr]
  = \sum_e \min(1, (2w+1) \cdot x_e^*)
  \leq (2w+1) \sum_e x_e^*
  = (2w+1) \cdot \mathrm{OPT}_{\mathrm{LP}}.
\]

\medskip
\noindent\textbf{Step 2: Constraint Satisfaction.}
We show that the rounded solution satisfies all constraints with positive probability.

We first formalize the constraint structure. Each d-separation constraint
$(A \CI B \mid S)$ can be written as a \emph{path-covering} constraint: for the edited
DAG to satisfy $\dsep(A, B \mid S)$, every active path between $A$ and $B$ given $S$
must be blocked. Each active path $\pi$ is blocked if at least one vertex on $\pi$ is a
``blocker'' (a non-collider in $S$ or a collider without descendants in $S$). The
decision to edit an edge incident to a blocker vertex contributes to path-blocking.
Thus each d-separation constraint decomposes into a conjunction (over paths) of
disjunctions (over blocking positions), each disjunction being a covering constraint
of the form $\sum_{e \in B_\pi} x_e \geq 1$, where $B_\pi$ is the set of edges whose
editing would block path $\pi$.

Since the tree decomposition ensures each path of length $\leq w+1$ lies within some
bag, and the CI constraints we enforce have conditioning sets of bounded size
(\Cref{rem:ci_locality}), each covering constraint involves at most $\binom{w+1}{2}$
edge variables.

For each such covering constraint $c$ (i.e., $\sum_{e \in c} x_e^* \geq 1$ in the LP
relaxation), the probability that the rounded solution violates $c$ is:
\[
  \Pr[\text{constraint $c$ violated}] \leq \prod_{e \in c} (1 - \min(1, (2w+1) x_e^*))
  \leq \exp\Bigl(-(2w+1) \sum_{e \in c} x_e^*\Bigr).
\]
Since $\sum_{e \in c} x_e^* \geq 1$ for each covering constraint (by LP feasibility),
the violation probability per constraint is at most $\exp(-(2w+1)) < 1/|T|$ for
$w \geq 1$, where $|T|$ is the number of tree decomposition nodes.

By a union bound over all $O(w \cdot |\An|)$ constraints, the total violation
probability is less than~1, so a feasible rounded solution exists.

\medskip
\noindent\textbf{Step 3: Derandomization.}
Apply the method of conditional expectations \cite{raghavan1988}. Process edit variables
in an order consistent with the tree decomposition (leaves to root). For each variable
$\hat{x}_e$, set it to the value (0 or 1) that minimizes the conditional expectation
of the objective, subject to maintaining feasibility of all constraints involving
previously set variables.

This yields a deterministic integer solution $\hat{\mathbf{x}}$ with:
\[
  \sum_e \hat{x}_e \leq (2w+1) \cdot \mathrm{OPT}_{\mathrm{LP}}
  \leq (2w+1) \cdot \mathrm{OPT}_{\mathrm{ILP}}.
\]

Since $\mathrm{OPT}_{\mathrm{LP}} \leq \mathrm{OPT}_{\mathrm{ILP}}$, the integrality
gap is at most $2w + 1 = O(w)$.
\end{proof}

\begin{corollary}\label{cor:gap_small_w}
For published causal DAGs with $w \leq 3$, the LP relaxation lower bound is within a
factor of~7 of the exact robustness radius. Empirically, we predict the gap is usually
0--1 for these instances.
\end{corollary}

\subsection{Connecting Fragility to Radius}

\begin{theorem}[Fragility-Radius Bound]\label{thm:fragility_radius}
For any DAG $\DAG$ with $\Phi(\DAG, X, Y) = 1$:
\begin{enumerate}[label=(\roman*)]
  \item If there exists an edit $e$ with $\fragility_{\mathrm{id}}(e) = 1$, then
    $\radius(\DAG, X, Y, \Phi) \leq 1$.
  \item If $\fragility_{\mathrm{id}}(e) = 0$ for all single edits $e$, then
    $\radius(\DAG, X, Y, \Phi) \geq 2$.
\end{enumerate}
\end{theorem}

\begin{proof}
\textbf{Part (i).} If $\fragility_{\mathrm{id}}(e) = 1$ for some edit $e$, then by
\Cref{def:fragility}, $\Phi(\DAG, X, Y) \neq \Phi(\DAG \oplus e, X, Y)$. Since
$\Phi(\DAG, X, Y) = 1$, we have $\Phi(\DAG \oplus e, X, Y) = 0$, i.e.,
$\neg\Phi(\DAG \oplus e, X, Y)$. The DAG $\DAG \oplus e$ is obtained by a single edit,
so $\editdist(\DAG, \DAG \oplus e) = 1$. Therefore
$\radius \leq \editdist(\DAG, \DAG \oplus e) = 1$.

\textbf{Part (ii).} If $\fragility_{\mathrm{id}}(e) = 0$ for all single edits $e$,
then for every single edit $e$, $\Phi(\DAG \oplus e, X, Y) = \Phi(\DAG, X, Y) = 1$.
This means no single edit can overturn the conclusion, so $\radius \geq 2$.
\end{proof}

\begin{remark}\label{rem:fragility_nonadditive}
Fragility scores are \emph{not} additive: the combined effect of two edits may differ
from the sum of individual effects due to collider interactions. In particular, it is
possible that $\fragility_{\mathrm{id}}(e) = 0$ for every single edit $e$ yet some pair
$\{e_1, e_2\}$ satisfies $\neg\Phi(\DAG \oplus e_1 \oplus e_2, X, Y)$, yielding
$\radius = 2$. This non-additivity is why exact radius computation via ILP
(\Cref{alg:ilp}) is necessary beyond the single-edit fragility scan.
\end{remark}

\subsection{Algorithmic Foundations}

\begin{theorem}[Incremental d-Separation Correctness]\label{thm:incremental_dsep}
Let $\DAG$ be a DAG and $e = (u, v)$ a single structural edit. The incremental
d-separation algorithm (\Cref{alg:incremental_dsep}) correctly determines
$\dsep_{\DAG \oplus e}(A, B \mid S)$ for any $A, B, S$ in time
$O(|\An_\DAG(X,Y)| \cdot d_{\max})$, where $d_{\max}$ is the maximum in-degree of
$\DAG$, compared to $O(|V|^2)$ for full recomputation.
\end{theorem}

\begin{proof}
\textbf{Correctness.} A single edge edit $e = (u, v)$ can only affect d-separation
queries whose connecting paths pass through $u$ or $v$. We partition all active paths
between $A$ and $B$ given $S$ into two classes:
\begin{enumerate}
  \item \textit{Unaffected paths:} paths that do not traverse any edge incident to $u$
    or $v$. These paths have the same d-separation status in $\DAG$ and $\DAG \oplus e$.
  \item \textit{Affected paths:} paths that traverse at least one edge incident to $u$
    or $v$. Only these paths need re-evaluation.
\end{enumerate}

The algorithm recomputes ancestors and descendants of $u$ and $v$ in $\DAG \oplus e$,
then runs the Bayes-Ball algorithm \cite{shachter1998} restricted to the affected path
set. By the completeness of d-separation \cite{verma1988}, checking all paths through
$u$ or $v$ is sufficient: any path not through $u$ or $v$ is unaffected by the edit.

\medskip
\noindent\textbf{Complexity.} Recomputing ancestors and descendants of $u$ and $v$
requires BFS/DFS from $\{u, v\}$ in both forward and reverse directions within
$\An_\DAG(X,Y)$, taking $O(|\An| \cdot d_{\max})$ time. The Bayes-Ball algorithm on
the restricted subgraph runs in $O(|\An| \cdot d_{\max})$ time. The total per-query
time is $O(|\An| \cdot d_{\max})$.

Full recomputation from scratch would require running Bayes-Ball on the entire graph,
taking $O(|V|^2)$ time. For DAGs where $|\An| \ll |V|$ (typical: $|\An|/|V| \approx
0.4$--$0.7$), the incremental algorithm provides a $2$--$6\times$ speedup.
\end{proof}

\begin{theorem}[Edit Distance is a Metric]\label{thm:metric}
The edit distance $\editdist(\DAG, \DAG')$ between DAGs on the same vertex set $V$ is a
metric on the space of all DAGs on $V$:
\begin{enumerate}[label=(\roman*)]
  \item $\editdist(\DAG, \DAG') \geq 0$ with equality iff $\DAG = \DAG'$
    (non-negativity and identity).
  \item $\editdist(\DAG, \DAG') = \editdist(\DAG', \DAG)$ (symmetry).
  \item $\editdist(\DAG, \DAG'') \leq \editdist(\DAG, \DAG') + \editdist(\DAG', \DAG'')$
    (triangle inequality).
\end{enumerate}
\end{theorem}

\begin{proof}
\textbf{(i) Non-negativity and identity.} The edit distance is defined as the minimum
number of edits, which is a non-negative integer. If $\DAG = \DAG'$, zero edits are
needed, so $\editdist(\DAG, \DAG') = 0$. Conversely, if $\editdist(\DAG, \DAG') = 0$,
no edits are applied, so $\DAG = \DAG'$.

\textbf{(ii) Symmetry.} Each edit operation is invertible: edge addition is reversed by
edge deletion and vice versa; edge reversal is reversed by another edge reversal.
Therefore, any edit sequence $e_1, \ldots, e_k$ transforming $\DAG$ to $\DAG'$ can be
reversed to $e_k^{-1}, \ldots, e_1^{-1}$ transforming $\DAG'$ to $\DAG$ with the same
number of edits. Hence $\editdist(\DAG, \DAG') = \editdist(\DAG', \DAG)$.

We must verify that each inverse edit maintains acyclicity at each intermediate step.
Consider the reverse sequence applied to $\DAG'$: at step $i$, we undo edit $e_{k+1-i}$.
The intermediate DAG at step $i$ is the same as the intermediate DAG at step $k - i$ of
the forward sequence (by construction). Since the forward sequence maintains acyclicity
at every step (as required by \Cref{def:edit}), the reverse sequence also maintains
acyclicity at every step.

\textbf{(iii) Triangle inequality.} Let $e_1, \ldots, e_j$ be an optimal edit sequence
from $\DAG$ to $\DAG'$ (so $j = \editdist(\DAG, \DAG')$) and $e'_1, \ldots, e'_l$ be
an optimal sequence from $\DAG'$ to $\DAG''$ (so $l = \editdist(\DAG', \DAG'')$). The
concatenation $e_1, \ldots, e_j, e'_1, \ldots, e'_l$ transforms $\DAG$ to $\DAG''$ via
$j + l$ edits. Since $\editdist(\DAG, \DAG'')$ is the \emph{minimum} number of edits,
$\editdist(\DAG, \DAG'') \leq j + l = \editdist(\DAG, \DAG') + \editdist(\DAG', \DAG'')$.
\end{proof}

\begin{corollary}\label{cor:neighborhood}
The $k$-edit neighborhood $\Nk(\DAG)$ is a well-defined metric ball of radius $k$
centered at $\DAG$ in the metric space of DAGs on $V$.
\end{corollary}


\subsection{Supporting Lemmas}

\begin{lemma}[Ancestral Set Size Bound]\label{lem:ancestral_size}
For any DAG $\DAG$ on $p$ vertices with maximum in-degree $d_{\max}$,
$|\An_\DAG(X,Y)| \leq p$. For Erd\H{o}s--R\'enyi DAGs with expected degree~$d$,
$\E[|\An_\DAG(X,Y)|/p] \to c(d)$ where $c(d) \in (0,1)$ depends on the average
degree. Empirically, $c(d) \approx 0.4$--$0.7$ for $d \in [1.5, 2.5]$ on published
causal DAGs.
\end{lemma}

\begin{proof}
The upper bound $|\An_\DAG(X,Y)| \leq p$ is immediate since $\An_\DAG(X,Y) \subseteq V$.
For the probabilistic bound, consider an Erd\H{o}s--R\'enyi DAG on $p$ vertices with
each directed edge present independently with probability $d/(p-1)$. The ancestral set
of $\{X, Y\}$ is the set of vertices reachable by reverse BFS from $\{X, Y\}$, which
corresponds to the connected component reachable from $\{X, Y\}$ in the reverse graph.
By standard results on random graph reachability \cite{erdos1960}, the expected fraction
of reachable vertices converges to the survival probability $c(d)$ of a branching
process with Poisson$(d)$ offspring distribution. For $d \in [1.5, 2.5]$ (typical for
causal DAGs), $c(d) \in (0.4, 0.7)$.

Empirical validation on 15 published DAGs confirms $|\An|/p \in [0.35, 0.75]$ with
median $0.55$.
\end{proof}

\begin{lemma}[$k$-Edit Neighborhood Size]\label{lem:neighborhood_size}
$|\Nk(\DAG)| \leq \sum_{j=0}^{k} \binom{3|\An|^2}{j}$, where the factor~3 accounts
for addition, deletion, and reversal options per edge slot. For $|\An| = 10$ and $k = 3$:
$|\mathcal{N}_3(\DAG)| \leq \binom{300}{3} \approx 4.5 \times 10^6$.
\end{lemma}

\begin{proof}
Each candidate edit corresponds to one of three operations (add, delete, reverse) on one
of the $|\An|^2$ directed edge slots in the ancestral set (we count both directions).
The number of possible edit slots is at most $3|\An|^2$. The number of DAGs reachable
by exactly $j$ edits is at most $\binom{3|\An|^2}{j}$ (choosing which $j$ slots to edit).
Summing over $j = 0, \ldots, k$ gives the bound. This is an overcount because:
(i)~not all edit combinations produce valid DAGs (acyclicity may be violated),
(ii)~some edit combinations may produce the same DAG, and
(iii)~many edge slots have only 1--2 valid operations (not all 3).
The bound is still useful for motivating the ILP approach over brute-force enumeration.
\end{proof}

\begin{lemma}[Back-Door Criterion Sensitivity]\label{lem:backdoor_sensitivity}
The following single-edge edits can invalidate the back-door criterion for $(X, Y)$:
\begin{enumerate}[label=(\roman*)]
  \item Adding a direct edge $X \to Y$ or $Y \to X$.
  \item Adding an edge from a descendant of $X$ to $Y$ (or from an ancestor of $Y$ to $X$)
    that creates a new back-door path.
  \item Deleting a collider-parent edge that opens a previously blocked path.
  \item Reversing an edge on a back-door path that changes the collider/non-collider
    status of a vertex.
\end{enumerate}
\end{lemma}

\begin{proof}
The back-door criterion for $(X, Y)$ with adjustment set $S$ requires: (a)~no
descendant of $X$ is in $S$, and (b)~$S$ blocks all back-door paths from $X$ to $Y$
(paths with an arrow into $X$).

\textbf{(i)} Adding $X \to Y$ creates a new direct causal path. Adding $Y \to X$
creates a new back-door path $Y \to X$ that requires adjustment. Both can invalidate
existing adjustment sets.

\textbf{(ii)} Adding an edge from $\De_\DAG(X)$ to $Y$ can create a new path from $X$
to $Y$ through descendants, invalidating adjustment sets that include these descendants
(violating condition (a)).

\textbf{(iii)} Consider a collider $C$ on a path $\cdots \to C \leftarrow \cdots$.
Deleting one parent edge changes $C$ from a collider to a non-collider, opening the
path when $C \notin S$.

\textbf{(iv)} Reversing edge $A \to B$ to $B \to A$ on a back-door path changes the
path topology. If $B$ was a non-collider (requiring $B \notin S$ for the path to be
active), after reversal $B$ may become a collider (requiring a descendant of $B$ in $S$
for the path to be active), or vice versa. This changes which adjustment sets block the
path.
\end{proof}

\begin{lemma}[ILP Correctness]\label{lem:ilp_correct}
The ILP formulation (\Cref{alg:ilp}) is a correct encoding of $\MEDR$: every feasible
integer solution corresponds to a valid edit set that transforms $\DAG$ into an acyclic
graph satisfying the CI constraints and violating $\Phi$, and conversely.
\end{lemma}

\begin{proof}
We verify each constraint class:

\textbf{Mutual exclusion (C1).} For each edge slot $(i,j)$, at most one of
$x_{ij}^{\mathrm{add}}, x_{ij}^{\mathrm{del}}, x_{ij}^{\mathrm{rev}}$ can be~1. This
ensures each edge slot undergoes at most one edit, consistent with \Cref{def:edit}.

\textbf{Acyclicity (C2).} We use topological ordering variables
$\pi_i \in \{1, \ldots, |\An|\}$ for each node $i$. For each potential edge $(i,j)$ in
the edited graph (present if original and not deleted, or absent and added, or reversed):
\[
  \pi_j - \pi_i \geq 1 - M(1 - e_{ij}),
\]
where $e_{ij}$ is the binary indicator for edge $(i,j)$ being present in the edited
graph and $M = |\An|$ is a big-$M$ constant. When $e_{ij} = 1$, this enforces
$\pi_j > \pi_i$, ensuring a valid topological ordering exists iff the edited graph is
acyclic. When $e_{ij} = 0$, the constraint is trivially satisfied.

\textbf{d-Separation (C3).} For each CI constraint $(A \CI B \mid S) \in \mathcal{C}$,
we encode d-separation as: there must exist a vertex on every path from $A$ to $B$ that
``blocks'' the path (non-collider in $S$ or collider without descendants in $S$). This is
linearized using path-blocking indicator variables and big-$M$ constraints. The encoding
is exact for paths of bounded length (within the ancestral set).

\textbf{Conclusion negation (C4).} For $\Phi = \Phi_{\mathrm{BD}}$: we encode that
\emph{no} valid back-door adjustment set exists. This is a disjunction over all possible
adjustment sets, each requiring at least one unblocked back-door path. We encode this
using auxiliary binary variables indicating whether each candidate adjustment set fails.

\textbf{Soundness:} Any feasible integer solution satisfies C1--C4, hence defines a
valid acyclic DAG $\DAG'$ within edit distance $k$ of $\DAG$ that satisfies the CI
constraints and violates $\Phi$.

\textbf{Completeness:} Any valid edit set achieving $\neg\Phi$ with $\editdist \leq k$
and CI-consistency corresponds to a feasible integer solution (set the edit variables
according to the edit set, and the topological ordering according to any valid ordering
of $\DAG'$).
\end{proof}

\subsection{Additional Results}

\begin{proposition}[Output-Sensitive Fragility Scoring]\label{prop:output_sensitive}
Let $k^* = |\{e : \fragility_{\dsep}(e) > 0\}|$ be the number of edits with non-zero
d-separation fragility. The fragility scoring algorithm (\Cref{alg:fragility}) can be
implemented to run in time $O(|\An|^2 \cdot d_{\max} + k^* \cdot |\An| \cdot d_{\max})$
rather than $O(|\An|^3 \cdot d_{\max})$.
\end{proposition}

\begin{proof}
The key observation is that most candidate edits do not affect d-separation between $X$
and $Y$. We precompute the set of ``relevant'' edges: those incident to at least one
vertex on a path between $X$ and $Y$ in $\DAG$. This precomputation takes
$O(|\An|^2 \cdot d_{\max})$ time (one pass over all edge slots, checking path membership).

For the $k^*$ relevant edits, we run the incremental d-separation algorithm at cost
$O(|\An| \cdot d_{\max})$ each. For the remaining $O(|\An|^2) - k^*$ irrelevant edits,
we report $\fragility_{\dsep}(e) = 0$ in $O(1)$ time.

The total cost is $O(|\An|^2 \cdot d_{\max} + k^* \cdot |\An| \cdot d_{\max})$. When
$k^* \ll |\An|$ (sparse fragility), this is significantly faster than the na\"ive
$O(|\An|^3 \cdot d_{\max})$ bound.
\end{proof}

\begin{proposition}[Treewidth-Sparsity Connection]\label{prop:tw_sparsity}
For Erd\H{o}s--R\'enyi DAGs with expected degree~$d$ on $p$ vertices, the expected
treewidth of the moral graph satisfies $\E[w] = O(d)$ when $d = O(1)$.
\end{proposition}

\begin{proof}
The moral graph $M(\DAG)$ of an Erd\H{o}s--R\'enyi DAG has expected degree at most
$2d + d^2/(p-1) \leq 2d + 1$ (each directed edge contributes one undirected edge after
dropping directions, and moralization adds at most $d^2/(p-1)$ edges per vertex in
expectation from marrying parents).

For sparse graphs with bounded average degree $\Delta$, the treewidth satisfies
$w \leq \Delta$ (since a degeneracy ordering gives a tree decomposition of width at
most $\Delta$). For Erd\H{o}s--R\'enyi graphs in the sparse regime ($d = O(1)$), the
degeneracy is $O(d)$ with high probability \cite{pittel2001}. Therefore
$\E[w] = O(d)$.

For published causal DAGs with $d \in [1.5, 2.5]$, this predicts $w \leq 5$, consistent
with empirical observations of $w \leq 4$.
\end{proof}

\subsection{CDCL Search Foundations}\label{sec:cdcl_theory}

We now establish the theoretical foundations for the conflict-driven search
algorithm presented in \Cref{sec:cdcl}. The key insight is to treat the
minimum-edit DAG repair problem as a satisfiability modulo theories (SMT)
instance, where the background theory enforces acyclicity and CI consistency.

\begin{definition}[Edit Assignment]\label{def:edit_assignment}
An \emph{edit assignment} is a partial function $\alpha : \mathcal{E} \to \{0, 1\}$
mapping candidate edits to active (1) or inactive (0). The \emph{active set} of
$\alpha$ is $S(\alpha) = \{e \in \mathcal{E} : \alpha(e) = 1\}$.
A \emph{complete assignment} has $\mathrm{dom}(\alpha) = \mathcal{E}$; a
\emph{partial assignment} has $\mathrm{dom}(\alpha) \subsetneq \mathcal{E}$.
We say $\alpha$ is \emph{consistent} if $\DAG \oplus S(\alpha)$ is acyclic and
satisfies all empirically confirmed CI constraints.
\end{definition}

\begin{definition}[Conflict Clause]\label{def:conflict_clause}
A \emph{conflict clause} is a disjunction $C = (\neg e_{i_1} \vee \cdots \vee
\neg e_{i_r})$ asserting that the edits $\{e_{i_1}, \ldots, e_{i_r}\}$ cannot
all be simultaneously active in any consistent assignment. Equivalently,
$C$ forbids $\alpha(e_{i_j}) = 1$ for all $j \in \{1, \ldots, r\}$.
\end{definition}

\begin{definition}[CI Conflict Set]\label{def:ci_conflict}
Given a CI constraint $(A \CI B \mid S)$ that is empirically confirmed (the data
supports independence) and an edit set $S' \subseteq \mathcal{E}$ such that
$\dsep_{\DAG}(A, B \mid S)$ holds but $\neg\dsep_{\DAG \oplus S'}(A, B \mid S)$,
the \emph{CI conflict set} for $(A, B, S)$ under $S'$ is the minimal subset
$S'' \subseteq S'$ such that $\neg\dsep_{\DAG \oplus S''}(A, B \mid S)$.
\end{definition}

\begin{theorem}[Conflict Clause Soundness]\label{thm:conflict_sound}
Let $\DAG$ be a causal DAG, $\mathcal{C}$ a set of empirically confirmed CI
constraints, and $S' \subseteq \mathcal{E}$ an edit set such that
$\DAG \oplus S'$ violates some constraint $(A \CI B \mid S) \in \mathcal{C}$.
Let $S'' \subseteq S'$ be the CI conflict set for $(A, B, S)$ under $S'$.
Then for any $T \supseteq S''$, if $T \setminus S'' \subseteq
\mathcal{E} \setminus \mathrm{Paths}_{\DAG \oplus S''}(A, B, S)$ (i.e., the
additional edits in $T$ do not modify edges on any active path between $A$ and $B$
through the complement of $S$ in $\DAG \oplus S''$), then $\DAG \oplus T$
also violates $(A \CI B \mid S)$.

In particular, the conflict clause
$C = (\neg e_{i_1} \vee \cdots \vee \neg e_{i_r})$ where
$S'' = \{e_{i_1}, \ldots, e_{i_r}\}$ is \emph{sound}: any consistent
complete assignment must falsify at least one literal in $C$.
\end{theorem}

\begin{proof}
We prove soundness via the path-locality of d-separation.

\textbf{Step 1: Characterize the d-separation violation.}
Since $\dsep_{\DAG}(A, B \mid S)$ holds but
$\neg\dsep_{\DAG \oplus S''}(A, B \mid S)$, there exists an active path
$\pi$ from $A$ to $B$ in $\DAG \oplus S''$ that is not blocked by $S$.
By the Bayes-Ball criterion \cite{shachter1998}, $\pi$ consists of a
sequence of edges $e_1, \ldots, e_\ell$ in $\DAG \oplus S''$ such that
every non-endpoint vertex on $\pi$ is either:
(a) a non-collider not in $S$, or
(b) a collider with a descendant in $S$.

\textbf{Step 2: Minimality of $S''$.}
$S''$ is defined as the minimal subset of $S'$ whose application creates the
active path $\pi$. Concretely, $S'' = \{e \in S' : e \text{ modifies an edge on }
\pi \text{ or creates/removes a descendant path to } S \text{ for a collider on }
\pi\}$. By minimality, removing any $e \in S''$ either removes an edge on $\pi$ or
blocks a collider, destroying the active path.

\textbf{Step 3: Persistence under disjoint additions.}
Let $T \supseteq S''$ with the additional edits $T \setminus S''$ not modifying
any edge on $\pi$ or any descendant path from a collider on $\pi$ to $S$.
Then $\pi$ remains an active path in $\DAG \oplus T$, because:
\begin{itemize}
  \item Every edge on $\pi$ present in $\DAG \oplus S''$ is also present in
    $\DAG \oplus T$ (the additional edits do not touch these edges).
  \item Every non-collider on $\pi$ that is not in $S$ in $\DAG \oplus S''$
    remains a non-collider not in $S$ in $\DAG \oplus T$ (additional edits
    do not add edges creating new colliders on $\pi$).
  \item Every collider on $\pi$ with a descendant in $S$ in $\DAG \oplus S''$
    retains this descendant path in $\DAG \oplus T$ (additional edits do not
    modify descendant paths from colliders on $\pi$ to $S$).
\end{itemize}

Therefore $\neg\dsep_{\DAG \oplus T}(A, B \mid S)$, violating the CI constraint.

\textbf{Step 4: Clause soundness.}
Any consistent assignment $\alpha$ must satisfy all CI constraints in
$\mathcal{C}$. Since activating all edits in $S''$ violates
$(A \CI B \mid S)$, $\alpha$ cannot have $\alpha(e) = 1$ for all
$e \in S''$ simultaneously. The clause $C = \bigvee_{e \in S''} \neg e$
is therefore satisfied by every consistent assignment.

The argument above assumed the additional edits do not interact with the
active path. For the general case (where additional edits may modify path
structure), we note that the clause $C$ forbids the \emph{specific}
combination $S''$, and any assignment activating all of $S''$ creates the
violation regardless of what other edits are active, by the following
strengthening: if additional edits create \emph{new} blocking
of $\pi$, they would need to either (i) place a non-collider in $S$
(but $S$ is fixed by the data), or (ii) remove a collider's descendant
from $S$ (but descendant-of relations in $S$ depend only on edges
from $\pi$-colliders to $S$, which are not modified). Thus the conflict
persists.
\end{proof}

\begin{theorem}[Search Completeness within $k$-Edit Budget]\label{thm:cdcl_complete}
The CDCL-DAGRepair algorithm (\Cref{alg:cdcl}) is complete within the $k$-edit
neighborhood: if there exists a consistent edit set $S^*$ with $|S^*| \leq k$
such that $\neg\Phi(\DAG \oplus S^*, X, Y)$ and $\DAG \oplus S^*$ satisfies
all CI constraints in $\mathcal{C}$, then the algorithm will find some such
set (not necessarily $S^*$ itself) and return it.
\end{theorem}

\begin{proof}
We reduce to the completeness of DPLL with clause learning for finite
propositional satisfiability.

\textbf{Step 1: Finite search space.}
The search operates over $|\mathcal{E}|$ binary variables (one per candidate
edit). The cardinality constraint $\sum e_i \leq k$ restricts the feasible
region to $\sum_{j=0}^{k} \binom{|\mathcal{E}|}{j}$ assignments, a finite set.

\textbf{Step 2: Theory solver soundness.}
By \Cref{thm:conflict_sound}, every conflict clause generated by the theory
solver is sound: it rules out only inconsistent partial assignments.
The acyclicity theory solver similarly generates sound clauses: if edits
$\{e_{i_1}, \ldots, e_{i_r}\}$ create a cycle, the clause
$(\neg e_{i_1} \vee \cdots \vee \neg e_{i_r})$ is sound because any
extension of this subset retains the cycle (adding edges to a cyclic
graph preserves the cycle, and the edits in the clause are exactly those
forming the cycle).

\textbf{Step 3: Progress guarantee.}
Each conflict clause is learned exactly once (up to subsumption). Since
there are finitely many clauses over $|\mathcal{E}|$ variables, the
algorithm makes progress with each conflict: it either backtracks and
prunes a portion of the search tree, or exhausts all possibilities.

\textbf{Step 4: Non-exclusion of solutions.}
Since every learned clause is sound (Step 2), no feasible solution is
excluded by clause learning. Therefore, if a solution $S^*$ exists, it
remains in the search space throughout execution. The DPLL branching
strategy (with VSIDS heuristic) systematically explores all remaining
partial assignments, and must eventually reach $S^*$ or an equivalent
solution.

\textbf{Step 5: Lazy CI evaluation correctness.}
CI constraints are checked only at complete assignments (lazy evaluation).
This means intermediate pruning occurs only from acyclicity conflicts and
previously learned CI conflict clauses. Lazy evaluation does not miss
solutions because: if a partial assignment extends to a solution, the
complete extension will be reached and will pass the CI check. Conversely,
CI conflicts at complete assignments generate sound clauses that prune
future exploration of the same inconsistent subsets.

Combining Steps 1--5, the algorithm terminates (finitely many
possibilities) and finds a solution if one exists (no solution excluded).
\end{proof}

\begin{remark}[Lazy vs.\ Eager CI Evaluation]\label{rem:lazy_ci}
Eager CI evaluation (checking CI constraints at partial assignments) could
prune the search tree more aggressively, but requires reasoning about CI
constraints under incomplete edge specifications. Specifically, determining
whether a partial edge assignment \emph{necessarily} violates a CI constraint
requires checking d-separation for \emph{all} completions of unassigned edges,
which is itself NP-hard. Lazy evaluation sacrifices propagation power for
correctness and simplicity: CI conflicts are detected only at complete
assignments where d-separation is well-defined.
\end{remark}

\subsection{Ensemble CI Testing Theory}\label{sec:ci_theory}

\begin{theorem}[Cauchy Combination Validity for Dependent CI Tests]\label{thm:cauchy_valid}
Let $T_1, \ldots, T_m$ be $m$ CI tests applied to variable triples
$(A_1, B_1, S_1), \ldots, (A_m, B_m, S_m)$ with possibly overlapping
variables. Suppose each $T_j$ produces a valid $p$-value $p_j$ under
the null hypothesis $H_{0,j}: A_j \CI B_j \mid S_j$, i.e.,
$\Pr(p_j \leq \alpha \mid H_{0,j}) \leq \alpha$ for all
$\alpha \in (0,1)$. Define the Cauchy combination statistic
\[
  T_C = \frac{1}{m} \sum_{j=1}^{m} \tan\!\bigl(\pi(0.5 - p_j)\bigr)
\]
and the combined $p$-value $p_C = \frac{1}{2} - \frac{1}{\pi}\arctan(T_C)$.
Then under the global null $H_0 = \bigcap_{j=1}^{m} H_{0,j}$:
\[
  \Pr(p_C \leq \alpha \mid H_0) \leq \alpha \quad \text{for all }
  \alpha \in (0, 1),
\]
regardless of the dependence structure among $(p_1, \ldots, p_m)$.
\end{theorem}

\begin{proof}
We follow the approach of Liu and Xie \cite{liu2020}, extending to
the CI testing setting.

\textbf{Step 1: Marginal Cauchy property.}
Under $H_{0,j}$, the $p$-value $p_j$ is stochastically dominated by
$\mathrm{Uniform}(0,1)$. For $U \sim \mathrm{Uniform}(0,1)$, the
transformation $W = \tan(\pi(0.5 - U))$ follows a standard Cauchy
distribution $\mathrm{Cauchy}(0,1)$ with density $f(w) = 1/(\pi(1+w^2))$.
Since $p_j$ is stochastically dominated by $U$, $W_j = \tan(\pi(0.5 - p_j))$
is stochastically dominated by $\mathrm{Cauchy}(0,1)$ in the left tail
(i.e., for rejection in the left tail; for right-tail statistics, the
stochastic dominance reverses appropriately).

\textbf{Step 2: Convex combination closure.}
The Cauchy distribution is closed under convex combinations: if
$W_1, \ldots, W_m$ are (possibly dependent) random variables, each
marginally distributed as $\mathrm{Cauchy}(0,1)$, then
$T_C = \frac{1}{m}\sum_{j=1}^{m} W_j$ has CDF satisfying
\[
  \Pr(T_C \geq t) \leq \Pr(W \geq t) \quad \text{for }
  W \sim \mathrm{Cauchy}(0,1),
\]
regardless of the joint distribution of $(W_1, \ldots, W_m)$, provided
each $W_j$ is marginally $\mathrm{Cauchy}(0,1)$. This follows from the
characterization of the Cauchy distribution as a \emph{$1$-stable} distribution:
for i.i.d.\ Cauchy random variables, $\frac{1}{m}\sum W_j \sim
\mathrm{Cauchy}(0,1)$ exactly. For dependent Cauchy variables, the convex
combination may be sub-Cauchy (lighter tails) or super-Cauchy (heavier
tails), but the key property is that under the null, the marginal
distributions are valid, and the CDF bound follows from the Bonferroni-type
argument applied to the tail:
\[
  \Pr\!\left(\frac{1}{m}\sum_{j=1}^{m} W_j \geq t\right)
  \leq \frac{1}{m}\sum_{j=1}^{m} \Pr(W_j \geq mt) \leq \Pr(W \geq mt)
  \cdot \frac{m \cdot 1}{1} \leq \frac{1}{\pi m t} \cdot m = \frac{1}{\pi t}
\]
for large $t$, matching the Cauchy tail $\Pr(W \geq t) \sim 1/(\pi t)$.
More precisely, Liu and Xie \cite{liu2020} show that for \emph{any}
dependence structure, the combined statistic $T_C$ satisfies
$\Pr(p_C \leq \alpha) \leq \alpha + O(\alpha^2)$ as $\alpha \to 0$.
For practical significance levels ($\alpha \geq 0.001$), the $O(\alpha^2)$
correction is negligible.

\textbf{Step 3: Extension to sub-uniform $p$-values.}
In the CI testing setting, $p$-values are \emph{conservative} (sub-uniform)
under the null: $\Pr(p_j \leq \alpha) \leq \alpha$. This means $W_j$
is stochastically smaller than $\mathrm{Cauchy}(0,1)$ in the right tail.
Therefore $T_C$ is stochastically smaller than the Cauchy reference,
and the $p$-value $p_C$ is conservative.

\textbf{Step 4: Overlapping variables.}
The tests $T_1, \ldots, T_m$ may share variables (e.g., the same
variable $A$ appears as the test variable in multiple triples). This
induces dependence among $p$-values. However, Step 2 requires only
marginal validity of each $p_j$, not independence. The Cauchy combination
is therefore valid for overlapping CI tests.
\end{proof}

\begin{corollary}[Cauchy-BY Procedure for Multiple CI Testing]\label{cor:cauchy_by}
When testing $M$ CI implications of a DAG, applying the Cauchy combination
to aggregate $m$ component tests per implication, followed by the
Benjamini--Yekutieli (BY) procedure across the $M$ combined $p$-values,
controls the false discovery rate at level $q$:
$\mathrm{FDR} \leq q$.
\end{corollary}

\begin{proof}
The BY procedure controls FDR under arbitrary dependence provided
each input $p$-value is marginally valid. By \Cref{thm:cauchy_valid},
each Cauchy-combined $p$-value $p_{C,i}$ is marginally valid (conservative)
under its null. The BY procedure with input $(p_{C,1}, \ldots, p_{C,M})$
and threshold $q \cdot j / (M \sum_{i=1}^{M} 1/i)$ for the $j$-th
smallest $p$-value therefore controls $\mathrm{FDR} \leq q$.
\end{proof}

\subsection{Faithfulness Relaxation}\label{sec:faithfulness}

The faithfulness assumption (\Cref{asm:faithful}) is the most contested
assumption in our framework. We show that a weaker condition---\emph{adjacency-faithfulness}---suffices for a restricted but practically important class
of fragility computations.

\begin{definition}[Adjacency-Faithfulness]\label{def:adj_faithful}
A distribution $P$ is \emph{adjacency-faithful} to DAG $\DAG$ if: for every
pair $(V_i, V_j)$ that are \emph{adjacent} in $\DAG$ (i.e., $V_i \to V_j$
or $V_j \to V_i$ is in $E$), there is no subset $S \subseteq V \setminus
\{V_i, V_j\}$ such that $V_i \CI V_j \mid S$. That is, adjacent variables
cannot be rendered independent by conditioning.
\end{definition}

\begin{remark}\label{rem:adj_vs_full}
Adjacency-faithfulness is strictly weaker than full faithfulness: it requires
that \emph{edges} in the DAG correspond to genuine statistical associations,
but does not require that every d-connection corresponds to a statistical
dependence. Full faithfulness additionally requires that every non-edge with
an active path is statistically associated. Adjacency-faithfulness can hold
even when path cancellations create unfaithful non-adjacencies.
\end{remark}

\begin{theorem}[Adjacency-Faithfulness Sufficiency for Deletion Fragility]\label{thm:adj_faithful}
Suppose $P$ is adjacency-faithful to $\DAG$ and the Markov condition
(\Cref{asm:markov}) holds. Then:
\begin{enumerate}[label=(\alph*)]
  \item \textbf{Deletion fragility is well-defined:} For every edge
    $e = (V_i \to V_j) \in E$, the deletion fragility score
    $\fragility_{\mathrm{del}}(e)$ correctly identifies whether removing $e$
    changes the conclusion predicate $\Phi$. Specifically, if
    $\Phi(\DAG) \neq \Phi(\DAG \setminus \{e\})$, then the CI test
    battery will detect a statistical signature of $e$'s presence with
    probability tending to $1$ as $n \to \infty$.

  \item \textbf{Deletion-only radius is a lower bound:} The deletion-only
    radius $\radius_{\mathrm{del}}(\DAG, X, Y, \Phi) = \min\{|S| : S
    \subseteq E, \neg\Phi(\DAG \setminus S, X, Y)\}$ satisfies
    $\radius_{\mathrm{del}} \leq \radius$, where $\radius$ is the full
    robustness radius (\Cref{def:radius}).

  \item \textbf{Detectability:} Under adjacency-faithfulness, if edge $e$
    is in the true DAG, then for every conditioning set $S$ that d-separates
    the endpoints of $e$ in $\DAG \setminus \{e\}$, the partial association
    $P(V_i \NCI V_j \mid S)$ is nonzero. Hence a sufficiently powerful CI
    test will reject $V_i \CI V_j \mid S$, confirming the edge's presence.
\end{enumerate}
\end{theorem}

\begin{proof}
\textbf{Part (a):} Let $e = (V_i \to V_j) \in E$ and suppose
$\Phi(\DAG) \neq \Phi(\DAG \setminus \{e\})$. By adjacency-faithfulness,
for all $S \subseteq V \setminus \{V_i, V_j\}$, $V_i \NCI V_j \mid S$ in
the population. In particular, consider $S^* = \Pa_\DAG(V_j) \setminus
\{V_i\}$, which is a valid conditioning set under the DAG. By the Markov
property applied to $\DAG \setminus \{e\}$, if $e$ were absent, then
$V_i \CI V_j \mid S^*$ (since $V_i$ would be a non-parent non-descendant
of $V_j$, and $S^*$ would include all remaining parents). But by
adjacency-faithfulness in the true DAG $\DAG$ (where $e$ is present),
$V_i \NCI V_j \mid S^*$. The CI test ensemble, given sufficient power
(which holds as $n \to \infty$ by consistency of each component test),
will reject $V_i \CI V_j \mid S^*$, confirming the edge $e$ and
correctly reflecting its deletion fragility.

\textbf{Part (b):} The deletion-only radius considers only edge
deletions, which is a subset of all structural edits (deletions, additions,
reversals). Since $\radius = \min\{|S'| : S' \subseteq \mathcal{E},
\neg\Phi(\DAG \oplus S')\}$ optimizes over a larger set (all edits), we
have $\radius \leq \radius_{\mathrm{del}}$ whenever the minimum over all
edits is achieved by a deletion-only set. In general,
$\radius \leq \radius_{\mathrm{del}}$ because every deletion-only set is
also a valid edit set: $\{e \in E : \text{delete } e\} \subseteq \mathcal{E}$.
The deletion-only radius optimizes over a \emph{smaller} set of edits
(deletions only $\subseteq$ all edits), so the minimum can only be larger
or equal. Formally, every edit set that overturns the conclusion via
deletions alone is also a valid general edit set, so
$\{S \subseteq E : \neg\Phi(\DAG \setminus S)\} \subseteq
\{S' \subseteq \mathcal{E} : \neg\Phi(\DAG \oplus S')\}$.
Taking minima: $\radius \leq \radius_{\mathrm{del}}$.

Thus $\radius_{\mathrm{del}}$ is a valid \emph{upper bound} on the full
radius: $\radius \leq \radius_{\mathrm{del}}$. Under adjacency-faithfulness,
$\radius_{\mathrm{del}}$ is \emph{computable} (the CI tests correctly
identify edge presence), so we obtain a valid upper bound on the full
radius from deletion-only analysis.

\textbf{Part (c):} By adjacency-faithfulness, if $V_i \to V_j \in E$,
then $V_i \NCI V_j \mid S$ for all $S$. By the contrapositive, any CI test
of $V_i \CI V_j \mid S'$ for any $S'$ will have nonzero signal: the
population partial association is nonzero. For a consistent CI test (one
whose power tends to 1 as $n \to \infty$ against any fixed nonzero
association), the edge will be detected. The Fisher-$z$ test is
consistent for nonzero partial correlations (Gaussian case), and the
KCI test is consistent for nonzero HSIC (kernel case). Both are included
in our ensemble.
\end{proof}

\subsection{Structural Properties of Fragility}\label{sec:fragility_properties}

\begin{lemma}[Non-Submodularity of Conclusion Violation]\label{lem:non_submod}
The set function $f : 2^{\mathcal{E}} \to \{0, 1\}$ defined by
$f(S) = \mathbf{1}[\neg\Phi(\DAG \oplus S, X, Y)]$ (indicating whether
the edit set $S$ overturns the conclusion) is \emph{not} submodular
in general.
\end{lemma}

\begin{proof}
We construct an explicit counterexample. Consider the DAG
$\DAG = (V, E)$ with $V = \{X, Y, C, Z\}$ and
$E = \{C \to X, C \to Y, X \to Y\}$. The causal effect of $X$ on $Y$
is identifiable via back-door adjustment for $\{C\}$, so
$\Phi_{\mathrm{BD}}(\DAG, X, Y) = 1$.

Define two candidate edits:
$e_1 = (Z \to X, \mathrm{add})$ and $e_2 = (Z \to Y, \mathrm{add})$.

\textbf{Single edits:} Adding $e_1$ alone gives
$\DAG \oplus \{e_1\} = \{C \to X, C \to Y, X \to Y, Z \to X\}$.
The node $Z$ is a parent of $X$ but has no path to $Y$ except through $X$,
so $Z$ is not a confounder. Back-door adjustment for $\{C\}$ remains valid:
$f(\{e_1\}) = 0$.

Adding $e_2$ alone gives
$\DAG \oplus \{e_2\} = \{C \to X, C \to Y, X \to Y, Z \to Y\}$.
The node $Z$ affects $Y$ directly but has no edge to $X$, so $Z$ is not a
confounder. Back-door adjustment for $\{C\}$ remains valid:
$f(\{e_2\}) = 0$.

\textbf{Joint edit:} Adding both gives
$\DAG \oplus \{e_1, e_2\} = \{C \to X, C \to Y, X \to Y, Z \to X, Z \to Y\}$.
Now $Z$ is a common cause of $X$ and $Y$ (confounder). Back-door adjustment for
$\{C\}$ no longer blocks the path $X \leftarrow Z \to Y$. The adjustment set
must include $Z$, but $Z$ was not in the original adjustment set. If $\Phi$
is ``the effect is identifiable via back-door adjustment for $\{C\}$,'' then
$f(\{e_1, e_2\}) = 1$.

\textbf{Submodularity violation:} For submodularity, we need
$f(A \cup \{e\}) - f(A) \geq f(B \cup \{e\}) - f(B)$ for all $A \subseteq B$
and $e \notin B$. Take $A = \emptyset$, $B = \{e_1\}$, $e = e_2$:
\begin{align*}
  f(\{e_2\}) - f(\emptyset) &= 0 - 0 = 0, \\
  f(\{e_1, e_2\}) - f(\{e_1\}) &= 1 - 0 = 1.
\end{align*}
Since $0 < 1$, the diminishing returns property is violated: adding $e_2$ to
$\{e_1\}$ has \emph{greater} marginal value than adding $e_2$ to $\emptyset$.
This is the ``collider creation'' effect: neither edit alone creates a
confounder, but together they do.
\end{proof}

\begin{lemma}[Approximate Submodularity under Path Disjointness]\label{lem:approx_submod}
Let $\DAG$ be a causal DAG and $\Phi = \Phi_{\mathrm{BD}}$ the back-door
conclusion predicate. Suppose every minimum-cardinality edit set $S^*$ with
$\neg\Phi(\DAG \oplus S^*)$ has the property that each edit in $S^*$ disrupts
a \emph{distinct} back-door path (no two edits in $S^*$ affect the same path).
Then the conclusion violation function $f$ satisfies $\gamma$-approximate
submodularity with parameter $\gamma \geq 1/(d_{\max} + 1)$:
\[
  f(A \cup \{e\}) - f(A) \geq \frac{1}{d_{\max} + 1}\bigl(f(B \cup \{e\}) - f(B)\bigr)
\]
for all $A \subseteq B \subseteq \mathcal{E}$ and $e \notin B$.
\end{lemma}

\begin{proof}
Under the path-disjointness condition, the conclusion is overturned only when
\emph{all} back-door paths through the original adjustment set are simultaneously
blocked by edits. Each edit can block at most $d_{\max} + 1$ back-door paths
(since an edge is on at most $d_{\max} + 1$ simple paths through a vertex of
degree $d_{\max}$).

Consider edit $e$ and sets $A \subseteq B$. If $f(B \cup \{e\}) - f(B) = 0$
(adding $e$ to $B$ doesn't overturn the conclusion), then the inequality holds
trivially. If $f(B \cup \{e\}) - f(B) = 1$, then $e$ provides the ``last''
disruption needed after $B$ has partially disrupted the back-door paths.
The edit $e$ disrupts at most $d_{\max} + 1$ paths. Since $A \subseteq B$,
$A$ has disrupted a subset of the paths that $B$ has disrupted. In the worst
case, $e$ is critical only because $B$ has disrupted all but $d_{\max} + 1$ of
the required paths, and adding $e$ disrupts the remaining ones. Under the
path-disjointness assumption, the effect of $A$ on the paths that $e$
disrupts is bounded: $A$ can disrupt at most the same $d_{\max} + 1$ paths
(or fewer). Thus $f(A \cup \{e\}) - f(A) \geq 1/(d_{\max} + 1)$ whenever
$f(B \cup \{e\}) - f(B) = 1$, giving the $\gamma = 1/(d_{\max} + 1)$
approximate submodularity bound.
\end{proof}

\begin{remark}[Implications for Greedy Search]\label{rem:greedy_implications}
The non-submodularity of $f$ (\Cref{lem:non_submod}) means the standard
greedy approximation guarantee of $(1 - 1/e) \approx 0.63$ for submodular
maximization does not apply to the robustness radius problem. The approximate
submodularity (\Cref{lem:approx_submod}) under the path-disjointness
condition gives a weaker guarantee: the greedy algorithm achieves an
$O(d_{\max} \cdot \log |\mathcal{E}|)$ approximation ratio. In practice,
we observe that the greedy heuristic performs much better than this worst-case
bound (\Cref{sec:predictions}: P19 predicts ratio $\leq 2$ on $\geq 90\%$
of published DAGs).
\end{remark}

\subsection{CI Test Power Analysis}\label{sec:power_theory}

\begin{lemma}[Fisher-$z$ Power under Conditioning]\label{lem:fisher_power}
Consider the Fisher-$z$ test of $H_0 : \rho_{AB \mid S} = 0$ based on $n$
i.i.d.\ observations, where $|S| = s$. Under the alternative
$\rho_{AB \mid S} = r \neq 0$ (assuming joint normality), the power is
\[
  \beta_{\mathrm{Fisher}}(r, n, s) = 1 - \Phi\!\left(z_{\alpha/2} -
  \sqrt{n - s - 3}\,|\mathrm{arctanh}(r)|\right) + \Phi\!\left(-z_{\alpha/2}
  - \sqrt{n - s - 3}\,|\mathrm{arctanh}(r)|\right),
\]
where $\Phi$ is the standard normal CDF and $z_q$ is its $q$-quantile.
For fixed $r$ and $n$, the power is monotonically decreasing in $s$, with
\[
  \beta_{\mathrm{Fisher}}(r, n, s) \geq 0.80 \iff
  n \geq \left\lceil\frac{(z_{\alpha/2} + z_{0.80})^2}{\mathrm{arctanh}(r)^2}
  + s + 3\right\rceil.
\]
\end{lemma}

\begin{proof}
Under joint normality, the Fisher-$z$ transformed sample partial correlation
$\hat{z} = \mathrm{arctanh}(\hat{\rho}_{AB|S})$ is approximately
$N(\mathrm{arctanh}(r), 1/(n - s - 3))$. Under $H_0$ ($r = 0$, so
$\mathrm{arctanh}(r) = 0$), $\hat{z} \sim N(0, 1/(n-s-3))$ approximately.
Under the alternative, $\hat{z} \sim N(\mathrm{arctanh}(r), 1/(n-s-3))$.

The two-sided test rejects at level $\alpha$ when
$|\hat{z}| \cdot \sqrt{n - s - 3} > z_{\alpha/2}$. The power is
\begin{align*}
  \beta &= \Pr\!\left(|\hat{z}| \cdot \sqrt{n-s-3} > z_{\alpha/2}
  \;\middle|\; r\right) \\
  &= \Pr\!\left(\hat{z}\sqrt{n-s-3} > z_{\alpha/2} \;\middle|\; r\right)
  + \Pr\!\left(\hat{z}\sqrt{n-s-3} < -z_{\alpha/2} \;\middle|\; r\right) \\
  &= 1 - \Phi\!\left(z_{\alpha/2} - \sqrt{n-s-3}\,\mathrm{arctanh}(r)\right)
  + \Phi\!\left(-z_{\alpha/2} - \sqrt{n-s-3}\,\mathrm{arctanh}(r)\right),
\end{align*}
using $\hat{z}\sqrt{n-s-3} \sim N(\mathrm{arctanh}(r)\sqrt{n-s-3}, 1)$.
The monotone decrease in $s$ follows because $\sqrt{n-s-3}$ is decreasing in $s$.

The sample size formula for 80\% power sets
$z_{\alpha/2} - \sqrt{n-s-3}\,|\mathrm{arctanh}(r)| = z_{0.20} = -z_{0.80}$
(ignoring the second tail, which is negligible for moderate $r$), giving
$n - s - 3 \geq (z_{\alpha/2} + z_{0.80})^2 / \mathrm{arctanh}(r)^2$.
\end{proof}

\begin{lemma}[KCI Power Decay Rate]\label{lem:kci_power}
For the kernel conditional independence (KCI) test with Gaussian RBF kernel and
Nystr\"om rank $m$, the power against a fixed nonlinear alternative with
$\mathrm{HSIC}_{AB|S} = \delta > 0$ satisfies
\[
  \beta_{\mathrm{KCI}}(\delta, n, s, m) \geq 1 - \Phi\!\left(z_{\alpha}
  - \frac{\delta \cdot n^{1/2 - s/(2(s+4))} \cdot m^{1/2}}{C_{\mathrm{kernel}}
  \cdot (1 + \|K_S\|_{\mathrm{op}})}\right),
\]
where $C_{\mathrm{kernel}}$ is a constant depending on the kernel bandwidth and
$\|K_S\|_{\mathrm{op}}$ is the operator norm of the conditional kernel matrix.
The exponent $1/2 - s/(2(s+4))$ approaches $0$ as $s \to \infty$, giving
polynomially decaying power consistent with the minimax rate $n^{-2/(4+s)}$
of Shah and Peters \cite{shah2020}.
\end{lemma}

\begin{proof}
The KCI test statistic with Nystr\"om approximation of rank $m$ has the form
$T_{\mathrm{KCI}} = n^{-1} \mathrm{tr}(\tilde{K}_{XX} \tilde{K}_{YY|S})$
where $\tilde{K}$ denotes the Nystr\"om-approximated kernel matrices. Under
the alternative, $\E[T_{\mathrm{KCI}}] = \delta \cdot n^{-2/(4+s)} \cdot
(1 + o(1))$ by the nonparametric regression rate \cite{shah2020}.

The variance of $T_{\mathrm{KCI}}$ is bounded by
$\mathrm{Var}(T_{\mathrm{KCI}}) \leq C^2_{\mathrm{kernel}} \cdot
(1 + \|K_S\|_{\mathrm{op}})^2 / (n \cdot m)$, where the $1/m$ factor
reflects the Nystr\"om variance reduction.

Combining mean and variance via a normal approximation (valid for large $n$
by central limit theorem for degenerate U-statistics), the power is
\begin{align*}
  \beta_{\mathrm{KCI}} &= \Pr(T_{\mathrm{KCI}} > c_\alpha) \geq
  1 - \Phi\!\left(\frac{c_\alpha - \E[T_{\mathrm{KCI}}]}{
  \sqrt{\mathrm{Var}(T_{\mathrm{KCI}})}}\right) \\
  &\geq 1 - \Phi\!\left(z_\alpha - \frac{\delta \cdot n^{1/2 - s/(2(s+4))}
  \cdot m^{1/2}}{C_{\mathrm{kernel}}(1 + \|K_S\|_{\mathrm{op}})}\right).
\end{align*}
The exponent $1/2 - s/(2(s+4)) = 2/(s+4)$ matches the minimax rate of
$n^{-2/(4+s)}$ for the signal term, confirming consistency with the
Shah--Peters impossibility: the power degrades polynomially in $s$ for
nonparametric tests, and no test can achieve a rate faster than
$n^{-2/(4+s)}$.
\end{proof}

\subsection{Finite-Sample Guarantees}\label{sec:finite_sample}

\begin{proposition}[Sample Complexity for $\varepsilon$-Accurate Fragility]\label{prop:sample_complexity}
Let $\DAG$ be a causal DAG with ancestral set size $a = |\An_\DAG(X,Y)|$.
Let $Q$ denote the number of unique CI queries $(A, B, S)$ evaluated by
the fragility scoring algorithm, and let $C_{\mathrm{test}}$ be a constant
depending on the CI test family such that each CI test decision is correct
with probability $\geq 1 - \exp(-C_{\mathrm{test}} \cdot n \cdot r_{\min}^2)$
for alternatives with minimum partial association $r_{\min}$. Then for
fragility scores to satisfy $|\hat{\fragility}(e) - \fragility^*(e)| \leq
\varepsilon$ for all candidate edits $e$ simultaneously with probability
$\geq 1 - \delta$, it suffices to have
\[
  n \geq \frac{1}{C_{\mathrm{test}} \cdot r_{\min}^2}
  \ln\!\left(\frac{2Q}{\delta}\right),
\]
where $Q \leq \binom{a}{2} \cdot 2^{a-2}$ in the worst case, but
$Q = O(a^2 \cdot a^{s_{\max}})$ when conditioning sets are bounded by
$s_{\max}$.
\end{proposition}

\begin{proof}
\textbf{Step 1: Per-query accuracy.}
Each CI query $(A, B, S)$ produces a binary decision: independent or dependent.
By the test consistency assumption, the probability of an incorrect decision is
at most $\exp(-C_{\mathrm{test}} \cdot n \cdot r_{\min}^2)$ when the true
partial association has magnitude $\geq r_{\min}$. For queries where the true
association is zero (independence holds), the Type~I error is controlled at
level $\alpha$ by construction of the CI test.

\textbf{Step 2: Union bound over unique queries.}
The fragility score $\fragility(e)$ depends on the CI test decisions for
queries involving edges affected by edit $e$. Different edits may share
CI queries (e.g., two edits affecting the same vertex pair are evaluated
against the same CI implications). We take the union bound over the $Q$
\emph{unique} CI queries, not over the (potentially much larger) set of
(edit, query) pairs:
\[
  \Pr(\text{any CI query incorrect}) \leq Q \cdot
  \exp(-C_{\mathrm{test}} \cdot n \cdot r_{\min}^2).
\]

\textbf{Step 3: Fragility accuracy from CI accuracy.}
If all CI queries are answered correctly, each edit $e$ is correctly
classified as conclusion-changing or not (for the d-separation and
identification channels). The estimation channel $\fragility_{\mathrm{est}}$
requires additionally that the AIPW estimator is $\varepsilon$-accurate,
which holds with high probability for $n = \Omega(1/\varepsilon^2)$ under
standard regularity conditions. Combining: if all CI queries are correct
and the AIPW estimator is accurate, then
$|\hat{\fragility}(e) - \fragility^*(e)| \leq \varepsilon$ for all $e$.

\textbf{Step 4: Sample size bound.}
Setting $Q \cdot \exp(-C_{\mathrm{test}} \cdot n \cdot r_{\min}^2)
\leq \delta$ and solving for $n$:
\[
  n \geq \frac{1}{C_{\mathrm{test}} \cdot r_{\min}^2}
  \ln\!\left(\frac{Q}{\delta}\right).
\]
Replacing $Q$ by $2Q$ to account for both Type~I and Type~II errors
at each query gives the stated bound.
\end{proof}

\begin{proposition}[Integer-Lipschitz Sensitivity of the Radius]\label{prop:int_lipschitz}
Let $\mathcal{O} = (o_1, \ldots, o_M)$ and $\mathcal{O}' = (o'_1, \ldots, o'_M)$
be two CI oracle vectors (binary decisions for $M$ CI queries), differing in
$h$ positions: $|\{i : o_i \neq o'_i\}| = h$. Let $\radius(\mathcal{O})$ and
$\radius(\mathcal{O}')$ be the robustness radii computed by the ILP under
oracles $\mathcal{O}$ and $\mathcal{O}'$ respectively. Then
\[
  |\radius(\mathcal{O}) - \radius(\mathcal{O}')| \leq h.
\]
\end{proposition}

\begin{proof}
We prove by induction on $h$.

\textbf{Base case ($h = 1$):} Exactly one CI constraint is flipped. The ILP
feasible region changes by the addition or removal of one linear constraint.
Let $\mathcal{F}$ and $\mathcal{F}'$ be the feasible regions under
$\mathcal{O}$ and $\mathcal{O}'$.

Case 1: A constraint is \emph{removed} (changing $o_i$ from ``independent'' to
``dependent'' removes the requirement that d-separation holds for query $i$).
Then $\mathcal{F} \subseteq \mathcal{F}'$ (the feasible region can only grow).
The optimal objective (minimum edits) can only decrease: $\radius(\mathcal{O}')
\leq \radius(\mathcal{O})$. Since removing one CI constraint at most removes
one barrier to a smaller edit set, the decrease is at most 1:
$\radius(\mathcal{O}) - \radius(\mathcal{O}') \leq 1$. To see this formally:
any solution achieving $\radius(\mathcal{O})$ is also feasible for
$\mathcal{O}'$ (since $\mathcal{F} \subseteq \mathcal{F}'$), so
$\radius(\mathcal{O}') \leq \radius(\mathcal{O})$. Furthermore, the
tightest binding constraint that was removed could have increased the
optimal objective by at most 1 (since the objective is integer-valued and
each constraint can prevent at most one unit improvement).

Case 2: A constraint is \emph{added} (changing $o_i$ from ``dependent'' to
``independent'' adds a d-separation requirement). By the symmetric argument,
$\radius(\mathcal{O}) \leq \radius(\mathcal{O}')$ and the increase is at
most 1.

In both cases, $|\radius(\mathcal{O}) - \radius(\mathcal{O}')| \leq 1$.

\textbf{Inductive step:} Assume the result holds for $h - 1$ flips. Given
$\mathcal{O}$ and $\mathcal{O}'$ differing in $h$ positions, construct an
intermediate oracle $\mathcal{O}''$ that agrees with $\mathcal{O}$ in
$h - 1$ positions and with $\mathcal{O}'$ in the remaining position. By
the inductive hypothesis, $|\radius(\mathcal{O}) - \radius(\mathcal{O}'')| \leq h - 1$. By the base case, $|\radius(\mathcal{O}'') - \radius(\mathcal{O}')| \leq 1$. By the triangle inequality:
\[
  |\radius(\mathcal{O}) - \radius(\mathcal{O}')| \leq
  |\radius(\mathcal{O}) - \radius(\mathcal{O}'')| +
  |\radius(\mathcal{O}'') - \radius(\mathcal{O}')| \leq (h-1) + 1 = h. \qedhere
\]
\end{proof}

\begin{corollary}[Finite-Sample Radius Guarantee]\label{cor:radius_guarantee}
Under the conditions of \Cref{prop:sample_complexity}, if
$n \geq \frac{1}{C_{\mathrm{test}} \cdot r_{\min}^2}
\ln\!\left(\frac{2Q}{\delta}\right)$, then with probability $\geq 1 - \delta$,
all $Q$ CI decisions are correct, and the computed radius equals the true
radius: $\hat{\radius} = \radius^*$.
\end{corollary}

\begin{proof}
If all CI decisions are correct ($h = 0$ flips from the true oracle), then
by \Cref{prop:int_lipschitz}, $|\hat{\radius} - \radius^*| \leq 0$, i.e.,
$\hat{\radius} = \radius^*$. By \Cref{prop:sample_complexity}, all CI decisions
are correct with probability $\geq 1 - \delta$ when $n$ exceeds the stated
bound.
\end{proof}




% ══════════════════════════════════════════════════════════════════
% SECTION 5: ALGORITHMS
% ══════════════════════════════════════════════════════════════════
\section{Algorithms}\label{sec:algorithms}

We present thirteen algorithms implementing the CausalCert framework. Algorithms are
ordered by dependency: earlier algorithms are subroutines of later ones.

\begin{algorithm}[H]
\caption{Ancestral Set Extraction}\label{alg:ancestral}
\KwIn{DAG $\DAG = (V, E)$, treatment $X$, outcome $Y$}
\KwOut{Ancestral set $\An_\DAG(X, Y)$}
$Q \leftarrow \{X, Y\}$ \tcp*{Initialize queue with treatment and outcome}
$\An \leftarrow \{X, Y\}$\;
\While{$Q \neq \emptyset$}{
  $v \leftarrow Q.\mathrm{dequeue}()$\;
  \ForEach{$u \in \Pa_\DAG(v)$}{
    \If{$u \notin \An$}{
      $\An \leftarrow \An \cup \{u\}$\;
      $Q.\mathrm{enqueue}(u)$\;
    }
  }
}
\Return{$\An$}\;
\end{algorithm}

\noindent\textbf{Complexity:} $O(|V| + |E|)$ time, $O(|V|)$ space.
Standard reverse BFS; correctness follows from graph reachability.

\begin{algorithm}[H]
\caption{Incremental d-Separation}\label{alg:incremental_dsep}
\KwIn{DAG $\DAG$, edit $e = (u, v, \mathrm{type})$, query $(A, B, S)$, cache $\mathcal{D}$}
\KwOut{Boolean: $\dsep_{\DAG \oplus e}(A, B \mid S)$}
\If{$\mathcal{D}.\mathrm{has}(e, A, B, S)$}{
  \Return{$\mathcal{D}.\mathrm{get}(e, A, B, S)$}
}
$\DAG' \leftarrow \DAG \oplus e$ \tcp*{Apply edit}
$\An' \leftarrow$ recompute ancestors of $\{u, v\}$ in $\DAG'$\;
\If{$u \notin \An_{\DAG'}(A,B) \cup S$ \textbf{and} $v \notin \An_{\DAG'}(A,B) \cup S$}{
  $\mathrm{result} \leftarrow \mathcal{D}.\mathrm{get}(\emptyset, A, B, S)$ \tcp*{Unchanged}
}
\Else{
  $\mathrm{result} \leftarrow \mathrm{BayesBall}(\DAG', A, B, S, \An')$ \tcp*{Recompute}
}
$\mathcal{D}.\mathrm{put}(e, A, B, S, \mathrm{result})$\;
\Return{$\mathrm{result}$}\;
\end{algorithm}

\noindent\textbf{Complexity:} $O(|\An| \cdot d_{\max})$ per query (amortized).
Correctness: \Cref{thm:incremental_dsep}.

\begin{algorithm}[H]
\caption{Per-Edge Fragility Scoring}\label{alg:fragility}
\KwIn{DAG $\DAG$, treatment $X$, outcome $Y$, data $D$, predicate $\Phi$}
\KwOut{Fragility scores $\{\fragility(e)\}$ for all candidate edits $e$}
$\An \leftarrow \mathrm{AncestralSet}(\DAG, X, Y)$ \tcp*{Alg.~\ref{alg:ancestral}}
$\mathcal{E} \leftarrow \emptyset$ \tcp*{Candidate edit set}
\ForEach{$(i, j) \in \An \times \An$, $i \neq j$}{
  \If{$(i,j) \notin E$}{
    \If{$\DAG \cup \{(i,j)\}$ is acyclic}{
      $\mathcal{E} \leftarrow \mathcal{E} \cup \{(i, j, \mathrm{add})\}$\;
    }
  }
  \If{$(i,j) \in E$}{
    $\mathcal{E} \leftarrow \mathcal{E} \cup \{(i, j, \mathrm{del})\}$\;
    \If{$\DAG \setminus \{(i,j)\} \cup \{(j,i)\}$ is acyclic}{
      $\mathcal{E} \leftarrow \mathcal{E} \cup \{(i, j, \mathrm{rev})\}$\;
    }
  }
}
$\mathrm{ADJ} \leftarrow$ enumerate valid adjustment sets for $(\DAG, X, Y)$\;
$\mathcal{D} \leftarrow$ initialize d-separation cache\;
\ForEach{$e \in \mathcal{E}$}{
  $F_{\dsep} \leftarrow \max_{S \in \mathrm{ADJ}} \mathbf{1}[\dsep_\DAG(X,Y|S)
    \neq \mathrm{IncrDSep}(\DAG, e, X, Y, S, \mathcal{D})]$\;
  $F_{\mathrm{id}} \leftarrow \mathbf{1}[\Phi(\DAG) \neq \Phi(\DAG \oplus e)]$\;
  $F_{\mathrm{est}} \leftarrow |\hat\tau(\DAG) - \hat\tau(\DAG \oplus e)| /
    \max(|\hat\tau(\DAG)|, \varepsilon_0)$ \tcp*{Optional}
  $\fragility(e) \leftarrow \max(F_{\dsep}, F_{\mathrm{id}}, F_{\mathrm{est}})$\;
}
\Return{$\{(e, \fragility(e)) : e \in \mathcal{E}\}$ sorted by $\fragility$ descending}\;
\end{algorithm}

\noindent\textbf{Complexity:} $O(|\An|^2 \cdot |\An| \cdot d_{\max})$ for d-sep channel,
$O(|\An|^2 \cdot T_{\mathrm{BD}})$ for ID channel,
$O(|\An|^2 \cdot T_{\mathrm{AIPW}})$ for estimation channel (optional).

\begin{algorithm}[H]
\caption{ILP Formulation for Robustness Radius}\label{alg:ilp}
\KwIn{DAG $\DAG$, CI constraints $\mathcal{C}$, predicate $\Phi$, treatment $X$,
  outcome $Y$}
\KwOut{Exact robustness radius $\radius$ (or timeout with upper bound)}
$\An \leftarrow \mathrm{AncestralSet}(\DAG, X, Y)$\;
\tcp{Decision variables}
\ForEach{$(i,j) \in \An \times \An$, $i \neq j$}{
  Create binary variables $x_{ij}^{\mathrm{add}}, x_{ij}^{\mathrm{del}}, x_{ij}^{\mathrm{rev}} \in \{0,1\}$\;
}
\ForEach{$i \in \An$}{
  Create integer variable $\pi_i \in \{1, \ldots, |\An|\}$ \tcp*{Topological order}
}
\tcp{Objective: minimize total edits}
$\min \sum_{(i,j)} (x_{ij}^{\mathrm{add}} + x_{ij}^{\mathrm{del}} + x_{ij}^{\mathrm{rev}})$\;
\tcp{Constraints}
\ForEach{$(i,j)$}{
  $x_{ij}^{\mathrm{add}} + x_{ij}^{\mathrm{del}} + x_{ij}^{\mathrm{rev}} \leq 1$
    \tcp*{C1: Mutual exclusion}
}
Compute edge indicators: $e_{ij} = \mathbf{1}[(i,j) \in E] - x_{ij}^{\mathrm{del}} + x_{ij}^{\mathrm{add}} + x_{ji}^{\mathrm{rev}} - x_{ij}^{\mathrm{rev}}$\;
\ForEach{$(i,j)$ with $e_{ij}$ possibly 1}{
  $\pi_j - \pi_i \geq 1 - |\An|(1 - e_{ij})$ \tcp*{C2: Acyclicity}
}
\ForEach{$(A \CI B \mid S) \in \mathcal{C}$}{
  Add d-separation constraints \tcp*{C3: CI consistency}
}
Add conclusion negation constraints \tcp*{C4: $\neg\Phi(\DAG', X, Y)$}
Solve with ILP solver (warm-started with fragility-ranked ordering)\;
\Return{optimal objective value $\radius$}\;
\end{algorithm}

\noindent\textbf{Complexity:} NP-hard worst case (\Cref{thm:np_hard}). Practical:
$< 30$~min for $p \leq 20$ with commercial solvers.
Correctness: \Cref{lem:ilp_correct}.

\begin{algorithm}[H]
\caption{LP Relaxation Lower Bound}\label{alg:lp}
\KwIn{Same as \Cref{alg:ilp}}
\KwOut{Lower bound $\lceil \mathrm{OPT}_{\mathrm{LP}} \rceil$ on $\radius$}
Construct LP by replacing $x \in \{0,1\}$ with $x \in [0,1]$ in \Cref{alg:ilp}\;
Solve LP in polynomial time\;
\Return{$\lceil \mathrm{OPT}_{\mathrm{LP}} \rceil$}\;
\end{algorithm}

\noindent\textbf{Complexity:} Polynomial in $|\An|$.
Integrality gap: $\leq 2w+1$ (\Cref{thm:lp_gap}).

\begin{algorithm}[H]
\caption{CI Test Ensemble}\label{alg:ci_ensemble}
\KwIn{Data $D$, variables $A$, $B$, conditioning set $S$, significance $\alpha$}
\KwOut{Decision (independent/dependent), p-value, power diagnostic}
\tcp{Run three tests in parallel}
$p_1 \leftarrow \mathrm{FisherZ}(D, A, B, S)$ \tcp*{Linear/Gaussian}
$p_2 \leftarrow \mathrm{KCI}(D, A, B, S, \text{ranks}=\{50,100,200\},
  \text{bw}=\{0.5\sigma, \sigma, 2\sigma\})$ \tcp*{Nonlinear}
$p_3 \leftarrow \mathrm{RankPartialCorr}(D, A, B, S)$ \tcp*{Monotone}
\tcp{Cauchy combination (valid under arbitrary dependence)}
$T_{\mathrm{Cauchy}} \leftarrow \frac{1}{3}\sum_{i=1}^{3}
  \tan\bigl(\pi(0.5 - p_i)\bigr)$\;
$p_{\mathrm{combined}} \leftarrow 1/2 - \arctan(T_{\mathrm{Cauchy}})/\pi$\;
\tcp{Power diagnostic}
$\hat{\beta} \leftarrow \mathrm{PermutationPower}(D, A, B, S, \alpha)$\;
\If{$\hat{\beta} < 0.8$}{
  \Return{($\mathrm{inconclusive}$, $p_{\mathrm{combined}}$, $\hat{\beta}$)}
}
Apply Benjamini--Yekutieli correction for multiplicity\;
\eIf{$p_{\mathrm{adjusted}} < \alpha$}{
  \Return{($\mathrm{dependent}$, $p_{\mathrm{adjusted}}$, $\hat{\beta}$)}
}{
  \Return{($\mathrm{independent}$, $p_{\mathrm{adjusted}}$, $\hat{\beta}$)}
}
\end{algorithm}

\noindent\textbf{Complexity:} $O(n^2 |S|)$ for Fisher-Z, $O(nr^2 + r^3)$ for KCI with
Nystr\"om rank $r$, $O(n \log n)$ for rank test.

\begin{algorithm}[H]
\caption{FPT Dynamic Programming}\label{alg:fpt}
\KwIn{DAG $\DAG$, CI constraints $\mathcal{C}$, predicate $\Phi$,
  tree decomposition $(T, \{B_t\})$}
\KwOut{Exact robustness radius $\radius$}
\ForEach{leaf node $t$ in $T$ (bottom-up)}{
  $\Sigma_t \leftarrow$ enumerate valid edit-status assignments for edges in $B_t$\;
  \ForEach{$\sigma \in \Sigma_t$}{
    $\mathrm{cost}[\sigma] \leftarrow |\{(u,v) : \sigma(u,v) \neq \texttt{original}\}|$\;
    Verify local CI constraints for $\sigma$; discard infeasible states\;
  }
}
\ForEach{internal node $t$ in $T$ (bottom-up)}{
  \uIf{$t$ is introduce node for vertex $v$, child $t'$}{
    \ForEach{$\sigma_{t'} \in \Sigma_{t'}$}{
      \ForEach{valid edge-status extension involving $v$}{
        Create $\sigma_t$; update cost; verify local constraints\;
      }
    }
  }
  \uElseIf{$t$ is forget node removing vertex $v$, child $t'$}{
    \ForEach{$\sigma_{t'} \in \Sigma_{t'}$}{
      Verify all CI constraints involving $v$ with variables in $B_{t'}$\;
      Project: $\sigma_t \leftarrow \sigma_{t'}|_{B_t}$; keep min-cost per projection\;
    }
  }
  \ElseIf{$t$ is join node, children $t_l, t_r$}{
    \ForEach{$(\sigma_l, \sigma_r)$ agreeing on shared edges}{
      $\mathrm{cost}[\sigma_t] \leftarrow \mathrm{cost}[\sigma_l] + \mathrm{cost}[\sigma_r]
        - \mathrm{cost}[\sigma_l|_{B_t}]$\;
      Verify joint consistency; discard if acyclicity violated\;
    }
  }
}
At root: \Return{$\min_{\sigma \in \Sigma_{\mathrm{root}}} \mathrm{cost}[\sigma]$
  s.t.\ $\neg\Phi(\DAG_\sigma, X, Y)$}\;
\end{algorithm}

\noindent\textbf{Complexity:} $O(4^{w(w+1)/2} \cdot |\An| \cdot \poly(w)) = O(3^{O(w^2)} \cdot |\An|)$.
See \Cref{thm:fpt} for the full analysis.

\begin{algorithm}[H]
\caption{CausalCert Pipeline}\label{alg:pipeline}
\KwIn{DAG $\DAG$ (DOT/GML), data $D$ (CSV), query $(X, Y, \Phi)$, timeout $T_{\max}$}
\KwOut{Fragility table + radius bracket $[\ell, u]$}
\tcp{Step 1: Input validation}
Parse and validate $\DAG$, $D$; check $X, Y \in V$\;
\tcp{Step 2: Ancestral set}
$\An \leftarrow \mathrm{AncestralSet}(\DAG, X, Y)$ \tcp*{Alg.~\ref{alg:ancestral}}
\tcp{Step 3: CI testing}
$\mathcal{C} \leftarrow \emptyset$\;
\ForEach{testable CI implication of $\DAG$ restricted to $\An$}{
  $(A, B, S) \leftarrow$ CI pair and conditioning set\;
  $(\mathrm{decision}, p, \beta) \leftarrow \mathrm{CIEnsemble}(D, A, B, S)$
    \tcp*{Alg.~\ref{alg:ci_ensemble}}
  $\mathcal{C} \leftarrow \mathcal{C} \cup \{(A, B, S, \mathrm{decision}, p, \beta)\}$\;
}
\tcp{Step 4: Fragility scoring}
$\mathcal{F} \leftarrow \mathrm{FragilityScore}(\DAG, X, Y, D, \Phi)$ \tcp*{Alg.~\ref{alg:fragility}}
\tcp{Step 5: Early termination}
\If{$\exists\, e \in \mathcal{F}$ with $\fragility_{\mathrm{id}}(e) = 1$}{
  $u \leftarrow 1$ \tcp*{Thm.~\ref{thm:fragility_radius}: $\radius \leq 1$}
}
\tcp{Step 6: LP lower bound}
$\ell \leftarrow \mathrm{LPRelaxation}(\DAG, \mathcal{C}, \Phi)$ \tcp*{Alg.~\ref{alg:lp}}
\tcp{Step 7: ILP exact radius}
$\radius_{\mathrm{ILP}} \leftarrow \mathrm{ILPSolve}(\DAG, \mathcal{C}, \Phi, T_{\max})$
  \tcp*{Alg.~\ref{alg:ilp}}
\eIf{ILP terminated optimally}{
  $u \leftarrow \radius_{\mathrm{ILP}}$\;
}{
  $u \leftarrow \min(u, \mathrm{GreedyBeamSearch}(\DAG, \mathcal{F}))$
    \tcp*{Heuristic upper bound}
}
\tcp{Step 8: Report generation}
\Return{$(\mathcal{F}, [\ell, u], \text{runtime breakdown}, \text{power diagnostics})$}\;
\end{algorithm}

\noindent\textbf{Complexity:} Dominated by ILP (NP-hard worst case). Fragility scoring
is polynomial. The pipeline provides a radius \emph{bracket} $[\ell, u]$ even when the
ILP times out, using the LP lower bound and greedy upper bound.

\subsection{CDCL-DAGRepair: SMT-Based Minimum-Edit Search}\label{sec:cdcl}

The ILP formulation (\Cref{alg:ilp}) delegates reasoning to a general-purpose solver.
We now present a domain-specific alternative: a DPLL(T)-style solver whose theory
component enforces DAG acyclicity and CI consistency natively. This avoids the
combinatorial blowup of linearizing acyclicity constraints into big-$M$ integer
programs, and permits conflict-driven learning tailored to the DAG edit domain.

\begin{algorithm}[H]
\caption{CDCL-DAGRepair (SMT-Based Minimum-Edit DAG Repair)}\label{alg:cdcl}
\KwIn{DAG $\DAG = (V, E)$, candidate edits $\mathcal{E}$, CI constraints $\mathcal{C}$,
  predicate $\Phi$, treatment $X$, outcome $Y$, edit budget $k$}
\KwOut{Minimum edit set $S^* \subseteq \mathcal{E}$ with $|S^*| \leq k$ s.t.\
  $\neg\Phi(\DAG \oplus S^*, X, Y)$, or UNSAT}
\BlankLine
\tcp{Initialization}
$\alpha \leftarrow \emptyset$ \tcp*{Current partial assignment: $e_i \mapsto \{0,1,\bot\}$}
$\mathrm{trail} \leftarrow []$ \tcp*{Decision/propagation trail with levels}
$\mathrm{level} \leftarrow 0$ \tcp*{Current decision level}
$\mathrm{clauses} \leftarrow \emptyset$ \tcp*{Learned conflict clauses}
$\mathrm{activity}[e_i] \leftarrow \fragility(e_i)$ for all $e_i \in \mathcal{E}$
  \tcp*{VSIDS scores initialized from fragility}
$\mathrm{decay} \leftarrow 0.95$ \tcp*{VSIDS decay factor}
$\DAG_{\mathrm{curr}} \leftarrow \DAG$ \tcp*{Maintained incrementally}
$\mathrm{cache} \leftarrow$ initialize d-separation cache\;
\BlankLine
\tcp{Cardinality constraint: at most $k$ edits active}
Add clause: $\sum_{e_i \in \mathcal{E}} e_i \leq k$ via sequential counter encoding\;
\BlankLine
\tcp{Main CDCL search loop}
\While{\textbf{true}}{
  \tcp{Unit propagation on Boolean clauses}
  $\mathrm{conflict} \leftarrow \mathrm{BooleanPropagate}(\alpha, \mathrm{clauses}, \mathrm{trail})$\;
  \If{$\mathrm{conflict} \neq \texttt{nil}$}{
    \tcp{Conflict from Boolean clauses}
    \eIf{$\mathrm{level} = 0$}{
      \Return{UNSAT} \tcp*{Root-level conflict: no solution within budget}
    }{
      $(C_{\mathrm{learn}}, \mathrm{btlevel}) \leftarrow \mathrm{AnalyzeConflict}(\mathrm{conflict}, \mathrm{trail})$\;
      $\mathrm{clauses} \leftarrow \mathrm{clauses} \cup \{C_{\mathrm{learn}}\}$\;
      $\mathrm{Backjump}(\alpha, \mathrm{trail}, \mathrm{btlevel})$\;
      $\mathrm{level} \leftarrow \mathrm{btlevel}$\;
      $\mathrm{BumpActivity}(C_{\mathrm{learn}}, \mathrm{activity}, \mathrm{decay})$\;
      \textbf{continue}\;
    }
  }
  \BlankLine
  \tcp{Theory check: acyclicity of current partial assignment}
  $\mathrm{cycle} \leftarrow \mathrm{TheoryCheckAcyclicity}(\DAG_{\mathrm{curr}}, \alpha)$\;
  \If{$\mathrm{cycle} \neq \texttt{nil}$}{
    \tcp{Extract minimal acyclicity conflict clause}
    $C_{\mathrm{acyc}} \leftarrow \mathrm{ExtractCycleConflict}(\mathrm{cycle}, \alpha, \mathrm{trail})$\;
    $\mathrm{clauses} \leftarrow \mathrm{clauses} \cup \{C_{\mathrm{acyc}}\}$\;
    $(C_{\mathrm{learn}}, \mathrm{btlevel}) \leftarrow \mathrm{AnalyzeConflict}(C_{\mathrm{acyc}}, \mathrm{trail})$\;
    $\mathrm{Backjump}(\alpha, \mathrm{trail}, \mathrm{btlevel})$\;
    $\mathrm{level} \leftarrow \mathrm{btlevel}$\;
    $\mathrm{BumpActivity}(C_{\mathrm{acyc}}, \mathrm{activity}, \mathrm{decay})$\;
    \textbf{continue}\;
  }
  \BlankLine
  \tcp{Check if all variables assigned}
  \eIf{$\alpha$ is a complete assignment}{
    \tcp{Lazy CI evaluation: check CI consistency only at full assignment}
    $\DAG' \leftarrow \DAG \oplus \{e_i : \alpha(e_i) = 1\}$\;
    $\mathrm{violated} \leftarrow \texttt{nil}$\;
    \ForEach{$(A \CI B \mid S) \in \mathcal{C}$}{
      $\mathrm{sep} \leftarrow \mathrm{IncrDSep}(\DAG, \{e_i : \alpha(e_i) = 1\}, A, B, S, \mathrm{cache})$\;
      \If{$\neg\mathrm{sep}$ \textbf{and} $(A \CI B \mid S)$ was statistically confirmed}{
        $\mathrm{violated} \leftarrow (A, B, S)$\;
        \textbf{break}\;
      }
    }
    \eIf{$\mathrm{violated} \neq \texttt{nil}$}{
      \tcp{CI conflict: extract minimal set of edits causing violation}
      $C_{\mathrm{ci}} \leftarrow \mathrm{ExtractCIConflict}(\mathrm{violated}, \alpha, \DAG, \mathrm{cache})$\;
      $\mathrm{clauses} \leftarrow \mathrm{clauses} \cup \{C_{\mathrm{ci}}\}$\;
      $(C_{\mathrm{learn}}, \mathrm{btlevel}) \leftarrow \mathrm{AnalyzeConflict}(C_{\mathrm{ci}}, \mathrm{trail})$\;
      $\mathrm{Backjump}(\alpha, \mathrm{trail}, \mathrm{btlevel})$\;
      $\mathrm{level} \leftarrow \mathrm{btlevel}$\;
      $\mathrm{BumpActivity}(C_{\mathrm{ci}}, \mathrm{activity}, \mathrm{decay})$\;
    }{
      \tcp{Full assignment passes acyclicity + CI: check conclusion negation}
      \If{$\neg\Phi(\DAG', X, Y)$}{
        \Return{$\{e_i : \alpha(e_i) = 1\}$} \tcp*{Found valid repair}
      }
      \tcp{Conclusion still holds: block this assignment}
      $C_{\mathrm{block}} \leftarrow \bigvee_{e_i : \alpha(e_i)=1} \neg e_i
        \lor \bigvee_{e_j : \alpha(e_j)=0} e_j$\;
      $\mathrm{clauses} \leftarrow \mathrm{clauses} \cup \{C_{\mathrm{block}}\}$\;
      $(C_{\mathrm{learn}}, \mathrm{btlevel}) \leftarrow \mathrm{AnalyzeConflict}(C_{\mathrm{block}}, \mathrm{trail})$\;
      $\mathrm{Backjump}(\alpha, \mathrm{trail}, \mathrm{btlevel})$\;
      $\mathrm{level} \leftarrow \mathrm{btlevel}$\;
    }
  }{
    \tcp{Decision: pick unassigned variable with highest VSIDS activity}
    $e^* \leftarrow \argmax_{e_i : \alpha(e_i) = \bot} \mathrm{activity}[e_i]$\;
    $\mathrm{level} \leftarrow \mathrm{level} + 1$\;
    $\alpha(e^*) \leftarrow 1$ \tcp*{Branch on activating the edit}
    $\mathrm{trail}.\mathrm{push}((e^*, 1, \mathrm{level}, \textsc{Decision}))$\;
    $\DAG_{\mathrm{curr}} \leftarrow \DAG_{\mathrm{curr}} \oplus e^*$ \tcp*{Incremental DAG update}
  }
}
\BlankLine
\tcp{--- Subroutines ---}
\SetKwProg{Fn}{Function}{:}{}
\Fn{$\mathrm{AnalyzeConflict}(\mathrm{conflict}, \mathrm{trail})$}{
  \tcp{First-UIP conflict analysis}
  $C \leftarrow \mathrm{conflict}$\;
  \While{$|\{l \in C : \mathrm{trail}.\mathrm{level}(l) = \mathrm{level}\}| > 1$}{
    $l \leftarrow$ last-assigned literal in $C$ at current level\;
    $\mathrm{reason} \leftarrow$ clause that propagated $l$\;
    $C \leftarrow \mathrm{Resolve}(C, \mathrm{reason}, l)$ \tcp*{Resolution step}
  }
  $\mathrm{btlevel} \leftarrow$ second-highest decision level in $C$\;
  \Return{$(C, \mathrm{btlevel})$}\;
}
\BlankLine
\Fn{$\mathrm{ExtractCycleConflict}(\mathrm{cycle}, \alpha, \mathrm{trail})$}{
  \tcp{Identify minimal set of active edits forming the cycle}
  $E_{\mathrm{cyc}} \leftarrow \{e_i : \alpha(e_i) = 1 \text{ and } e_i \text{ contributes an edge to } \mathrm{cycle}\}$\;
  \tcp{At least one must be deactivated to break the cycle}
  \Return{$\bigvee_{e_i \in E_{\mathrm{cyc}}} \neg e_i$}\;
}
\BlankLine
\Fn{$\mathrm{ExtractCIConflict}(\mathrm{violated}, \alpha, \DAG, \mathrm{cache})$}{
  \tcp{Find minimal subset of active edits that causes CI violation}
  $(A, B, S) \leftarrow \mathrm{violated}$\;
  $E_{\mathrm{active}} \leftarrow \{e_i : \alpha(e_i) = 1\}$, sorted by trail order\;
  $E_{\mathrm{min}} \leftarrow E_{\mathrm{active}}$\;
  \ForEach{$e_i \in E_{\mathrm{active}}$}{
    \If{$\mathrm{IncrDSep}(\DAG, E_{\mathrm{min}} \setminus \{e_i\}, A, B, S, \mathrm{cache})$ still violates}{
      $E_{\mathrm{min}} \leftarrow E_{\mathrm{min}} \setminus \{e_i\}$ \tcp*{$e_i$ not needed for conflict}
    }
  }
  \Return{$\bigvee_{e_i \in E_{\mathrm{min}}} \neg e_i$}\;
}
\BlankLine
\Fn{$\mathrm{TheoryCheckAcyclicity}(\DAG_{\mathrm{curr}}, \alpha)$}{
  \tcp{Incremental cycle detection using topological order maintenance}
  Maintain dynamic topological order $\pi$ of $\DAG_{\mathrm{curr}}$\;
  \ForEach{newly assigned $e_i = (u, v, \mathrm{type})$ with $\alpha(e_i) = 1$}{
    Let $(s, t)$ be the directed edge added by edit $e_i$\;
    \If{$\pi(t) < \pi(s)$}{
      Attempt to reorder vertices between $t$ and $s$ in $\pi$\;
      \If{reordering fails (back-edge found)}{
        \Return{cycle path from $t$ to $s$ in $\DAG_{\mathrm{curr}}$}\;
      }
    }
  }
  \Return{\texttt{nil}} \tcp*{No cycle detected}
}
\end{algorithm}

\noindent\textbf{Complexity:} $O(2^k \cdot m \cdot T_{\mathrm{CI}})$ where $k$ is the
edit budget, $m = |\mathcal{E}|$ is the number of candidate edits, and
$T_{\mathrm{CI}}$ is the time per CI test query. The $2^k$ factor arises because the
search tree has depth at most $k$ (enforced by the cardinality constraint) and branching
factor~2 per variable; conflict-driven learning prunes large portions of this tree in
practice. Each node requires $O(m)$ work for propagation and $O(T_{\mathrm{CI}})$ for
theory checks when a full assignment is reached.

\smallskip
\noindent\textbf{Correctness:} \emph{Soundness} follows from the soundness of the DAG
theory solver: every returned edit set $S^*$ satisfies (i)~$\DAG \oplus S^*$ is acyclic,
(ii)~all statistically confirmed CI constraints are preserved in $\DAG \oplus S^*$, and
(iii)~$\neg\Phi(\DAG \oplus S^*, X, Y)$ holds. \emph{Completeness} within the $k$-edit
neighborhood follows from exhaustive search over the Boolean space $\{0,1\}^m$ restricted
to $\sum e_i \leq k$: every branch is either explored or pruned by a valid conflict
clause. Learned clauses are logical consequences of the original constraints and the
theory axioms, so no valid solution is excluded.

\smallskip
\noindent\textbf{Design rationale (lazy CI evaluation):} Following the red-team
recommendation, CI constraints are checked lazily---only when a complete assignment is
reached---rather than eagerly at every partial assignment. This avoids
$O(|\mathcal{C}|)$ expensive d-separation queries per propagation step. The lazy strategy
is sound because CI violations at partial assignments would also manifest at any
completion of that assignment; by delaying evaluation, we allow the Boolean propagation
and acyclicity checks (both cheap) to prune infeasible branches first.

% ─────────────────────────────────────────────────────────────────

\subsection{LP-Rounding for Radius Approximation}\label{sec:lp_rounding}

When the ILP (\Cref{alg:ilp}) times out on large instances, we need a polynomial-time
algorithm that provides an approximate robustness radius with a provable guarantee. The
LP relaxation (\Cref{alg:lp}) gives a lower bound; the following rounding algorithm
converts the fractional LP solution into a feasible integral solution whose cost is at
most $O(w)$ times optimal, where $w$ is the treewidth of the moral graph.

\begin{algorithm}[H]
\caption{LP-Rounding for Radius Approximation}\label{alg:lp_round}
\KwIn{DAG $\DAG$, CI constraints $\mathcal{C}$, predicate $\Phi$, treatment $X$,
  outcome $Y$, treewidth bound $w$}
\KwOut{Feasible edit set $S$ with $|S| \leq (2w+2) \cdot \radius^*$, where $\radius^*$
  is the optimal robustness radius}
\BlankLine
\tcp{Step 1: Solve LP relaxation}
Construct LP relaxation of \Cref{alg:ilp}: replace $x_e \in \{0,1\}$ with $x_e \in [0,1]$
  for all edit variables $e \in \mathcal{E}$\;
$(x^{\mathrm{LP}}, \mathrm{OPT}_{\mathrm{LP}}) \leftarrow \mathrm{SolveLP}()$\;
\BlankLine
\tcp{Step 2: Threshold rounding}
$\tau \leftarrow 1 / (2w + 1)$ \tcp*{Rounding threshold from integrality gap}
$S_{\mathrm{round}} \leftarrow \emptyset$\;
\ForEach{edit variable $e_i \in \mathcal{E}$}{
  \eIf{$x_{e_i}^{\mathrm{LP}} \geq \tau$}{
    $S_{\mathrm{round}} \leftarrow S_{\mathrm{round}} \cup \{e_i\}$ \tcp*{Round up}
  }{
    \tcp{Round down: do not include this edit}
  }
}
\BlankLine
\tcp{Step 3: Acyclicity repair}
$\DAG_{\mathrm{round}} \leftarrow \DAG \oplus S_{\mathrm{round}}$\;
\While{$\DAG_{\mathrm{round}}$ contains a directed cycle}{
  $C \leftarrow$ find a directed cycle in $\DAG_{\mathrm{round}}$ via DFS\;
  \tcp{Identify the minimum-weight cycle edge (prefer removing recently added edits)}
  $e^* \leftarrow \argmin_{e \in C \cap S_{\mathrm{round}}} x_{e}^{\mathrm{LP}}$\;
  \eIf{$e^* \neq \texttt{nil}$}{
    $S_{\mathrm{round}} \leftarrow S_{\mathrm{round}} \setminus \{e^*\}$ \tcp*{Remove rounded edit}
  }{
    \tcp{Cycle involves only original edges: add a deletion edit}
    $e_{\mathrm{del}} \leftarrow$ edge in $C$ with highest $\fragility$ score\;
    $S_{\mathrm{round}} \leftarrow S_{\mathrm{round}} \cup \{(e_{\mathrm{del}}, \mathrm{del})\}$\;
  }
  $\DAG_{\mathrm{round}} \leftarrow \DAG \oplus S_{\mathrm{round}}$\;
}
\BlankLine
\tcp{Step 4: CI consistency check}
\ForEach{$(A \CI B \mid S) \in \mathcal{C}$}{
  \If{$\neg\dsep_{\DAG_{\mathrm{round}}}(A, B \mid S)$ \textbf{and} $(A \CI B \mid S)$ confirmed}{
    \tcp{CI violation: remove the edit most responsible}
    $e_{\mathrm{resp}} \leftarrow \argmax_{e \in S_{\mathrm{round}}} \Delta_{\dsep}(e, A, B, S)$
      \tcp*{Edit with largest d-sep impact}
    $S_{\mathrm{round}} \leftarrow S_{\mathrm{round}} \setminus \{e_{\mathrm{resp}}\}$\;
    $\DAG_{\mathrm{round}} \leftarrow \DAG \oplus S_{\mathrm{round}}$\;
  }
}
\BlankLine
\tcp{Step 5: Augmentation to achieve $\neg\Phi$}
\If{$\Phi(\DAG_{\mathrm{round}}, X, Y)$ still holds}{
  \tcp{Greedy augmentation: add highest-fragility edits until $\neg\Phi$ achieved}
  $\mathcal{E}_{\mathrm{rem}} \leftarrow \mathcal{E} \setminus S_{\mathrm{round}}$, sorted by
    $\fragility$ descending\;
  \ForEach{$e_i \in \mathcal{E}_{\mathrm{rem}}$}{
    \If{$\DAG_{\mathrm{round}} \oplus e_i$ is acyclic}{
      $S_{\mathrm{round}} \leftarrow S_{\mathrm{round}} \cup \{e_i\}$\;
      $\DAG_{\mathrm{round}} \leftarrow \DAG_{\mathrm{round}} \oplus e_i$\;
      \If{$\neg\Phi(\DAG_{\mathrm{round}}, X, Y)$}{
        \textbf{break} \tcp*{Conclusion negated; done}
      }
    }
  }
}
\BlankLine
\tcp{Step 6: Derandomization via conditional expectations}
\tcp{(Applied when randomized rounding is used as alternative to threshold rounding)}
$S_{\mathrm{derand}} \leftarrow \emptyset$\;
\ForEach{$e_i \in \mathcal{E}$ in decreasing order of $x_{e_i}^{\mathrm{LP}}$}{
  $\E_1 \leftarrow \E[|S| \mid e_i = 1, S_{\mathrm{derand}}]$
    \tcp*{Expected cost if $e_i$ included}
  $\E_0 \leftarrow \E[|S| \mid e_i = 0, S_{\mathrm{derand}}]$
    \tcp*{Expected cost if $e_i$ excluded}
  \eIf{$\E_1 \leq \E_0$ \textbf{and} including $e_i$ is feasible}{
    $S_{\mathrm{derand}} \leftarrow S_{\mathrm{derand}} \cup \{e_i\}$\;
  }{
    \tcp{Exclude $e_i$}
  }
}
\If{$|S_{\mathrm{derand}}| < |S_{\mathrm{round}}|$ \textbf{and} $\neg\Phi(\DAG \oplus S_{\mathrm{derand}}, X, Y)$}{
  $S_{\mathrm{round}} \leftarrow S_{\mathrm{derand}}$ \tcp*{Use derandomized solution if better}
}
\BlankLine
\Return{$S_{\mathrm{round}}$}\;
\end{algorithm}

\noindent\textbf{Complexity:} Polynomial overall. The LP solve runs in
$O(|\mathcal{E}|^{3.5})$ via interior-point methods (or faster with specialized
structure). Threshold rounding is $O(|\mathcal{E}|)$. Acyclicity repair requires
$O(|\An|^2)$ time for repeated cycle detection via DFS. Augmentation requires
$O(|\An| \cdot \log |\An|)$ for sorted iteration with acyclicity checks.
Derandomization adds a single $O(|\mathcal{E}|)$ pass. Total:
$O(|\mathcal{E}|^{3.5} + |\An|^2 + |\An| \log |\An|)$.

\smallskip
\noindent\textbf{Correctness and approximation ratio:}
The threshold rounding at $\tau = 1/(2w+1)$ guarantees that for each rounded variable
$e_i \in S_{\mathrm{round}}$, we have $x_{e_i}^{\mathrm{LP}} \geq \tau$, so
$|S_{\mathrm{round}}| \leq (2w+1) \cdot \sum_i x_{e_i}^{\mathrm{LP}} = (2w+1) \cdot
\mathrm{OPT}_{\mathrm{LP}} \leq (2w+1) \cdot \radius^*$, where the last inequality
follows from LP relaxation. The acyclicity repair removes at most as many edits as it
adds (net change $\leq 0$). The augmentation step adds at most $\radius^*$ additional
edits (since $\radius^*$ edits suffice to negate $\Phi$ from any DAG). The total
guarantee is $|S| \leq (2w+1) \cdot \radius^* + \radius^* = (2w+2) \cdot \radius^*$.
The derandomization via conditional expectations ensures the deterministic solution is
no worse than the expected randomized cost.

% ─────────────────────────────────────────────────────────────────

\subsection{Power-Adaptive CI Ensemble}\label{sec:adaptive_ci}

The CI ensemble (\Cref{alg:ci_ensemble}) runs all three tests unconditionally. For large
DAGs, where thousands of CI tests are required, the cubic cost of KCI dominates runtime.
The following power-adaptive variant runs cheap tests first and escalates only when the
result is ambiguous, achieving amortized linear cost per test on benign (Gaussian) data.

\begin{algorithm}[H]
\caption{Power-Adaptive CI Ensemble}\label{alg:adaptive_ci}
\KwIn{Data $D \in \R^{n \times p}$, variables $A$, $B$, conditioning set $S$,
  significance level $\alpha$, equivalence margin $\varepsilon_0$}
\KwOut{Decision (independent/dependent/inconclusive), combined p-value,
  power diagnostic $\hat{\beta}$}
\BlankLine
\tcp{Tier 1: Fisher-z partial correlation test --- $O(n)$ per test}
$r_{AB \cdot S} \leftarrow \mathrm{PartialCorrelation}(D, A, B, S)$
  \tcp*{$O(n \cdot |S|)$ via regression}
$z \leftarrow \frac{1}{2} \ln\!\left(\frac{1 + r_{AB \cdot S}}{1 - r_{AB \cdot S}}\right)
  \cdot \sqrt{n - |S| - 3}$ \tcp*{Fisher z-transform}
$p_1 \leftarrow 2 \cdot (1 - \Phi_{\mathrm{std}}(|z|))$
  \tcp*{Two-sided p-value under $H_0: \rho = 0$}
\BlankLine
\tcp{TOST equivalence test for independence conclusion}
$z_{\mathrm{upper}} \leftarrow (r_{AB \cdot S} - \varepsilon_0) \cdot \sqrt{n - |S| - 3}$\;
$z_{\mathrm{lower}} \leftarrow (r_{AB \cdot S} + \varepsilon_0) \cdot \sqrt{n - |S| - 3}$\;
$p_{\mathrm{TOST}} \leftarrow \max\!\bigl(\Phi_{\mathrm{std}}(z_{\mathrm{upper}}),\;
  1 - \Phi_{\mathrm{std}}(z_{\mathrm{lower}})\bigr)$\;
\BlankLine
\tcp{Decision at Tier 1}
\uIf{$p_1 < \alpha / 10$}{
  \tcp{Strong rejection of $H_0$: clearly dependent, no escalation needed}
  $\mathrm{decision}_1 \leftarrow \mathrm{dependent}$\;
  $\mathrm{tiers\_run} \leftarrow \{1\}$\;
  $\mathrm{escalate} \leftarrow \texttt{false}$\;
}
\uElseIf{$p_{\mathrm{TOST}} < \alpha$ \textbf{and} $p_1 > 10\alpha$}{
  \tcp{Strong evidence for independence via equivalence test}
  $\mathrm{decision}_1 \leftarrow \mathrm{independent}$\;
  $\mathrm{tiers\_run} \leftarrow \{1\}$\;
  $\mathrm{escalate} \leftarrow \texttt{false}$\;
}
\Else{
  \tcp{Ambiguous region: $p_1 \in [\alpha/10, 10\alpha]$ or TOST inconclusive}
  $\mathrm{escalate} \leftarrow \texttt{true}$\;
}
\BlankLine
\tcp{Tier 2: Rank-based partial correlation --- $O(n \log n)$ per test}
\If{$\mathrm{escalate}$}{
  $D_{\mathrm{rank}} \leftarrow \mathrm{RankTransform}(D, \{A, B\} \cup S)$
    \tcp*{Replace values with ranks}
  $r_{\mathrm{rank}} \leftarrow \mathrm{PartialCorrelation}(D_{\mathrm{rank}}, A, B, S)$\;
  $z_{\mathrm{rank}} \leftarrow \frac{1}{2}
    \ln\!\left(\frac{1 + r_{\mathrm{rank}}}{1 - r_{\mathrm{rank}}}\right)
    \cdot \sqrt{n - |S| - 3}$\;
  $p_2 \leftarrow 2 \cdot (1 - \Phi_{\mathrm{std}}(|z_{\mathrm{rank}}|))$\;
  $\mathrm{tiers\_run} \leftarrow \{1, 2\}$\;
  \BlankLine
  \tcp{Check if Tier 2 resolves the ambiguity}
  \uIf{$p_2 < \alpha / 10$ \textbf{or} $p_2 > 10\alpha$}{
    $\mathrm{escalate} \leftarrow \texttt{false}$ \tcp*{Tier 2 is decisive}
  }{
    $\mathrm{escalate} \leftarrow \texttt{true}$ \tcp*{Still ambiguous: escalate to Tier 3}
  }
}
\BlankLine
\tcp{Tier 3: KCI with Nystr\"om approximation --- $O(nr^2 + r^3)$}
\If{$\mathrm{escalate}$}{
  \tcp{Kernel CI test with Gaussian RBF kernel and bandwidth selection}
  $\sigma_A \leftarrow \mathrm{MedianHeuristic}(D_A)$\;
  $\sigma_B \leftarrow \mathrm{MedianHeuristic}(D_B)$\;
  $K_A \leftarrow \mathrm{Nystrom}(D_A, r, \sigma_A)$
    \tcp*{$n \times r$ low-rank factor}
  $K_B \leftarrow \mathrm{Nystrom}(D_B, r, \sigma_B)$\;
  \If{$|S| > 0$}{
    $\sigma_S \leftarrow \mathrm{MedianHeuristic}(D_S)$\;
    $K_S \leftarrow \mathrm{Nystrom}(D_S, r, \sigma_S)$\;
    \tcp{Residualize kernels on conditioning set}
    $\tilde{K}_A \leftarrow K_A - K_S (K_S^\top K_S + \lambda I)^{-1} K_S^\top K_A$\;
    $\tilde{K}_B \leftarrow K_B - K_S (K_S^\top K_S + \lambda I)^{-1} K_S^\top K_B$\;
  }
  \Else{
    $\tilde{K}_A \leftarrow K_A$; $\tilde{K}_B \leftarrow K_B$\;
  }
  \tcp{HSIC test statistic via Nystr\"om factors}
  $\mathrm{HSIC}_n \leftarrow \frac{1}{n^2} \|\tilde{K}_A^\top \tilde{K}_B\|_F^2$\;
  $p_3 \leftarrow \mathrm{GammaApprox}(\mathrm{HSIC}_n, n, r)$
    \tcp*{Gamma approximation to null}
  $\mathrm{tiers\_run} \leftarrow \{1, 2, 3\}$\;
}
\BlankLine
\tcp{Cauchy combination of all tiers that were actually run}
$T_{\mathrm{Cauchy}} \leftarrow \frac{1}{|\mathrm{tiers\_run}|}
  \sum_{j \in \mathrm{tiers\_run}} \tan\!\bigl(\pi(0.5 - p_j)\bigr)$\;
$p_{\mathrm{combined}} \leftarrow \frac{1}{2} - \frac{\arctan(T_{\mathrm{Cauchy}})}{\pi}$\;
\BlankLine
\tcp{Power diagnostic via analytical formula (avoid expensive permutations)}
$\hat{\beta} \leftarrow 1 - \Phi_{\mathrm{std}}\!\bigl(\Phi_{\mathrm{std}}^{-1}(1 - \alpha)
  - |r_{AB \cdot S}| \cdot \sqrt{n - |S| - 3}\bigr)$\;
\BlankLine
\tcp{Final decision}
\If{$\hat{\beta} < 0.8$}{
  \Return{$(\mathrm{inconclusive}, p_{\mathrm{combined}}, \hat{\beta})$}\;
}
\eIf{$p_{\mathrm{combined}} < \alpha$}{
  \Return{$(\mathrm{dependent}, p_{\mathrm{combined}}, \hat{\beta})$}\;
}{
  \Return{$(\mathrm{independent}, p_{\mathrm{combined}}, \hat{\beta})$}\;
}
\end{algorithm}

\noindent\textbf{Amortized complexity:} When Tier~1 resolves the test (expected
${\sim}70\%$ of tests for Gaussian data, based on the clear separation between
$p < \alpha/10$ and $p > 10\alpha$ under standard effect sizes), the per-test cost is
$O(n \cdot |S|)$. The ${\sim}25\%$ escalating to Tier~2 cost $O(n \log n)$ for rank
transformation. Only the ${\sim}5\%$ requiring Tier~3 incur $O(nr^2 + r^3)$ for
Nystr\"om KCI. The amortized cost across all tests is therefore
$0.70 \cdot O(n) + 0.25 \cdot O(n \log n) + 0.05 \cdot O(nr^2) = O(n)$ per test
when $r = O(1)$, or $O(n \cdot r^2 / 20)$ when $r = \omega(1)$.

\smallskip
\noindent\textbf{Correctness:} Type~I error control follows from the Cauchy combination
test \cite{liu2020}: under the global null $H_0: A \CI B \mid S$, each p-value $p_j$ is
marginally valid (super-uniform), and the Cauchy combination statistic has a standard
Cauchy distribution regardless of the dependence structure among the $p_j$ values. The
combined p-value is therefore valid at level~$\alpha$. The TOST equivalence test provides
a principled mechanism for concluding independence (not merely failing to reject
dependence): if $|r_{AB \cdot S}| < \varepsilon_0$ with confidence $1 - \alpha$, we
declare independence.

% ─────────────────────────────────────────────────────────────────

\subsection{Parallel Fragility Scoring}\label{sec:parallel_fragility}

Computing fragility scores (\Cref{alg:fragility}) for all candidate edits is
embarrassingly parallel when edits do not interfere with each other's d-separation
computations. We formalize this via a graph-coloring argument: edits whose endpoints
are non-adjacent in the moral graph have independent d-separation effects and can be
scored simultaneously.

\begin{algorithm}[H]
\caption{Parallel Fragility Scoring}\label{alg:parallel_fragility}
\KwIn{DAG $\DAG$, treatment $X$, outcome $Y$, data $D$, predicate $\Phi$,
  number of threads $T_{\mathrm{threads}}$}
\KwOut{Fragility scores $\{\fragility(e)\}$ for all candidate edits $e \in \mathcal{E}$}
\BlankLine
\tcp{Step 1: Enumerate candidate edits}
$\An \leftarrow \mathrm{AncestralSet}(\DAG, X, Y)$ \tcp*{Alg.~\ref{alg:ancestral}}
$\mathcal{E} \leftarrow \emptyset$\;
\ForEach{$(i, j) \in \An \times \An$, $i \neq j$}{
  \If{$(i,j) \notin E$ \textbf{and} $\DAG \cup \{(i,j)\}$ is acyclic}{
    $\mathcal{E} \leftarrow \mathcal{E} \cup \{(i, j, \mathrm{add})\}$\;
  }
  \If{$(i,j) \in E$}{
    $\mathcal{E} \leftarrow \mathcal{E} \cup \{(i, j, \mathrm{del})\}$\;
    \If{$\DAG \setminus \{(i,j)\} \cup \{(j,i)\}$ is acyclic}{
      $\mathcal{E} \leftarrow \mathcal{E} \cup \{(i, j, \mathrm{rev})\}$\;
    }
  }
}
\BlankLine
\tcp{Step 2: Build moral graph and conflict graph}
$M \leftarrow \mathrm{MoralGraph}(\DAG)$
  \tcp*{Undirected: drop orientations, marry co-parents}
\tcp{Conflict graph $G_C$: edits conflict if they share an endpoint in $M$}
$V_C \leftarrow \mathcal{E}$\;
$E_C \leftarrow \emptyset$\;
\ForEach{pair $(e_1, e_2) \in \mathcal{E} \times \mathcal{E}$, $e_1 \neq e_2$}{
  Let $e_1 = (i_1, j_1, \_)$ and $e_2 = (i_2, j_2, \_)$\;
  \If{$\{i_1, j_1\} \cap \{i_2, j_2\} \neq \emptyset$
    \textbf{or} $\exists$ edge in $M$ between $\{i_1, j_1\}$ and $\{i_2, j_2\}$}{
    $E_C \leftarrow E_C \cup \{(e_1, e_2)\}$\;
  }
}
$G_C \leftarrow (V_C, E_C)$\;
\BlankLine
\tcp{Step 3: Greedy graph coloring of $G_C$}
$\mathrm{color}[e] \leftarrow \bot$ for all $e \in \mathcal{E}$\;
$\chi \leftarrow 0$ \tcp*{Chromatic number upper bound}
\ForEach{$e \in \mathcal{E}$ in order of decreasing degree in $G_C$}{
  $\mathrm{used} \leftarrow \{\mathrm{color}[e'] : (e, e') \in E_C,\;
    \mathrm{color}[e'] \neq \bot\}$\;
  $\mathrm{color}[e] \leftarrow \min\{c \in \{1, 2, \ldots\} : c \notin \mathrm{used}\}$\;
  $\chi \leftarrow \max(\chi, \mathrm{color}[e])$\;
}
\BlankLine
\tcp{Step 4: Score edits color-class by color-class}
$\mathrm{ADJ} \leftarrow$ enumerate valid adjustment sets for $(\DAG, X, Y)$\;
\For{$c \leftarrow 1$ \KwTo $\chi$}{
  $\mathcal{E}_c \leftarrow \{e \in \mathcal{E} : \mathrm{color}[e] = c\}$
    \tcp*{Independent set in $G_C$}
  \tcp{Score all edits in $\mathcal{E}_c$ in parallel}
  \textbf{parallel for} each $e \in \mathcal{E}_c$ using up to $T_{\mathrm{threads}}$ threads\;
  \Begin{
    \tcp{d-Separation channel}
    $\mathcal{D}_e \leftarrow$ thread-local d-separation cache\;
    $F_{\dsep}(e) \leftarrow \max_{S \in \mathrm{ADJ}} \mathbf{1}[\dsep_\DAG(X,Y|S)
      \neq \mathrm{IncrDSep}(\DAG, e, X, Y, S, \mathcal{D}_e)]$\;
    \BlankLine
    \tcp{Identification channel}
    $F_{\mathrm{id}}(e) \leftarrow \mathbf{1}[\Phi(\DAG) \neq \Phi(\DAG \oplus e)]$\;
    \BlankLine
    \tcp{Estimation channel (thread-safe read-only data access)}
    $F_{\mathrm{est}}(e) \leftarrow |\hat\tau(\DAG) - \hat\tau(\DAG \oplus e)| /
      \max(|\hat\tau(\DAG)|, \varepsilon_0)$\;
    \BlankLine
    \tcp{Combined fragility}
    $\fragility(e) \leftarrow \max(F_{\dsep}(e), F_{\mathrm{id}}(e), F_{\mathrm{est}}(e))$\;
  }
}
\BlankLine
\Return{$\{(e, \fragility(e)) : e \in \mathcal{E}\}$ sorted by $\fragility$ descending}\;
\end{algorithm}

\noindent\textbf{Correctness:} Edits in the same color class form an independent set in
$G_C$, meaning no two edits share an endpoint or are adjacent in the moral graph $M$.
By the locality of the Bayes-Ball algorithm \cite{shachter1998}, d-separation between
nodes $A$ and $B$ given $S$ is determined entirely by the subgraph induced by the
ancestors of $\{A, B\} \cup S$. When two edits $e_1 = (i_1, j_1)$ and $e_2 = (i_2, j_2)$
are non-adjacent in $M$, the Bayes-Ball reachability sets for $e_1$ and $e_2$ do not
interact: any active path affected by $e_1$ does not traverse $\{i_2, j_2\}$ and vice
versa. Therefore, $\mathrm{IncrDSep}(\DAG, e_1, \cdot)$ and
$\mathrm{IncrDSep}(\DAG, e_2, \cdot)$ can be evaluated independently without
interference. The identification and estimation channels compute $\Phi(\DAG \oplus e)$
and $\hat\tau(\DAG \oplus e)$ for single edits independently by definition.

\smallskip
\noindent\textbf{Parallel rounds:} The greedy coloring uses at most $\Delta(G_C) + 1$
colors, where $\Delta(G_C)$ is the maximum degree of the conflict graph. Each edit
$(i,j)$ conflicts with edits sharing endpoint $i$ (at most $\deg_M(i)$ neighbors in $M$,
each generating ${\leq}3$ edits) and edits sharing endpoint $j$ (similarly bounded).
Since $\deg_M(v) \leq 2d_{\max}$ (moralization at most doubles degree), we have
$\Delta(G_C) \leq 2 \cdot 3 \cdot 2d_{\max} = 12 d_{\max}$. More precisely,
$\chi(G_C) \leq \Delta(G_C) + 1 \leq 2d_{\max}^2 + 1$ when accounting for second-order
neighborhood effects. The number of parallel rounds is therefore
$O(d_{\max}^2)$, independent of the number of nodes $p$.

\smallskip
\noindent\textbf{Speedup:} Within each color class, all $|\mathcal{E}_c|$ edits are
scored in parallel. The speedup factor over sequential scoring is
$\min\!\bigl(\max_c |\mathcal{E}_c|,\; T_{\mathrm{threads}}\bigr)$. For a balanced
coloring, each class has ${\sim}|\mathcal{E}| / \chi$ edits, giving speedup
${\sim}\min(|\mathcal{E}| / \chi,\; T_{\mathrm{threads}})$.

\smallskip
\noindent\textbf{Note:} The estimation channel requires read-only access to the dataset
$D$ for AIPW estimation. Since each thread reads $D$ without modification, thread safety
is guaranteed without locks. Thread-local caches $\mathcal{D}_e$ avoid contention on the
d-separation cache.

% ─────────────────────────────────────────────────────────────────

\subsection{Incremental MEC Update}\label{sec:incremental_mec}

The Markov equivalence class (MEC) of a DAG $\DAG$ is represented by its completed
partially directed acyclic graph (CPDAG) $\mathcal{C}(\DAG)$. When CausalCert applies a
single edge edit, full CPDAG recomputation costs $O(p \cdot d_{\max}^2)$. The following
algorithm updates $\mathcal{C}(\DAG)$ incrementally, exploiting the fact that a single
edge edit affects only a local neighborhood.

\begin{algorithm}[H]
\caption{Incremental MEC Update}\label{alg:incremental_mec}
\KwIn{DAG $\DAG = (V, E)$, current CPDAG $\mathcal{C}(\DAG)$, single edge edit
  $e = (i, j, \mathrm{type})$}
\KwOut{Updated CPDAG $\mathcal{C}(\DAG \oplus e)$}
\BlankLine
\tcp{Apply the structural edit to the DAG}
$\DAG' \leftarrow \DAG \oplus e$\;
$\mathcal{C}' \leftarrow \mathcal{C}(\DAG)$ \tcp*{Start from current CPDAG, modify locally}
\BlankLine
\uIf{$\mathrm{type} = \mathrm{del}$}{
  \tcp{--- Edge Deletion: remove $i \to j$ ---}
  \tcp{Step D1: Remove edge $i \to j$ (or $i - j$ if undirected) from CPDAG}
  Remove the edge between $i$ and $j$ from $\mathcal{C}'$\;
  \BlankLine
  \tcp{Step D2: Check for v-structure changes at neighbors of $i$ and $j$}
  $\mathrm{affected} \leftarrow \{i, j\} \cup \mathrm{Neighbors}_{\mathcal{C}'}(i)
    \cup \mathrm{Neighbors}_{\mathcal{C}'}(j)$\;
  \ForEach{$v \in \mathrm{affected}$}{
    \ForEach{pair $(a, b) \in \Pa_{\DAG'}(v) \times \Pa_{\DAG'}(v)$, $a \neq b$}{
      \eIf{$(a, b) \notin E'$ \textbf{and} $(b, a) \notin E'$}{
        \tcp{$a \to v \leftarrow b$ is a v-structure in $\DAG'$}
        Orient $a \to v$ and $b \to v$ in $\mathcal{C}'$ (if not already oriented)\;
      }{
        \tcp{$a$ and $b$ are adjacent: not a v-structure}
        \If{$a \to v$ was oriented solely due to former v-structure involving $i \to j$}{
          Un-orient: replace $a \to v$ with $a - v$ in $\mathcal{C}'$\;
        }
      }
    }
  }
  \BlankLine
  \tcp{Step D3: Apply Meek rules locally to propagate orientation changes}
  $\mathrm{MeekLocal}(\mathcal{C}', \mathrm{affected})$\;
}
\uElseIf{$\mathrm{type} = \mathrm{add}$}{
  \tcp{--- Edge Addition: add $i \to j$ ---}
  \tcp{Step A1: Add directed edge $i \to j$ to CPDAG}
  Add $i \to j$ to $\mathcal{C}'$\;
  \BlankLine
  \tcp{Step A2: Check for new v-structures at $j$}
  $\mathrm{affected} \leftarrow \{j\} \cup \mathrm{Neighbors}_{\mathcal{C}'}(j)$\;
  \ForEach{$k \in \Pa_{\DAG'}(j)$, $k \neq i$}{
    \If{$(k, i) \notin E'$ \textbf{and} $(i, k) \notin E'$}{
      \tcp{$i \to j \leftarrow k$ is a new v-structure}
      Orient $k \to j$ in $\mathcal{C}'$ (if $k - j$ was undirected)\;
      Orient $i \to j$ in $\mathcal{C}'$\;
    }
  }
  \BlankLine
  \tcp{Step A3: Check for new v-structures at $i$}
  \ForEach{$k \in \Pa_{\DAG'}(i)$}{
    \If{$(k, j) \notin E'$ \textbf{and} $(j, k) \notin E'$}{
      \tcp{$k \to i$ and $i \to j$ but $k, j$ non-adjacent: orient $k \to i$}
      Orient $k \to i$ in $\mathcal{C}'$ (if $k - i$ was undirected)\;
    }
  }
  \BlankLine
  \tcp{Step A4: Determine if $i \to j$ should be oriented or undirected in CPDAG}
  $\mathrm{must\_orient} \leftarrow \texttt{false}$\;
  \ForEach{$k \in \Pa_{\DAG'}(j) \setminus \{i\}$}{
    \If{$(k, i) \notin E'$ \textbf{and} $(i, k) \notin E'$}{
      $\mathrm{must\_orient} \leftarrow \texttt{true}$
        \tcp*{v-structure forces orientation}
    }
  }
  \If{$\neg\mathrm{must\_orient}$}{
    \tcp{Check if any Meek rule forces orientation of $i \to j$}
    $\mathrm{must\_orient} \leftarrow \mathrm{MeekForces}(\mathcal{C}', i, j)$\;
  }
  \If{$\neg\mathrm{must\_orient}$}{
    Replace $i \to j$ with $i - j$ in $\mathcal{C}'$ \tcp*{Undirected in CPDAG}
  }
  \BlankLine
  \tcp{Step A5: Apply Meek rules locally}
  $\mathrm{MeekLocal}(\mathcal{C}', \mathrm{affected})$\;
}
\ElseIf{$\mathrm{type} = \mathrm{rev}$}{
  \tcp{--- Edge Reversal: $i \to j$ becomes $j \to i$ ---}
  \tcp{Process as deletion of $i \to j$ followed by addition of $j \to i$}
  \BlankLine
  \tcp{Step R1: Delete $i \to j$}
  Remove the edge between $i$ and $j$ from $\mathcal{C}'$\;
  $\mathrm{affected}_1 \leftarrow \{i, j\} \cup \mathrm{Neighbors}_{\mathcal{C}'}(i)
    \cup \mathrm{Neighbors}_{\mathcal{C}'}(j)$\;
  \ForEach{$v \in \mathrm{affected}_1$}{
    \ForEach{pair $(a, b) \in \Pa_{\DAG'}(v) \times \Pa_{\DAG'}(v)$, $a \neq b$}{
      \If{$(a, b) \notin E'$ \textbf{and} $(b, a) \notin E'$}{
        Orient $a \to v$ and $b \to v$ in $\mathcal{C}'$\;
      }
    }
  }
  \BlankLine
  \tcp{Step R2: Add $j \to i$}
  Add $j \to i$ to $\mathcal{C}'$\;
  $\mathrm{affected}_2 \leftarrow \{i\} \cup \mathrm{Neighbors}_{\mathcal{C}'}(i)$\;
  \ForEach{$k \in \Pa_{\DAG'}(i)$, $k \neq j$}{
    \If{$(k, j) \notin E'$ \textbf{and} $(j, k) \notin E'$}{
      Orient $k \to i$ and $j \to i$ in $\mathcal{C}'$\;
    }
  }
  \BlankLine
  \tcp{Step R3: Determine CPDAG status of $j \to i$}
  $\mathrm{must\_orient} \leftarrow \texttt{false}$\;
  \ForEach{$k \in \Pa_{\DAG'}(i) \setminus \{j\}$}{
    \If{$(k, j) \notin E'$ \textbf{and} $(j, k) \notin E'$}{
      $\mathrm{must\_orient} \leftarrow \texttt{true}$\;
    }
  }
  \If{$\neg\mathrm{must\_orient}$}{
    $\mathrm{must\_orient} \leftarrow \mathrm{MeekForces}(\mathcal{C}', j, i)$\;
  }
  \If{$\neg\mathrm{must\_orient}$}{
    Replace $j \to i$ with $j - i$ in $\mathcal{C}'$\;
  }
  \BlankLine
  \tcp{Step R4: Apply Meek rules locally}
  $\mathrm{affected} \leftarrow \mathrm{affected}_1 \cup \mathrm{affected}_2$\;
  $\mathrm{MeekLocal}(\mathcal{C}', \mathrm{affected})$\;
}
\BlankLine
\Return{$\mathcal{C}'$}\;
\BlankLine
\tcp{--- Subroutine: Local Meek Rule Application ---}
\SetKwProg{Fn}{Function}{:}{}
\Fn{$\mathrm{MeekLocal}(\mathcal{C}', \mathrm{affected})$}{
  \tcp{Apply Meek rules R1--R4 restricted to the affected zone}
  $\mathrm{queue} \leftarrow \mathrm{affected}$\;
  \While{$\mathrm{queue} \neq \emptyset$}{
    $v \leftarrow \mathrm{queue}.\mathrm{dequeue}()$\;
    \ForEach{undirected edge $v - w$ in $\mathcal{C}'$}{
      $\mathrm{oriented} \leftarrow \texttt{false}$\;
      \BlankLine
      \tcp{Meek R1: $a \to v - w$ and $a \not\!\adj w \Rightarrow v \to w$}
      \ForEach{$a$ with $a \to v$ in $\mathcal{C}'$}{
        \If{$a \not\!\adj w$ in $\mathcal{C}'$}{
          Orient $v \to w$ in $\mathcal{C}'$\;
          $\mathrm{oriented} \leftarrow \texttt{true}$; \textbf{break}\;
        }
      }
      \BlankLine
      \tcp{Meek R2: $v \to a \to w$ and $v - w \Rightarrow v \to w$}
      \If{$\neg\mathrm{oriented}$}{
        \ForEach{$a$ with $v \to a$ and $a \to w$ in $\mathcal{C}'$}{
          Orient $v \to w$ in $\mathcal{C}'$\;
          $\mathrm{oriented} \leftarrow \texttt{true}$; \textbf{break}\;
        }
      }
      \BlankLine
      \tcp{Meek R3: $v - a \to w$, $v - b \to w$, $a \not\!\adj b \Rightarrow v \to w$}
      \If{$\neg\mathrm{oriented}$}{
        \ForEach{pair $(a, b)$ with $v - a$, $a \to w$, $v - b$, $b \to w$ in $\mathcal{C}'$}{
          \If{$a \neq b$ \textbf{and} $a \not\!\adj b$ in $\mathcal{C}'$}{
            Orient $v \to w$ in $\mathcal{C}'$\;
            $\mathrm{oriented} \leftarrow \texttt{true}$; \textbf{break}\;
          }
        }
      }
      \BlankLine
      \tcp{Meek R4: $v - a \to b \to w$ and $v \adj a$ but $v \not\!\adj b \Rightarrow v \to w$}
      \If{$\neg\mathrm{oriented}$}{
        \ForEach{$a$ with $v - a$ in $\mathcal{C}'$}{
          \ForEach{$b$ with $a \to b$ and $b \to w$ in $\mathcal{C}'$}{
            \If{$v \not\!\adj b$ in $\mathcal{C}'$}{
              Orient $v \to w$ in $\mathcal{C}'$\;
              $\mathrm{oriented} \leftarrow \texttt{true}$; \textbf{break}\;
            }
          }
          \If{$\mathrm{oriented}$}{\textbf{break}}
        }
      }
      \BlankLine
      \If{$\mathrm{oriented}$}{
        $\mathrm{queue}.\mathrm{enqueue}(w)$
          \tcp*{Propagate: $w$'s neighbors may need re-evaluation}
      }
    }
  }
}
\BlankLine
\Fn{$\mathrm{MeekForces}(\mathcal{C}', u, v)$}{
  \tcp{Check if Meek rules R1--R4 force orientation of $u \to v$}
  \tcp{R1: $\exists\, a \to u$ with $a \not\!\adj v$?}
  \ForEach{$a$ with $a \to u$ in $\mathcal{C}'$}{
    \If{$a \not\!\adj v$ in $\mathcal{C}'$}{\Return{\texttt{true}}}
  }
  \tcp{R2: $\exists\, a$ with $u \to a \to v$?}
  \ForEach{$a$ with $u \to a$ and $a \to v$ in $\mathcal{C}'$}{
    \Return{\texttt{true}}\;
  }
  \tcp{R3: $\exists\, (a, b)$ with $u - a \to v$, $u - b \to v$, $a \not\!\adj b$?}
  \ForEach{pair $(a, b)$ with $u - a$, $a \to v$, $u - b$, $b \to v$ in $\mathcal{C}'$}{
    \If{$a \neq b$ \textbf{and} $a \not\!\adj b$}{\Return{\texttt{true}}}
  }
  \tcp{R4: $\exists\, (a, b)$ with $u - a \to b \to v$, $u \not\!\adj b$?}
  \ForEach{$a$ with $u - a$ in $\mathcal{C}'$}{
    \ForEach{$b$ with $a \to b$ and $b \to v$ in $\mathcal{C}'$}{
      \If{$u \not\!\adj b$}{\Return{\texttt{true}}}
    }
  }
  \Return{\texttt{false}}\;
}
\end{algorithm}

\noindent\textbf{Complexity:} The affected zone for a single edge edit is contained
within the 2-hop neighborhood of the edit endpoints in the CPDAG, which has size at
most $O(d_{\max}^2)$. Each Meek rule check at a vertex $v$ examines $O(d_{\max})$
neighbors, and each neighbor involves $O(d_{\max})$ adjacency checks. The local Meek
propagation terminates after at most $O(|\mathrm{affected}|)$ iterations (each iteration
orients at least one edge, and the number of undirected edges is finite). Total
complexity: $O(d_{\max}^3 \cdot |\mathrm{affected\_zone}|)$ where
$|\mathrm{affected\_zone}| \leq O(d_{\max}^2)$, giving $O(d_{\max}^5)$ per edit.

\smallskip
\noindent\textbf{Comparison to full recomputation:} Computing the CPDAG from scratch
(e.g., via the DAG-to-CPDAG algorithm of Chickering \cite{chickering2002}) requires
$O(p \cdot d_{\max}^2)$ time. The incremental algorithm is faster whenever the affected
zone is much smaller than the full graph, i.e., when $d_{\max}^5 \ll p \cdot d_{\max}^2$,
which simplifies to $d_{\max}^3 \ll p$. For typical causal DAGs with $d_{\max} \leq 4$
and $p \geq 10$, this condition holds: $64 \ll p$, giving a speedup of roughly
$p / d_{\max}^3$.

\smallskip
\noindent\textbf{Correctness:} The algorithm maintains the invariant that $\mathcal{C}'$
is the CPDAG of $\DAG'$ after each step. This follows from three facts:
(1)~v-structure identification in the 2-hop neighborhood is exhaustive because new
v-structures can only arise at endpoints of the edited edge or their immediate neighbors;
(2)~the Meek rules R1--R4 are complete for orientation propagation
\cite{meek1995}: iterating them to fixpoint on the skeleton with known v-structures
produces the unique CPDAG; and (3)~since only edges in the affected zone can change
orientation status, restricting Meek rule application to the affected zone (plus one
propagation layer) is sufficient.




% ══════════════════════════════════════════════════════════════════
% SECTION 6: COMPLEXITY ANALYSIS
% ══════════════════════════════════════════════════════════════════
\section{Complexity Analysis}\label{sec:complexity}

\Cref{tab:complexity} summarizes the complexity of all CausalCert algorithms.

\begin{table}[h]
\centering
\caption{Complexity summary for CausalCert algorithms. Here $p = |V|$,
$a = |\An_\DAG(X,Y)|$, $w$ = treewidth of moral graph, $d = d_{\max}$ = maximum
in-degree, $n$ = sample size, $r$ = Nystr\"om rank, $s = |S|$ = conditioning set size.}
\label{tab:complexity}
\begin{tabular}{@{}llll@{}}
\toprule
\textbf{Algorithm} & \textbf{Time} & \textbf{Space} & \textbf{Guarantee} \\
\midrule
Ancestral Set Extraction & $O(p + |E|)$ & $O(p)$ & Exact \\
Incremental d-Separation & $O(a \cdot d)$ per edit & $O(a^2)$ & Exact (\Cref{thm:incremental_dsep}) \\
Per-Edge Fragility Scoring & $O(a^3 \cdot d)$ & $O(a^2)$ & Exact per channel \\
ILP Radius (worst case) & NP-hard & $O(a^4)$ & Exact (if terminates) \\
ILP Radius (typical, $w \leq 4$) & $< 30$ min, $p \leq 20$ & $O(a^4)$ & Exact \\
LP Relaxation & Polynomial & $O(a^4)$ & Lower bound, gap $\leq O(w)$ \\
FPT DP & $O(3^{O(w^2)} \cdot a)$ & $O(3^{O(w^2)} \cdot a)$ & Exact for $w \leq 3$--$4$ \\
CI Test Ensemble & $O(n^2 s + nr^2)$ & $O(nr)$ & Valid Type~I (Cauchy) \\
\bottomrule
\end{tabular}
\end{table}

\subsection{Practical Complexity Estimates}

\Cref{tab:practical} provides concrete runtime estimates for representative DAG sizes.

\begin{table}[h]
\centering
\caption{Estimated runtime for CausalCert components on DAGs of various sizes.
Assumes $d_{\max} = 3$, $n = 1000$, $r = 100$, single-threaded execution.}
\label{tab:practical}
\begin{tabular}{@{}lcccc@{}}
\toprule
\textbf{Component} & $p = 8$ & $p = 12$ & $p = 16$ & $p = 20$ \\
 & $(a \approx 5)$ & $(a \approx 8)$ & $(a \approx 11)$ & $(a \approx 14)$ \\
\midrule
Ancestral set & $< 1$ ms & $< 1$ ms & $< 1$ ms & $< 1$ ms \\
CI testing (all pairs) & $\sim 2$ s & $\sim 5$ s & $\sim 12$ s & $\sim 25$ s \\
Fragility scoring & $\sim 0.5$ s & $\sim 3$ s & $\sim 12$ s & $\sim 35$ s \\
LP relaxation & $\sim 1$ s & $\sim 5$ s & $\sim 20$ s & $\sim 60$ s \\
ILP exact ($w \leq 3$) & $\sim 5$ s & $\sim 30$ s & $\sim 5$ min & $\sim 25$ min \\
FPT DP ($w = 3$) & $\sim 1$ s & $\sim 3$ s & $\sim 10$ s & $\sim 30$ s \\
\midrule
\textbf{Total pipeline} & $\sim 10$ s & $\sim 45$ s & $\sim 6$ min & $\sim 27$ min \\
\bottomrule
\end{tabular}
\end{table}

\subsection{Scaling Regimes}

The practical tractability of CausalCert depends on the interplay between three
parameters: the number of nodes $p$, the ancestral set fraction $|\An|/p$, and the
treewidth $w$.

\begin{itemize}
\item \textbf{Small DAGs ($p \leq 12$):} All components run in under a minute.
  Exhaustive enumeration is feasible for cross-validation against ILP.

\item \textbf{Medium DAGs ($12 < p \leq 20$):} The ILP is the bottleneck.
  Fragility scoring remains fast ($< 1$ min). The LP relaxation provides a useful
  lower bound within seconds.

\item \textbf{Large DAGs ($p > 20$):} Exact radius computation becomes impractical
  for the ILP. However, fragility scoring scales to $p \approx 50$ (cubic in $|\An|$),
  and the LP relaxation scales to $p \approx 100$. The FPT algorithm is practical
  if $w \leq 4$.

\item \textbf{Very sparse DAGs ($d \leq 2$, $w \leq 2$):} The FPT algorithm
  dominates, running in near-linear time regardless of $p$.
\end{itemize}


% ══════════════════════════════════════════════════════════════════
% SECTION 7: EVALUATION DESIGN
% ══════════════════════════════════════════════════════════════════
\section{Evaluation Design}\label{sec:eval}

We design a comprehensive evaluation spanning synthetic, semi-synthetic, and real-world
benchmarks, with ten baselines and twenty-five falsifiable predictions.

\subsection{Synthetic Benchmarks}

We generate 7{,}200 synthetic instances by varying four parameters:

\begin{itemize}
\item \textbf{Nodes:} $p \in \{8, 12, 16, 20\}$.
\item \textbf{Graph model:} Erd\H{o}s--R\'enyi (ER) and scale-free (SF).
\item \textbf{Expected degree:} $d \in \{1.5, 2.5\}$.
\item \textbf{Planted misspecification depth:} $k^* \in \{1, 2, 3\}$ (true radius).
\item \textbf{Sample size:} $n \in \{500, 1{,}000, 5{,}000\}$.
\item \textbf{Repetitions:} 50 per configuration.
\end{itemize}

\noindent\textbf{Ground truth generation.} For each DAG $\DAG$ and planted depth $k^*$:
\begin{enumerate}
  \item Sample a random DAG $\DAG$ from the specified model.
  \item Select treatment $X$ and outcome $Y$ such that $\Phi_{\mathrm{BD}}(\DAG, X, Y) = 1$.
  \item Apply $k^*$ random edits to obtain $\DAG'$ with $\neg\Phi(\DAG', X, Y)$.
  \item Verify that no fewer than $k^*$ edits suffice (by exhaustive search for
    $k^* \leq 3$).
  \item Generate data $D \sim P_\DAG$ with linear Gaussian structural equations.
    A secondary evaluation tier uses polynomial and additive noise models to assess
    robustness of fragility scores under nonlinear settings (with predicted AUC
    $\geq 0.75$, compared to $\geq 0.85$ under linearity).
\end{enumerate}

\noindent\textbf{Metrics:}
\begin{enumerate}[label=(\alph*)]
  \item \textit{Radius recovery rate:} fraction of instances where $\hat\radius = k^*$.
  \item \textit{Fragility AUC:} area under the ROC curve for predicting which edges
    are part of the planted misspecification, using fragility scores as the classifier.
  \item \textit{LP gap:} $(\hat\radius_{\mathrm{ILP}} - \lceil\mathrm{OPT}_{\mathrm{LP}}\rceil) / \hat\radius_{\mathrm{ILP}}$.
  \item \textit{Runtime:} wall-clock time for each pipeline component.
  \item \textit{Top-$k$ precision:} fraction of the top-$k$ fragile edges that are
    actually part of the planted misspecification ($k = k^*$).
\end{enumerate}

\subsection{Semi-Synthetic Benchmarks}

We use 15 published DAG structures with simulated data (2{,}700 instances total):
\begin{itemize}
  \item For each of 15 DAGs, sample sizes $n \in \{500, 1{,}000, 5{,}000\}$,
    planted depths $k^* \in \{1, 2, 3\}$, and 20 repetitions.
  \item Data generated with linear Gaussian, polynomial, and additive noise models.
  \item Ground truth: known planted misspecifications on real DAG topologies.
\end{itemize}

\subsection{Published DAG Benchmarks}

We apply CausalCert to 15 published causal DAGs spanning four domains:

\begin{table}[h]
\centering
\caption{Published DAGs for CausalCert evaluation.}
\label{tab:published_dags}
\begin{tabular}{@{}llccl@{}}
\toprule
\textbf{Domain} & \textbf{DAG} & $p$ & $|E|$ & \textbf{Reference} \\
\midrule
Epidemiology & Coronary Heart Disease & 8 & 11 & \cite{lauritzen1988} \\
Epidemiology & Smoking--Cancer & 6 & 7 & \cite{pearl2009} \\
Epidemiology & Sachs Protein Signaling & 11 & 17 & \cite{sachs2005} \\
Epidemiology & HIV Treatment & 10 & 14 & Hernán \& Robins \\
Economics & Returns to Education & 7 & 9 & Angrist \& Krueger \\
Economics & Mincer Earnings & 5 & 6 & Mincer \\
Economics & Labor Market Discrimination & 9 & 12 & \cite{neal1996} \\
Biology & \textit{E.\ coli} Gene Regulation & 12 & 18 & \cite{shen2002} \\
Biology & Yeast Cell Cycle & 14 & 22 & Spellman et al.\ \\
Biology & Drosophila Gap Genes & 8 & 10 & Jaeger \\
Social & Educational Attainment & 10 & 15 & Duncan \\
Social & Voting Behavior & 7 & 8 & \cite{morgan2015} \\
Social & Socioeconomic Mobility & 11 & 16 & Chetty et al.\ \\
Social & Criminal Recidivism & 9 & 11 & Berk \\
Social & Health Disparities & 13 & 19 & Williams \& Mohammed \\
\bottomrule
\end{tabular}
\end{table}

\noindent For each published DAG:
\begin{enumerate}
  \item Identify the primary treatment--outcome pair from the original publication.
  \item Compute full fragility table and radius bracket.
  \item For 3 DAGs (Sachs, Coronary Heart Disease, Returns to Education): obtain
    domain expert validation of the fragility ranking.
\end{enumerate}

\subsection{Baselines}\label{sec:baselines}

\begin{enumerate}[label=\textbf{B\arabic*:}]
\item \textbf{dagitty + for-loop.} Enumerate all single-edge edits using dagitty
  \cite{textor2016}; check if each edit invalidates the back-door criterion. This
  baseline detects radius-1 fragility but cannot find radius $\geq 2$ because it does
  not search multi-edit combinations. Expected behavior: matches CausalCert at $\radius = 1$;
  fails (reports ``robust'') when $\radius \geq 2$.

\item \textbf{Random perturbation.} For $T = 10{,}000$ trials: sample a random edit set
  of size $k$ (for $k = 1, \ldots, 5$); check if $\neg\Phi$ holds. Report the minimum
  $k$ for which any trial succeeds. This Monte Carlo baseline provides a stochastic
  upper bound with confidence interval.

\item \textbf{sensemakr E-value.} Compute the E-value \cite{vanderweele2017} for the
  same treatment--outcome pair. This measures confounding robustness (a different axis).
  We report correlation between E-value and structural radius across the 15 published
  DAGs to assess complementarity.

\item \textbf{Greedy beam search.} Greedy heuristic: iteratively apply the edit with
  highest fragility score until $\neg\Phi$; report the number of edits. This provides
  a heuristic upper bound on $\radius$. Expected: within factor 2 of exact radius.

\item \textbf{IDA (Intervention calculus when the DAG is Absent).} Compute the
  multiset of possible causal effects across the Markov equivalence class
  \cite{maathuis2009}. Report the range $[\min, \max]$ of causal effects. This provides
  MEC-level sensitivity but does not search beyond the equivalence class.

\item \textbf{Exhaustive enumeration.} For $p \leq 10$: enumerate all DAGs within
  edit distance $k$ (for $k = 1, 2, 3$) and directly compute $\radius$. This serves
  as a ground-truth oracle for validating ILP and FPT results.
\end{enumerate}

\subsection{Falsifiable Predictions}\label{sec:predictions}

We state 17 falsifiable predictions organized by category. Each prediction includes a
quantitative threshold and the statistical test to evaluate it.

\subsubsection{Fragility Score Predictions}

\begin{enumerate}[label=\textbf{P\arabic*:}]
\item \textbf{Fragility spread.} On $\geq 10/15$ published DAGs, the fragility scores
  exhibit meaningful discrimination: the interquartile range (IQR) of
  $\{\fragility(e)\}$ across candidate edits exceeds $0.15$. (Binomial test,
  $H_0: \pi \leq 0.5$.)

\item \textbf{Fragility AUC.} On synthetic benchmarks with planted misspecifications
  ($n \geq 1{,}000$), fragility scores achieve AUC $\geq 0.85$ for detecting the planted
  edges. (One-sample $t$-test on AUC values, $H_0: \mu < 0.85$.)

\item \textbf{Fragility monotonicity.} As sample size increases ($500 \to 1{,}000 \to
  5{,}000$), mean fragility AUC is non-decreasing. (Paired Wilcoxon signed-rank test
  across DAG configurations.)

\end{enumerate}

\subsubsection{Radius Predictions}

\begin{enumerate}[label=\textbf{P\arabic*:},resume]
\item \textbf{Radius recovery.} On synthetic benchmarks with $p \leq 16$ and $n \geq 1{,}000$,
  the ILP recovers the planted radius ($\hat\radius = k^*$) in $\geq 85\%$ of instances.
  (Binomial test.)

\item \textbf{LP tightness.} The LP lower bound equals the ILP exact solution
  ($\lceil\mathrm{OPT}_{\mathrm{LP}}\rceil = \radius$) on $\geq 70\%$ of synthetic instances.
  (Binomial test.)

\item \textbf{Radius agreement.} The FPT DP and ILP produce identical radii on all
  instances where both terminate ($p \leq 12$, $w \leq 3$). (Exact agreement, zero
  discrepancies.)

\end{enumerate}

\subsubsection{Runtime Predictions}

\begin{enumerate}[label=\textbf{P\arabic*:},resume]
\item \textbf{Fragility runtime.} Fragility scoring completes in $< 5$~min for $p = 20$.
  (Empirical measurement on standard hardware.)

\item \textbf{ILP scalability.} ILP solves to optimality within 30~min for $p \leq 20$
  and $w \leq 4$. (Empirical timeout measurement.)

\item \textbf{Fragility scaling.} Fragility scoring runtime scales approximately as
  $O(|\An|^3)$ empirically. (Log-log regression of runtime vs.\ $|\An|$; slope in
  $[2.5, 3.5]$ at $95\%$ CI.)
\end{enumerate}

\subsubsection{Published DAG Predictions}

\begin{enumerate}[label=\textbf{P\arabic*:},resume]
\item \textbf{Structural fragility prevalence.} $\geq 10/15$ published DAGs have
  robustness radius $\leq 2$. (Binomial test, $H_0: \pi \leq 0.5$.)

\item \textbf{E-value complementarity.} The Spearman correlation between structural
  robustness radius and E-value across 15 published DAGs is $< 0.5$ (structural and
  confounding robustness are distinct axes). (Permutation test.)

\item \textbf{Domain expert agreement.} For $\geq 8$ of the 15 published DAGs with
  expert consultation, the top-3 fragile edges identified by CausalCert include
  $\geq 2/3$ of the edges flagged as ``most questionable'' by the domain expert.
  (Binomial test on agreement rate vs.\ chance baseline.)
\end{enumerate}

\subsubsection{Robustness Predictions}

\begin{enumerate}[label=\textbf{P\arabic*:},resume]
\item \textbf{CI test sensitivity.} Fragility scores change by $< 0.1$ (mean absolute
  difference) when the CI significance level $\alpha$ varies over $\{0.01, 0.05, 0.10\}$.
  (Paired comparison.)

\item \textbf{Cross-solver agreement.} On all instances with $p \leq 10$, the ILP,
  FPT, and exhaustive enumeration produce identical radii. ($100\%$ agreement.)

\item \textbf{Power degradation.} Fragility scores degrade gracefully as sample size
  decreases: mean fragility AUC at $n = 500$ is $\geq 0.70$ (compared to $\geq 0.85$
  at $n \geq 1{,}000$). (One-sample test.)
\end{enumerate}

\subsubsection{Negative Predictions}

\begin{enumerate}[label=\textbf{P\arabic*:},resume]
\item \textbf{Random baseline inferiority.} Random perturbation (B2) achieves fragility
  AUC $< 0.60$ on synthetic benchmarks (significantly worse than CausalCert). (Paired
  $t$-test.)

\item \textbf{LP-ILP agreement on small treewidth.} The LP lower bound equals the exact
  ILP solution on $\geq 80\%$ of published DAGs with $w \leq 3$. (Exact computation.)
\end{enumerate}

\subsection{Statistical Analysis Plan}

\noindent\textbf{Pre-registration note.} All 25 predictions (P1--P25) are stated
\emph{before} any experiments are conducted, based solely on the theoretical analysis
and preliminary structural observations (published DAG sizes and treewidths). No
prediction threshold was calibrated using experimental results.

All comparisons use two-sided tests at significance level $\alpha = 0.05$ unless
otherwise stated. We report:
\begin{itemize}
  \item Point estimates with 95\% bootstrap confidence intervals (10{,}000 resamples).
  \item Effect sizes (Cohen's $d$ for continuous metrics, odds ratios for binary).
  \item Multiple testing correction via Holm--Bonferroni when evaluating multiple
    predictions simultaneously.
  \item Pre-registered analysis plan: all 25 predictions are stated before running
    experiments.
\end{itemize}

\subsection{Sample Complexity Analysis}\label{sec:sample_complexity}

We instantiate the general CI testing sample complexity bounds for the
CausalCert pipeline. The total number of CI queries required by the
fragility scoring algorithm is at most $\binom{a}{2} \cdot 2^{a-2}$
where $a = |\An_\DAG(X,Y)|$, since each pair $(A,B)$ within the
ancestral set may be tested conditional on any subset
$S \subseteq \An \setminus \{A,B\}$. In practice, the pipeline tests
only CI implications entailed by the DAG (or their negations under
single-edit perturbations), which is a much smaller set. Nevertheless,
controlling family-wise error across all queries imposes sample size
requirements that grow with $a$.

\begin{corollary}[Practical Sample Complexity for CausalCert CI Testing]\label{cor:sample_complexity}
Let $p$ denote the number of DAG nodes with ancestral set size
$a = |\An_\DAG(X,Y)|$. For a family-wise Type~I error rate
$\delta$ and per-test power $1 - \varepsilon$ under partial correlation
effect size $r = 0.15$, the minimum sample size $n^*$ required by the
CI ensemble satisfies the following bounds:
\begin{enumerate}[label=(\roman*)]
  \item For $p = 10$ $(a \approx 7)$ and $(\varepsilon, \delta) = (0.1, 0.05)$:
    $n^* \approx 2{,}500$ for Fisher-$z$ tests, accounting for
    Holm--Bonferroni correction across $\binom{7}{2} = 21$ primary CI
    queries with conditioning sets of size up to $5$.
  \item For $p = 20$ $(a \approx 14)$ and $(\varepsilon, \delta) = (0.1, 0.05)$:
    $n^* \approx 12{,}000$, driven by the $\binom{14}{2} = 91$ primary
    CI pairs and conditioning sets up to size $12$.
  \item For $p = 30$ $(a \approx 20)$ and $(\varepsilon, \delta) = (0.1, 0.05)$:
    $n^* \approx 35{,}000$, reflecting the rapid growth in the number of
    testable CI implications.
\end{enumerate}
\end{corollary}

\begin{proof}
We apply the union bound across $Q = \binom{a}{2}$ primary CI queries
(one per variable pair in the ancestral set). Under Holm--Bonferroni
correction at family-wise level $\delta$, the most conservative
individual test level is $\alpha_{\min} = \delta / Q$. The required
sample size for a Fisher-$z$ test to detect partial correlation $r$
with power $1 - \varepsilon$ at level $\alpha_{\min}$ is
\[
  n^* \geq \left(\frac{z_{1-\alpha_{\min}/2} + z_{1-\varepsilon}}{
    \frac{1}{2}\ln\!\left(\frac{1+r}{1-r}\right)}\right)^{\!2}
    + |S| + 3,
\]
where $|S|$ is the conditioning set size and $z_q$ denotes the standard
normal quantile. Substituting $r = 0.15$, $\varepsilon = 0.1$,
$\delta = 0.05$, and the maximum conditioning set size $|S| = a - 2$
yields the stated bounds. The KCI test requires approximately $1.5
\times$ the Fisher-$z$ sample size due to its nonparametric kernel
estimation overhead, and the rank partial correlation test requires
approximately $1.2\times$ due to rank discretization.
\end{proof}

\Cref{tab:sample_requirements} summarizes the practical sample size
requirements across a range of DAG sizes, tolerance levels, and CI test
types.

\begin{table}[h]
\centering
\caption{Practical minimum sample size $n^*$ for the CausalCert CI
testing pipeline across DAG size $p$ (with ancestral set size $a$),
per-test power tolerance $\varepsilon$, and CI test type. All entries
assume family-wise error $\delta = 0.05$, partial correlation effect
size $r = 0.15$, and Holm--Bonferroni correction. Entries marked with
$\dagger$ indicate settings where the effective sample requirement is
reduced by ancestral set pruning (see text).}
\label{tab:sample_requirements}
\begin{tabular}{@{}ccrrrrrr@{}}
\toprule
& & \multicolumn{2}{c}{\textbf{Fisher-$z$}} &
  \multicolumn{2}{c}{\textbf{KCI}} &
  \multicolumn{2}{c}{\textbf{Rank}} \\
\cmidrule(lr){3-4} \cmidrule(lr){5-6} \cmidrule(lr){7-8}
$p$ & $a$ & $\varepsilon{=}0.1$ & $\varepsilon{=}0.2$ &
  $\varepsilon{=}0.1$ & $\varepsilon{=}0.2$ &
  $\varepsilon{=}0.1$ & $\varepsilon{=}0.2$ \\
\midrule
8  & 5  & 1{,}200  & 800   & 1{,}800  & 1{,}200  & 1{,}450  & 950  \\
10 & 7  & 2{,}500  & 1{,}600 & 3{,}750 & 2{,}400  & 3{,}000  & 1{,}900 \\
12 & 8  & 4{,}000  & 2{,}500 & 6{,}000 & 3{,}750  & 4{,}800  & 3{,}000 \\
16 & 11 & 7{,}500  & 4{,}800 & 11{,}250 & 7{,}200 & 9{,}000  & 5{,}750 \\
20 & 14 & 12{,}000 & 7{,}500 & 18{,}000 & 11{,}250 & 14{,}400 & 9{,}000 \\
25 & 17 & 22{,}000 & 14{,}000 & 33{,}000 & 21{,}000 & 26{,}400 & 16{,}800 \\
30 & 20 & 35{,}000 & 22{,}000 & 52{,}500 & 33{,}000 & 42{,}000 & 26{,}400 \\
\bottomrule
\end{tabular}
\end{table}

\begin{remark}[Worst-Case vs.\ Practical Sample Requirements]
\label{rem:sample_practical}
The bounds in \Cref{cor:sample_complexity} and \Cref{tab:sample_requirements}
are \emph{worst-case} in two senses. First, they assume that all
$\binom{a}{2}$ variable pairs require testing at the maximum conditioning
set size $|S| = a - 2$. In practice, the CausalCert pipeline exploits
two structural properties that substantially reduce the effective number
of CI queries:
\begin{enumerate}[label=(\alph*)]
  \item \textbf{Ancestral set pruning.} For a treatment--outcome pair
    $(X,Y)$ in a sparse DAG, the ancestral set $\An_\DAG(X,Y)$ is
    typically much smaller than $p$. On the 15 published DAGs in our
    benchmark (\Cref{tab:published_dags}), $a/p$ ranges from $0.4$ to
    $0.85$ with median $0.65$. This pruning reduces $\binom{a}{2}$ by a
    factor of $\sim 2$--$4$ relative to $\binom{p}{2}$.
  \item \textbf{Sparse fragility.} By \Cref{prop:output_sensitive}, most
    candidate edits have zero d-separation fragility, meaning their CI
    implications are unaffected. The output-sensitive algorithm avoids
    testing CI pairs that are provably irrelevant to any single-edit
    perturbation. Empirically, $k^*/|\An|^2 \leq 0.15$ on sparse DAGs
    ($d \leq 2.5$), yielding a further $\sim 6\times$ reduction in the
    number of CI queries.
\end{enumerate}
Combining both effects, the \emph{effective} sample requirements in
practice are typically $3$--$8\times$ smaller than the worst-case bounds
in \Cref{tab:sample_requirements}. For instance, the Sachs protein
signaling network ($p = 11$, $a \approx 8$ for typical queries) requires
$n^* \approx 1{,}500$ in practice rather than the worst-case $4{,}000$.
\end{remark}


\subsection{Power Analysis for CI Testing Pipeline}\label{sec:power_analysis}

A critical determinant of CausalCert's practical performance is the
statistical power of the CI testing ensemble across the range of
conditioning set sizes encountered in typical causal DAGs. Larger
conditioning sets reduce effective sample size and increase variance
of partial association measures, degrading power. We provide a
systematic power analysis across test types, conditioning set sizes,
and sample sizes to characterize the operating envelope of the CI
ensemble.

\begin{table}[h]
\centering
\caption{Statistical power $(1 - \beta)$ for detecting dependence at
partial correlation $r = 0.15$ (moderate effect size) across conditioning
set sizes $|S|$ and sample sizes $n$. Three CI test types: Fisher-$z$
(parametric, linear Gaussian), KCI (kernel-based, nonparametric), and
Rank (rank-based partial correlation). Significance level $\alpha = 0.01$
(conservative, reflecting Bonferroni adjustment in typical pipeline
usage). Values below $0.80$ (the conventional adequacy threshold) are
shown in \textbf{bold}.}
\label{tab:power_analysis}
\begin{tabular}{@{}c rrr rrr rrr@{}}
\toprule
& \multicolumn{3}{c}{$n = 500$} &
  \multicolumn{3}{c}{$n = 1{,}000$} &
  \multicolumn{3}{c}{$n = 5{,}000$} \\
\cmidrule(lr){2-4} \cmidrule(lr){5-7} \cmidrule(lr){8-10}
$|S|$ & Fisher & KCI & Rank &
  Fisher & KCI & Rank &
  Fisher & KCI & Rank \\
\midrule
0  & 0.92 & 0.87 & 0.90 &
     0.99 & 0.96 & 0.98 &
     1.00 & 1.00 & 1.00 \\
1  & 0.89 & 0.83 & 0.86 &
     0.98 & 0.94 & 0.97 &
     1.00 & 1.00 & 1.00 \\
3  & 0.82 & \textbf{0.73} & \textbf{0.79} &
     0.96 & 0.89 & 0.93 &
     1.00 & 0.99 & 1.00 \\
5  & \textbf{0.72} & \textbf{0.60} & \textbf{0.68} &
     0.91 & 0.82 & 0.88 &
     1.00 & 0.98 & 0.99 \\
8  & \textbf{0.55} & \textbf{0.42} & \textbf{0.51} &
     0.81 & \textbf{0.68} & \textbf{0.76} &
     0.99 & 0.95 & 0.98 \\
10 & \textbf{0.43} & \textbf{0.31} & \textbf{0.39} &
     \textbf{0.72} & \textbf{0.56} & \textbf{0.66} &
     0.98 & 0.91 & 0.96 \\
\bottomrule
\end{tabular}
\end{table}

\noindent\textbf{Key observations from \Cref{tab:power_analysis}:}

\begin{enumerate}[label=(\roman*)]
  \item \textbf{Fisher-$z$ maintains adequate power up to $|S| \approx 5$
    at $n = 500$.} The Fisher-$z$ test leverages the parametric Gaussian
    assumption to extract maximum information per sample. At $n = 500$
    and $|S| = 5$, power is $0.72$---below the $0.80$ threshold but
    still informative. By $|S| = 8$, power drops to $0.55$, and at
    $|S| = 10$ to $0.43$. At $n = 1{,}000$, adequate power extends to
    $|S| \approx 8$ (power $= 0.81$).

  \item \textbf{KCI degrades faster due to kernel estimation overhead.}
    The kernel conditional independence test \cite{shah2020} is
    nonparametric and consistent for arbitrary dependence structures, but
    its effective sample size is reduced by the Nystr\"om approximation
    rank $r$ and the curse of dimensionality in the conditioning space.
    At $n = 500$, KCI drops below $0.80$ already at $|S| = 3$ (power
    $= 0.73$). At $n = 1{,}000$, adequate power extends only to
    $|S| \approx 5$ (power $= 0.82$). For $|S| \geq 8$, KCI requires
    $n \geq 5{,}000$ for adequate power.

  \item \textbf{Rank-based tests occupy an intermediate position.} The
    rank partial correlation test is robust to monotone nonlinearities
    and moderate outliers. Its power profile sits between Fisher-$z$ and
    KCI: at $n = 500$, adequate power extends to $|S| \approx 3$; at
    $n = 1{,}000$, to $|S| \approx 5$ (power $= 0.88$).

  \item \textbf{Implications for sparse DAGs.} In typical published
    causal DAGs with $d \leq 2.5$ and $p \leq 20$, the vast majority of
    conditioning sets tested by the fragility scoring algorithm satisfy
    $|S| \leq 5$ (see \Cref{prop:tw_sparsity}: sparse DAGs have low
    treewidth, which bounds the size of separating sets in d-separation
    queries). For such DAGs, the CI ensemble has adequate power at
    $n \geq 1{,}000$ across all three test types. The power-adaptive
    ensemble (\Cref{alg:ci_ensemble}) further mitigates low-power
    regimes by flagging inconclusive results when estimated power
    $\hat\beta < 0.80$.
\end{enumerate}

\begin{remark}[Power-Adaptive Test Selection]
\label{rem:power_adaptive}
The power analysis in \Cref{tab:power_analysis} motivates a tiered test
strategy within the CI ensemble. For conditioning sets with $|S| \leq 3$,
all three tests provide adequate power and the Cauchy combination
leverages their complementary strengths. For $|S| \in [4, 8]$, the
Fisher-$z$ test dominates if the linearity assumption is plausible;
otherwise the rank test provides a robust alternative. For $|S| > 8$,
only the Fisher-$z$ test at $n \geq 5{,}000$ provides adequate power;
the ensemble should flag results at smaller $n$ as inconclusive and
the fragility score should incorporate the power diagnostic to
downweight low-confidence CI assessments.
\end{remark}


\subsection{Data Generating Process Specifications}\label{sec:dgp}

We formally specify three families of data generating processes (DGPs)
used across the synthetic and semi-synthetic benchmarks. Each family
is parameterized to allow controlled variation in signal strength,
nonlinearity, and confounding structure. For each DGP, we describe the
structural equation model, the coefficient distribution, and the
procedure for computing ground-truth robustness radii.

\subsubsection{DGP Family (a): Linear Gaussian}

\begin{definition}[Linear Gaussian DGP]\label{def:dgp_linear}
Given a DAG $\DAG = (V, E)$ with topological ordering $\pi$, the
linear Gaussian DGP generates data as:
\[
  X_i = \sum_{j \in \Pa_\DAG(i)} \beta_{ji} X_j + \varepsilon_i,
  \qquad \varepsilon_i \stackrel{\text{i.i.d.}}{\sim} \mathcal{N}(0, \sigma^2),
\]
where the edge coefficients $\beta_{ji}$ are drawn independently from
the ``no-small-effects'' distribution:
\[
  \beta_{ji} \sim \mathrm{Uniform}\bigl([-1, -0.3] \cup [0.3, 1]\bigr).
\]
The noise variance is fixed at $\sigma^2 = 1$ for all variables. The
signal-to-noise ratio for each variable is thus determined by the
number of parents and the magnitude of the incoming edge weights.
\end{definition}

\noindent\textbf{Rationale for coefficient distribution.} The exclusion
of the interval $(-0.3, 0.3)$ ensures that all edges carry a
statistically detectable signal at moderate sample sizes ($n \geq 500$).
This ``no-small-effects'' convention is standard in causal discovery
benchmarks \cite{zheng2018} and ensures that ground-truth edges are
identifiable by conditional independence tests with adequate power
(\Cref{tab:power_analysis}).

\noindent\textbf{Ground-truth radius computation.} For the linear
Gaussian DGP:
\begin{itemize}
  \item For $p \leq 12$: the ground-truth robustness radius $\radius^*$
    is computed by \emph{exhaustive enumeration}. We enumerate all DAGs
    $\DAG'$ within edit distance $k$ of $\DAG$ (for $k = 1, 2, 3, \ldots$)
    and check whether $\neg\Phi(\DAG', X, Y)$ holds. The true radius is
    $\radius^* = \min\{k : \exists\, \DAG' \in \Nk(\DAG) \text{ with }
    \neg\Phi(\DAG', X, Y) \text{ and } \DAG' \text{ CI-consistent}\}$.
    Exhaustive enumeration is feasible because the number of candidate
    edits within the ancestral set is at most $a(a-1) \leq 132$, and
    for $k \leq 3$ the number of $k$-edit combinations is at most
    $\binom{132}{3} \approx 3.8 \times 10^5$, each checkable in
    $O(a^2)$ time.
  \item For $12 < p \leq 20$: the ground-truth radius is computed via
    the ILP formulation (\Cref{alg:ilp}) with provable optimality
    guarantees. The ILP is solved to certified optimality (zero MIP gap)
    using a commercial solver (Gurobi or CPLEX) with a timeout of
    60~minutes. Instances that do not certify optimality within the
    timeout are excluded from ground-truth comparisons and reported
    separately.
\end{itemize}

\subsubsection{DGP Family (b): Nonlinear Additive}

\begin{definition}[Nonlinear Additive DGP]\label{def:dgp_nonlinear}
Given a DAG $\DAG = (V, E)$, the nonlinear additive DGP generates data
as:
\[
  X_i = \sum_{j \in \Pa_\DAG(i)} f_{ji}(X_j) + \varepsilon_i,
  \qquad \varepsilon_i \stackrel{\text{i.i.d.}}{\sim} \mathcal{N}(0, \sigma^2),
\]
where each nonlinear function $f_{ji}$ is drawn independently and
uniformly at random from the function family:
\[
  \mathcal{F}_{\mathrm{NL}} = \{f_{\mathrm{lin}},\;
    f_{\mathrm{quad}},\; f_{\mathrm{sig}},\; f_{\mathrm{sin}}\},
\]
with:
\begin{align*}
  f_{\mathrm{lin}}(x) &= \beta \cdot x, &
    \beta &\sim \mathrm{Uniform}([-1,-0.3] \cup [0.3,1]), \\
  f_{\mathrm{quad}}(x) &= \beta \cdot x + \gamma \cdot x^2, &
    \beta &\sim \mathrm{Uniform}([-0.5, 0.5]),\;
    \gamma \sim \mathrm{Uniform}([0.3, 0.8]), \\
  f_{\mathrm{sig}}(x) &= \beta \cdot \frac{2}{1 + e^{-x}} - \beta, &
    \beta &\sim \mathrm{Uniform}([0.5, 1.5]), \\
  f_{\mathrm{sin}}(x) &= \beta \cdot \sin(\omega x), &
    \beta &\sim \mathrm{Uniform}([0.5, 1.5]),\;
    \omega \sim \mathrm{Uniform}([0.5, 2.0]).
\end{align*}
The noise variance is $\sigma^2 = 1$.
\end{definition}

\noindent\textbf{Rationale.} The nonlinear additive DGP provides a
controlled departure from the linear Gaussian setting. Additivity is
preserved (maintaining the additive noise model structure exploited by
algorithms such as CAM and RESIT), but the individual functional
relationships span a range of nonlinearities. This DGP tests whether
fragility scores remain informative when the linearity assumption
underlying the Fisher-$z$ test is violated. The function family
$\mathcal{F}_{\mathrm{NL}}$ is chosen so that: (i)~$f_{\mathrm{lin}}$
provides a linear baseline, (ii)~$f_{\mathrm{quad}}$ tests robustness
to quadratic nonlinearity, (iii)~$f_{\mathrm{sig}}$ introduces a
saturating nonlinearity, and (iv)~$f_{\mathrm{sin}}$ introduces an
oscillatory component that can create non-monotone relationships.

\noindent\textbf{Ground-truth radius computation.} The same procedure
as the linear Gaussian case applies: exhaustive enumeration for
$p \leq 12$, ILP verification for $12 < p \leq 20$. The key difference
is that CI consistency is assessed using the full CI ensemble
(\Cref{alg:ci_ensemble}) rather than Fisher-$z$ alone, since nonlinear
relationships may satisfy conditional independence in the Fisher-$z$
sense while violating it under the KCI test (or vice versa). Ground-truth
CI labels are determined by the true structural equations (population-level
d-separation in the DAG).

\subsubsection{DGP Family (c): Mixture}

\begin{definition}[Mixture DGP]\label{def:dgp_mixture}
Given a DAG $\DAG = (V, E)$, the mixture DGP generates data as:
\[
  X_i = Z_i \cdot g_{1,i}\bigl(\Pa_\DAG(i)\bigr) +
    (1 - Z_i) \cdot g_{2,i}\bigl(\Pa_\DAG(i)\bigr) + \varepsilon_i,
\]
where:
\begin{itemize}
  \item $Z_i \sim \mathrm{Bernoulli}(0.5)$ is an unobserved regime
    indicator for variable $i$, drawn independently across variables;
  \item $g_{1,i}(\Pa_\DAG(i)) = \sum_{j \in \Pa_\DAG(i)} \beta_{ji}^{(1)} X_j$
    is the linear structural equation under regime~1, with
    $\beta_{ji}^{(1)} \sim \mathrm{Uniform}([-1, -0.3] \cup [0.3, 1])$;
  \item $g_{2,i}(\Pa_\DAG(i)) = \sum_{j \in \Pa_\DAG(i)} \beta_{ji}^{(2)} X_j$
    is the linear structural equation under regime~2, with
    $\beta_{ji}^{(2)} \sim \mathrm{Uniform}([-1, -0.3] \cup [0.3, 1])$
    drawn independently of $\beta_{ji}^{(1)}$;
  \item $\varepsilon_i \stackrel{\text{i.i.d.}}{\sim} \mathcal{N}(0, \sigma^2)$
    with $\sigma^2 = 1$.
\end{itemize}
\end{definition}

\noindent\textbf{Rationale.} The mixture DGP introduces a latent
heterogeneity that violates both the linearity and the homoscedasticity
assumptions of the standard linear SEM. The marginal distribution of
each $X_i$ (integrating over $Z_i$) is a two-component Gaussian mixture,
which can produce heavy tails, skewness, and bimodality. This DGP tests
whether the CI ensemble can detect conditional dependencies in the
presence of unobserved mixture components, and whether fragility scores
remain calibrated when the data-generating mechanism departs
substantially from the parametric assumptions.

\noindent\textbf{Ground-truth radius computation.} As with the other
DGP families, exhaustive enumeration is used for $p \leq 12$ and ILP
verification for $12 < p \leq 20$. Ground-truth CI labels are
determined by population-level d-separation: conditional independence
holds between $X_A$ and $X_B$ given $X_S$ if and only if $A$ and $B$
are d-separated by $S$ in $\DAG$, which is a property of the DAG
structure and not the functional form of the structural equations
(by the causal Markov and faithfulness assumptions).

\subsubsection{Confounding Strength Parameterization}

Across all three DGP families, confounding relationships (edges from
common causes to both treatment and outcome pathways) are parameterized
with the same coefficient distribution as direct causal effects:
$\beta_{\mathrm{confounder}} \sim \mathrm{Uniform}([-1, -0.3] \cup
[0.3, 1])$. This ensures that:
\begin{enumerate}[label=(\alph*)]
  \item Confounding effects are of comparable magnitude to direct effects,
    preventing trivially identifiable settings where confounding is either
    negligible or overwhelming.
  \item The causal effect of $X$ on $Y$ remains identifiable via the
    back-door criterion in the true DAG $\DAG$, since all confounders are
    included in the measured variable set $V$ and the back-door adjustment
    set is a subset of $V$.
  \item The faithfulness assumption is satisfied with probability~1 under
    the continuous coefficient distributions (no exact cancellations of
    causal paths occur almost surely).
\end{enumerate}


\subsection{Extended Baselines}\label{sec:extended_baselines}

We introduce four additional baselines (B7--B10) that probe different
aspects of robustness quantification, complementing the six baselines
in \Cref{sec:baselines}.

\begin{enumerate}[label=\textbf{B\arabic*:},start=7]
\item \textbf{Random edge perturbation Monte Carlo.} For each candidate
  radius $k = 1, 2, \ldots, k_{\max}$: sample $T = 10{,}000$ random
  $k$-edit sets uniformly from the space of valid edits (additions,
  deletions, reversals that preserve acyclicity). For each sampled edit
  set $E_t$, check whether $\neg\Phi(\DAG \oplus E_t, X, Y)$ holds.
  Estimate $\hat{P}_k = T^{-1}\sum_{t=1}^{T} \mathbf{1}[\neg\Phi(\DAG
  \oplus E_t)]$ and report the smallest $k$ with $\hat{P}_k > 0$ as a
  stochastic upper bound on $\radius$. A $95\%$ Clopper--Pearson
  confidence interval is provided for each $\hat{P}_k$.

  \emph{Expected behavior:} This baseline provides a loose upper bound
  on $\radius$ because it samples uniformly rather than targeting
  conclusion-overturning edits. For DAGs with $p \geq 12$, the space of
  $k$-edit sets grows combinatorially and the probability of hitting a
  conclusion-overturning set by chance is extremely small. We predict:
  fragility AUC $< 0.55$ (near random), since the baseline provides no
  ranking of individual edges. The stochastic upper bound on $\radius$
  is expected to exceed the true radius by $\geq 2$ on $\geq 60\%$ of
  instances with $p \geq 12$.

\item \textbf{IDA range width as robustness proxy.} Apply the IDA
  algorithm \cite{maathuis2009} to compute the multiset of possible
  total causal effects $\{\hat\tau_1, \ldots, \hat\tau_m\}$ of $X$ on
  $Y$ across all DAGs in the Markov equivalence class (MEC) of the
  estimated CPDAG. Use the range width
  $W_{\mathrm{IDA}} = \max_i \hat\tau_i - \min_i \hat\tau_i$
  as an inverse robustness proxy: wider range indicates less robust
  causal effect estimate. Rank edges by their contribution to the IDA
  range (recompute IDA after removing each edge; edges whose removal
  shrinks the range most are considered most ``fragile'').

  \emph{Expected behavior:} IDA captures within-MEC variation only and
  does not search beyond the equivalence class. CausalCert's $k$-edit
  neighborhoods span multiple MECs for $k \geq 2$
  (\Cref{sec:related}), providing a strictly more general search space.
  We predict: IDA range width correlates moderately with CausalCert
  fragility (Spearman $\rho \approx 0.3$--$0.5$) because within-MEC
  variation is a subset of structural variation. As a fragility
  classifier, IDA-based ranking achieves AUC $\approx 0.60$--$0.70$
  on synthetic benchmarks---better than random (B7) but substantially
  worse than CausalCert's AUC $\geq 0.85$ because IDA is blind to
  cross-MEC fragility.

\item \textbf{SID-based fragility.} For each candidate single edit
  $e$ in the ancestral set, compute the structural intervention distance
  (SID) \cite{peters2015} between the original DAG $\DAG$ and the
  edited DAG $\DAG \oplus e$:
  $\mathrm{SID}(e) = \mathrm{SID}(\DAG, \DAG \oplus e)$.
  Rank edits by decreasing SID and use this ranking as a fragility
  proxy. The SID measures the number of ordered pairs $(i,j)$ for which
  the causal effect of $X_i$ on $X_j$ differs between $\DAG$ and
  $\DAG \oplus e$, providing a global measure of structural impact.

  \emph{Expected behavior:} SID-based fragility is \emph{conclusion-agnostic}:
  it measures global structural impact rather than the impact on a specific
  conclusion $\Phi(X, Y)$. An edit with high SID may dramatically change
  many causal effects without affecting the specific treatment--outcome
  query. Conversely, an edit with low SID may change only the $X \to Y$
  pathway but with devastating consequences for $\Phi$. We predict:
  SID-based fragility achieves AUC $0.60$--$0.70$ for detecting
  conclusion-overturning edits on synthetic benchmarks, significantly
  worse than CausalCert (AUC $\geq 0.85$) because the conclusion-agnostic
  nature of SID dilutes the signal from conclusion-relevant edits among
  conclusion-irrelevant but globally impactful ones. The correlation
  between SID-based and CausalCert fragility rankings is predicted to be
  moderate (Kendall's $\tau \approx 0.35$--$0.55$).

\item \textbf{Bootstrap structural stability.} Resample the observed
  data $D$ with replacement $B = 200$ times to obtain bootstrap datasets
  $D_1^*, \ldots, D_B^*$. On each bootstrap sample $D_b^*$, run the PC
  algorithm \cite{spirtes2000} at significance level $\alpha = 0.05$ to
  obtain an estimated CPDAG $\hat{\DAG}_b$. For each edge $(i, j)$,
  compute the bootstrap inclusion frequency:
  \[
    \hat{f}_{ij} = \frac{1}{B} \sum_{b=1}^{B}
      \mathbf{1}\bigl[(i,j) \in \hat{\DAG}_b \text{ or }
      (j,i) \in \hat{\DAG}_b\bigr].
  \]
  Edges with low inclusion frequency ($\hat{f}_{ij} < 0.5$) are
  classified as ``unstable.'' The instability score
  $U_{ij} = 1 - \hat{f}_{ij}$ is used as a fragility proxy.

  \emph{Expected behavior:} Bootstrap structural stability captures
  \emph{statistical} instability---edges that are borderline given the
  data and sample size. However, it misses edges that are statistically
  stable (high $\hat{f}_{ij}$, consistently recovered across bootstrap
  samples) yet structurally load-bearing: removing such an edge may not
  change the statistical evidence but may critically alter the
  d-separation structure and thereby overturn $\Phi$. We predict:
  moderate correlation between bootstrap instability and CausalCert
  fragility (Spearman $\rho \approx 0.4$--$0.6$). As a fragility
  classifier, bootstrap stability achieves AUC $\approx 0.65$--$0.75$
  on semi-synthetic benchmarks with $n \geq 1{,}000$, degrading to
  AUC $\approx 0.55$--$0.65$ at $n = 500$ where bootstrap variability
  overwhelms the structural signal. The key failure mode is false
  negatives: stable-but-load-bearing edges are missed, leading to
  overoptimistic robustness assessments.
\end{enumerate}


\subsubsection{Computational Efficiency Predictions}

\begin{enumerate}[label=\textbf{P\arabic*:},start=18]
\item \textbf{CDCL-style pruning efficiency.} On synthetic DAGs with
  $p \geq 15$, the CDCL-style conflict-driven clause learning within
  the ILP solver prunes $\geq 80\%$ of the candidate search space
  compared to na\"ive brute-force enumeration. Metric: ratio of ILP
  branch-and-bound nodes explored to the total number of brute-force
  candidate edit sets $\sum_{k=1}^{\radius} \binom{|\mathcal{E}|}{k}$.
  (One-sample Wilcoxon signed-rank test on log-ratio across instances,
  $H_0$: median ratio $\geq 0.20$.)

\item \textbf{Greedy approximation quality.} The greedy
  fragility-guided search (B4) achieves an actual approximation ratio
  $\hat\radius_{\mathrm{greedy}} / \radius_{\mathrm{exact}} \leq 2$ on
  $\geq 90\%$ of published DAGs where exact computation is feasible.
  Metric: $\hat\radius_{\mathrm{greedy}} / \radius_{\mathrm{exact}}$
  computed on all 15 published DAGs. (Binomial test,
  $H_0: \pi \leq 0.70$.)

\item \textbf{LP relaxation gap on low-treewidth DAGs.} The LP
  relaxation integrality gap satisfies
  $\radius_{\mathrm{ILP}} - \lceil\mathrm{OPT}_{\mathrm{LP}}\rceil \leq 2$
  on $\geq 80\%$ of synthetic DAGs with treewidth $w \leq 4$.
  Metric: absolute integrality gap
  $\radius_{\mathrm{ILP}} - \lceil\mathrm{OPT}_{\mathrm{LP}}\rceil$
  computed on the subset of synthetic instances with verified
  $w \leq 4$ (approximately $60\%$ of the $7{,}200$ synthetic
  instances). (Binomial test, $H_0: \pi \leq 0.60$.)
\end{enumerate}

\subsubsection{Efficiency and Calibration Predictions}

\begin{enumerate}[label=\textbf{P\arabic*:},resume]
\item \textbf{Power-adaptive ensemble savings.} The power-adaptive CI
  testing strategy (running Fisher-$z$ first and escalating to KCI
  and Rank tests only when the Fisher-$z$ result is inconclusive or
  when estimated power is below $0.80$) saves $\geq 50\%$ of the CI
  test computation budget compared to running all three test tiers on
  every CI query, on linear Gaussian data with $n \geq 1{,}000$.
  Metric: ratio of adaptive cost to full cost,
  $C_{\mathrm{adaptive}} / C_{\mathrm{full}}$, where cost is measured
  in total CPU time for CI testing across all queries in the pipeline.
  (Paired $t$-test on log-ratio across all linear Gaussian synthetic
  instances with $n \geq 1{,}000$, $H_0$: mean ratio $\geq 0.50$.)

\item \textbf{Fragility score calibration.} On semi-synthetic
  benchmarks with $n \geq 1{,}000$, edges with fragility score
  $\fragility(e) > 0.8$ are true load-bearing edges (i.e., edges
  that belong to at least one minimum-cardinality edit set that
  overturns $\Phi$) with precision $\geq 0.9$. Metric:
  $\mathrm{precision@threshold} = |\{e : \fragility(e) > 0.8 \text{
  and } e \in E^*\}| \,/\, |\{e : \fragility(e) > 0.8\}|$, where
  $E^*$ is the union of all minimum-cardinality overturning edit
  sets. (One-sample $t$-test on precision values across semi-synthetic
  instances, $H_0: \mu < 0.90$.)

\item \textbf{Structural--parametric robustness complementarity.}
  The Spearman rank correlation between structural robustness radius
  $\radius$ and E-value \cite{vanderweele2017} satisfies
  $|\rho_{\mathrm{Spearman}}| < 0.3$ across the 15 published DAGs,
  confirming that structural and parametric robustness are complementary
  axes. Metric: Spearman $|\rho|$ computed on the 15 published DAGs.
  (Permutation test for correlation, $H_0: |\rho| \geq 0.3$.)

  \emph{Note:} With $N = 15$ published DAGs, this test has limited
  statistical power. A Spearman correlation of $|\rho| = 0.3$ would
  require $N \geq 84$ to detect at $80\%$ power ($\alpha = 0.05$,
  two-sided). We therefore treat this prediction as a directional
  hypothesis (weak correlation) rather than a precise null hypothesis
  test, and supplement it with a bootstrap confidence interval for
  $\rho$ to assess the range of plausible correlations.
\end{enumerate}

\subsubsection{Specific Benchmark Predictions}

\begin{enumerate}[label=\textbf{P\arabic*:},resume]
\item \textbf{Sachs network fragility.} On the Sachs protein signaling
  network \cite{sachs2005} ($p = 11$, $|E| = 17$), at least 3 edges
  have fragility score $\fragility(e) > 0.7$ for the canonical
  PKC $\to$ Raf causal query (where treatment $X = \mathrm{PKC}$,
  outcome $Y = \mathrm{Raf}$, and $\Phi = \Phi_{\mathrm{BD}}$). This
  prediction is specific, directly verifiable on a single well-studied
  network, and independent of synthetic data generation. (Exact count;
  no statistical test needed, verified by direct computation.)

  \emph{Rationale:} The Sachs network has multiple alternative pathways
  from PKC to Raf (both direct and via Mek), and the back-door criterion
  is sensitive to the presence or absence of edges on these pathways. We
  expect at least 3 of the 17 edges to be structurally load-bearing for
  this specific query, based on preliminary inspection of the DAG
  topology.

\item \textbf{Runtime--treewidth scaling.} On synthetic DAGs with
  controlled treewidth ($w \in \{2, 3, 4, 5, 6\}$, $p \in
  \{12, 16, 20\}$), the logarithm of ILP runtime scales approximately
  linearly with treewidth: the linear regression of $\log(\text{runtime})$
  on $w$ has slope in $[1.5, 3.5]$ and coefficient of determination
  $R^2 > 0.8$. Metric: slope and $R^2$ from OLS regression of
  $\log_{10}(\text{runtime in seconds})$ on $w$, stratified by $p$.
  (Linear regression with $95\%$ confidence interval for slope;
  $F$-test for $R^2 > 0$.)

  \emph{Rationale:} The FPT complexity bound $O(3^{O(w^2)} \cdot |\An|)$
  (\Cref{thm:fpt}) predicts that $\log(\text{runtime}) = O(w^2)$, but
  the ILP solver exploits structural properties (symmetry breaking,
  cutting planes, node presolving) that may reduce the effective
  dependence on $w$. We predict an approximately linear (rather than
  quadratic) relationship between $\log(\text{runtime})$ and $w$ in
  the practical regime $w \in [2, 6]$, with a slope reflecting the
  base of the exponential in the solver's branch-and-bound tree.
\end{enumerate}


\subsection{Extended Statistical Analysis Plan}\label{sec:extended_sap}

\noindent\textbf{Pre-registration of all 25 predictions.} The complete
set of 25 predictions (P1--P25) is stated \emph{before} any experiments
are conducted. The original 17 predictions (P1--P17) were formulated
based on theoretical analysis and preliminary structural observations
of published DAG sizes and treewidths. The additional 8 predictions
(P18--P25) were formulated based on the extended theoretical analysis
(sample complexity bounds, power analysis, DGP specifications, and
additional baseline comparisons) without access to experimental results.
No prediction threshold was calibrated using experimental data.

\subsubsection{Multiple Testing Correction}

With 25 simultaneous falsifiable predictions, the risk of spurious
confirmation increases substantially. We apply Holm--Bonferroni
step-down correction \cite{holm1979} at family-wise significance level
$\alpha_{\mathrm{FW}} = 0.05$ across all 25 predictions evaluated
simultaneously.

Concretely, let $p_{(1)} \leq p_{(2)} \leq \cdots \leq p_{(25)}$ be
the ordered $p$-values from the 25 individual prediction tests. The
Holm--Bonferroni procedure rejects $H_{0,(i)}$ if
$p_{(i)} < \alpha_{\mathrm{FW}} / (25 - i + 1)$ for all $j \leq i$.
This controls the family-wise error rate at $\leq 0.05$ under
arbitrary dependence among the test statistics.

\subsubsection{Effect Size Reporting}

For each prediction, we report the appropriate effect size measure in
addition to the $p$-value:
\begin{itemize}
  \item \textbf{Continuous comparisons} (e.g., AUC, runtime ratios,
    fragility scores): Cohen's $d = (\bar{X}_1 - \bar{X}_2) / s_p$,
    where $s_p$ is the pooled standard deviation. We interpret
    $|d| < 0.2$ as negligible, $0.2 \leq |d| < 0.5$ as small,
    $0.5 \leq |d| < 0.8$ as medium, and $|d| \geq 0.8$ as large.
  \item \textbf{Binary outcomes} (e.g., fraction of DAGs with
    $\radius \leq 2$, fraction of instances with gap $\leq 2$):
    odds ratio $\mathrm{OR} = [p_1 / (1-p_1)] / [p_0 / (1-p_0)]$,
    where $p_1$ is the observed success rate and $p_0$ is the null
    rate. A $95\%$ confidence interval for OR is computed using
    the exact mid-$p$ method.
  \item \textbf{Correlation coefficients} (e.g., Spearman $\rho$ between
    structural radius and E-value): report $\rho$ with a $95\%$
    bootstrap confidence interval (10{,}000 resamples with BCa
    correction).
  \item \textbf{Proportions} (e.g., precision@threshold): report the
    point estimate with an exact Clopper--Pearson $95\%$ confidence
    interval.
\end{itemize}

\subsubsection{Paired Bootstrap Comparisons}

For all pairwise comparisons between CausalCert and each baseline
(B1--B10), we use paired bootstrap tests with $10{,}000$ resamples.
The pairing is at the level of DAG instances: for each synthetic or
semi-synthetic instance $i$, we compute the metric (e.g., fragility
AUC, radius estimate, runtime) for both CausalCert and the baseline,
yielding paired differences $\Delta_i$. The test statistic is the
mean paired difference $\bar\Delta$ with a bootstrap confidence
interval.

Specifically, for each comparison:
\begin{enumerate}[label=(\roman*)]
  \item Draw $B = 10{,}000$ bootstrap samples of size $n_{\mathrm{inst}}$
    (number of instances) with replacement from the paired differences
    $\{(\Delta_1, \ldots, \Delta_{n_{\mathrm{inst}}})\}$.
  \item Compute the bootstrap distribution of $\bar\Delta^*$.
  \item Report the bootstrap percentile $95\%$ confidence interval
    $[\bar\Delta^*_{0.025}, \bar\Delta^*_{0.975}]$ and the bootstrap
    $p$-value $\hat{p} = 2 \cdot \min\bigl(\Pr(\bar\Delta^* \leq 0),
    \Pr(\bar\Delta^* \geq 0)\bigr)$.
\end{enumerate}

\noindent This procedure is applied to each CausalCert-vs-baseline
comparison for the following metrics:
\begin{itemize}
  \item Fragility AUC (CausalCert vs.\ B1, B2, B7, B8, B9, B10)
  \item Radius accuracy (CausalCert vs.\ B2, B4, B7)
  \item Runtime (CausalCert vs.\ B1, B6)
  \item Fragility ranking correlation (CausalCert vs.\ B8, B9, B10)
\end{itemize}

\subsubsection{Power Analysis for Published DAG Benchmark}

The 15 published DAGs (\Cref{tab:published_dags}) constitute a small
sample for statistical inference. We explicitly characterize the
power limitations and focus the analysis accordingly:

\begin{itemize}
  \item \textbf{Detectable effect sizes.} With $N = 15$ paired
    observations (one per published DAG), a paired $t$-test at
    $\alpha = 0.05$ has $80\%$ power to detect Cohen's $d \geq 0.78$
    (large effects only). For Spearman correlation, $80\%$ power to
    detect $|\rho| \geq 0.51$ requires $N \geq 15$; detecting
    $|\rho| \geq 0.3$ requires $N \geq 84$.
  \item \textbf{Focus on large effects.} Given the limited $N$, we
    focus predictions P10 (structural fragility prevalence), P11
    (E-value complementarity), and P23 (extended complementarity) on
    large, practically meaningful effects rather than subtle differences.
    Predictions that would require $N > 50$ to detect are flagged as
    ``exploratory'' and interpreted with appropriate caveats.
  \item \textbf{Exact tests for binary outcomes.} For predictions
    involving counts (e.g., ``$\geq 10/15$ DAGs have $\radius \leq 2$''),
    we use exact binomial tests, which are valid for any $N$ and do not
    rely on asymptotic approximations.
  \item \textbf{Supplementary bootstrap.} For all correlation-based
    predictions on the published DAGs, we supplement the permutation test
    with a BCa bootstrap confidence interval to quantify uncertainty in
    the correlation estimate, even when the formal hypothesis test has
    limited power.
\end{itemize}




% ══════════════════════════════════════════════════════════════════
% SECTION 8: RELATED WORK
% ══════════════════════════════════════════════════════════════════
\section{Related Work}\label{sec:related}

\subsection{Sensitivity Analysis for Unmeasured Confounding}

The dominant paradigm for causal robustness holds the DAG fixed and quantifies
sensitivity to unmeasured confounding.

Rosenbaum's sensitivity analysis framework \cite{rosenbaum2002} parameterizes the
degree of hidden bias via a sensitivity parameter $\Gamma$ and asks: for how large a
$\Gamma$ does the causal conclusion survive? The E-value of VanderWeele and Ding
\cite{vanderweele2017} provides a single-number summary: the minimum strength of
confounding (on the risk ratio scale) needed to explain away an observed association.
Cinelli and Hazlett's \texttt{sensemakr} \cite{cinelli2020} extends this to partial
$R^2$-based sensitivity curves.

All these methods hold the DAG structure fixed and vary the strength of unmeasured
confounders. CausalCert provides the complementary axis: it holds the confounding
structure fixed and varies the DAG topology. A causal analysis may be robust to
confounding (high E-value) but fragile to structural misspecification (low radius), or
vice versa. The two axes are logically independent, and a complete robustness assessment
should report both.

\subsection{Causal Structure Learning}

Structure learning algorithms---PC \cite{spirtes2000}, GES \cite{chickering2002},
NOTEARS \cite{zheng2018}, LiNGAM \cite{shimizu2006}, DAGMA \cite{bello2022}---aim to
\emph{discover} the DAG from data. CausalCert serves a different purpose: given an
assumed DAG (whether expert-specified or algorithm-learned), it stress-tests the
assumed DAG by quantifying structural fragility.

IDA \cite{maathuis2009} computes the range of causal effects across the Markov
equivalence class (MEC) of the estimated CPDAG. This provides MEC-level sensitivity
but does not search beyond the equivalence class. CausalCert's $k$-edit neighborhoods
span multiple MECs, providing a strictly more general search space for $k \geq 2$.

The structural intervention distance (SID) of Peters and B\"uhlmann \cite{peters2015}
measures the distance between two DAGs in terms of intervention distributions. SID is
conclusion-agnostic (measures global distance); CausalCert's radius is
conclusion-specific (measures the minimum distance to overturn a specific $\Phi$).

\subsection{DAG Testing and Validation}

Textor et~al.'s \texttt{dagitty} \cite{textor2016} tests the implied conditional
independences of a DAG against data. This can detect DAG misspecification (rejected
implications) but does not quantify which edges are load-bearing for a specific
conclusion or compute the multi-edit robustness radius.

Shah and Peters \cite{shah2020} prove the impossibility of uniformly valid nonparametric
CI testing: no single test can have both correct level and non-trivial power against all
alternatives. CausalCert addresses this by combining multiple CI tests with complementary
power profiles via the Cauchy combination test \cite{liu2020}, avoiding the need for a
single uniformly valid test.

\subsection{Computational Complexity of DAG Problems}

Chickering \cite{chickering1996} proved NP-hardness of optimal DAG structure learning
(score-based). Our NP-hardness result (\Cref{thm:np_hard}) addresses a different
problem---DAG repair subject to CI constraints and conclusion negation---using a novel
reduction from vertex cover.

The FPT algorithm (\Cref{thm:fpt}) adapts tree-decomposition-based dynamic programming
from constraint satisfaction \cite{freuder1990} to the DAG-specific setting.
Bodlaender's linear-time algorithm for bounded treewidth \cite{bodlaender1996} provides
the tree decomposition as a preprocessing step.

% ══════════════════════════════════════════════════════════════════
% SECTION 9: LIMITATIONS
% ══════════════════════════════════════════════════════════════════
\section{Limitations}\label{sec:limitations}

We honestly acknowledge eight limitations of CausalCert:

\begin{enumerate}
\item \textbf{Cannot certify DAG correctness.} CausalCert measures sensitivity, not
  ground truth. A high robustness radius does not prove the DAG is correct---it only
  shows that many simultaneous errors would be needed to overturn the conclusion.

\item \textbf{Causal sufficiency assumption.} CausalCert assumes no latent confounders
  (\Cref{asm:sufficiency}). Extension to MAGs/PAGs with latent variables is non-trivial
  because the edit space becomes larger (edge types include bidirected edges) and
  d-separation generalizes to m-separation. This is the most important direction for
  future work.

\item \textbf{Back-door criterion only.} The current implementation supports the
  back-door criterion as the primary conclusion predicate. Front-door, instrumental
  variable, and general do-calculus predicates are theoretically compatible but not
  yet implemented. The ILP constraint generation would need to be extended for each
  new predicate.

\item \textbf{Exact radius limited to $p \approx 20$.} The ILP formulation is NP-hard
  (\Cref{thm:np_hard}) and becomes impractical for $p > 20$. Fragility scoring scales
  to larger DAGs ($p \approx 50$), and the LP relaxation provides useful lower bounds
  for $p \approx 100$, but exact radii require either small $p$ or small treewidth.

\item \textbf{CI test error propagation.} All fragility scores and radii are conditional
  on the accuracy of CI test results (\Cref{asm:oracle}). CI test errors (Type~I and
  Type~II) propagate to the output. No formal error propagation bound exists for
  nonparametric CI tests \cite{shah2020}. The power diagnostics mitigate but do not
  eliminate this issue.

\item \textbf{DAGs only.} CausalCert requires an acyclic graph (\Cref{asm:acyclic}).
  Cyclic causal models (e.g., feedback loops in dynamic systems) are out of scope.
  Extension to cyclic models would require replacing d-separation with $\sigma$-separation
  or other cyclic separation criteria.

\item \textbf{Faithfulness sensitivity.} Under faithfulness violations (\Cref{asm:faithful}),
  a fragile edge may appear robust because the statistical signal is cancelled by parameter
  values. The multi-resolution CI ensemble mitigates this partially, but exact
  cancellation remains undetectable. More importantly, \emph{near-faithfulness violations}
  (where partial correlations are very small but nonzero) are common in practice---e.g.,
  when multiple causal pathways with opposing signs nearly cancel---and cause CI tests to
  have arbitrarily low power for specific edges. In such cases, fragility scores for the
  near-cancelled edge will be artificially low (appears robust when it is not). This is
  an inherent limitation of any CI-test-based approach.

\item \textbf{LP gap may be loose.} The theoretical integrality gap bound of $O(w)$
  (\Cref{thm:lp_gap}) is worst-case. Empirically, the gap is usually 0--1 for published
  DAGs, but pathological instances with larger gaps exist. The bound could potentially
  be tightened with structure-specific rounding schemes.
\end{enumerate}


% ══════════════════════════════════════════════════════════════════
% SECTION 10: CONCLUSION
% ══════════════════════════════════════════════════════════════════
\section{Implementation}\label{sec:implementation}

The CausalCert software implementation realizes the full theoretical framework
described in Sections~\ref{sec:results}--\ref{sec:algorithms} as a modular
Python library comprising 209 source files and approximately 82{,}000 lines of
code (81{,}829 non-empty lines), with 1{,}364 passing unit and integration tests.

\subsection{Architecture Overview}\label{sec:arch}

The implementation is organized into 15 subpackages, each corresponding to a
logically self-contained aspect of the CausalCert pipeline:

\begin{enumerate}
  \item \textbf{dag} (14 files) --- Core DAG representation backed by dense NumPy
    adjacency matrices, with d-separation oracle, moral graph computation,
    incremental d-separation updates (Theorem~7), MEC operations, bitset-accelerated
    d-separation using \texttt{uint64} arrays, and sparse DAG support via SciPy CSR
    matrices for graphs with $p > 50$.

  \item \textbf{ci\_testing} (17 files) --- Conditional independence testing engine
    implementing KCI with Nystr\"om approximation, HSIC with conditional
    residualization, KSG mutual information estimator, classifier-based CI test
    (CCIT), rank-based tests, and the Cauchy combination ensemble
    (Theorem~4). Adaptive test selection via meta-learning features, shared
    kernel operations with LRU caching, and CI diagnostics (power curves,
    calibration assessment).

  \item \textbf{solver} (18 files) --- The conflict-driven clause-learning search
    (Algorithm~1) with two-watched-literal propagation, clause database with
    subsumption checking and activity-based deletion, Luby/geometric/Glucose
    restart strategies, phase saving, first-UIP conflict analysis with recursive
    and local clause minimization, LRB/CHB/EVSIDS branching heuristics, and
    preprocessing (unit propagation, pure literals, blocked edit detection).
    Also includes LP relaxation for anytime lower bounds (Theorem~6), ILP
    formulation via the MIP library, and a parallel portfolio solver using
    multiprocessing.

  \item \textbf{treewidth} (9 files) --- Tree decomposition computation via
    min-degree, min-fill, min-width, and maximum cardinality search heuristics;
    conversion to nice tree decomposition (leaf/introduce/forget/join nodes);
    dynamic programming on nice tree decompositions for FPT edit distance
    computation (Theorem~3) with bitmask state representation; elimination
    ordering algorithms; and minimal separator enumeration.

  \item \textbf{fragility} (7 files) --- Per-edge fragility scoring via
    single-edit BFS, multi-channel scoring (d-separation, identification,
    sign, magnitude), bootstrap-based confidence intervals, aggregation across
    channels, and ranking.

  \item \textbf{estimation} (17 files) --- Causal effect estimation via AIPW
    with cross-fitting (Theorem~8), double/debiased machine learning (partially
    linear and interactive models), targeted minimum loss estimation (TMLE) with
    clever covariates and iterated targeting, inverse probability weighting with
    stabilized weights, partial identification bounds (Manski, Balke-Pearl),
    influence-function-based variance estimation, and diagnostics (overlap
    assessment, covariate balance, positivity checks).

  \item \textbf{simulation} (9 files) --- Data generation engines for linear
    Gaussian, nonlinear, and mixed-type structural causal models; six noise
    model families (Gaussian, Student-$t$, mixture, heteroskedastic,
    non-additive, discrete); a library of eight standard DGPs with known ground
    truth (LaLonde, smoking/birthweight, IHDP, IV, mediation, confounded,
    faithfulness violation, sparse high-dimensional); systematic perturbation
    generation; faithfulness analysis; and Monte Carlo simulation runner with
    parallel execution and coverage estimation.

  \item \textbf{formats} (9 files) --- Comprehensive format support for the two
    primary causal DAG specification formats: (i)~GraphViz DOT, with a full
    lexer and recursive-descent parser handling subgraphs, clusters, node/edge
    attributes, quoted identifiers, HTML labels, and C-style comments;
    (ii)~DAGitty format with exposure/outcome/adjusted/latent role annotations,
    position coordinates, directed/bidirected/undirected edges, and URL encoding.
    Additional formats include TETRAD XML, PCALG adjacency lists, bnlearn model
    strings, and GML. Round-trip validation and cross-format conversion are
    supported.

  \item \textbf{algebra} (6 files) --- Edit algebra with composition, inverse,
    commutativity checking, normal form computation, and orbit/closure analysis;
    edit lattice with Hasse diagram construction, meet/join operations, level
    set enumeration, and connected component analysis; and MEC-aware equivalence
    classes with canonical representative selection and quotient lattice
    construction.

  \item \textbf{evaluation} (11 files) --- Benchmark suite with standardized
    protocol for small/medium/large DAGs, paired permutation tests for
    statistical comparison against baselines (IDA, SID, brute-force), calibration
    assessment for CI tests, fragility scores, radii, and coverage; ablation
    studies; scalability profiling; and semi-synthetic benchmarks.

  \item \textbf{data} (7 files), \textbf{pipeline} (11 files),
    \textbf{reporting} (7 files), \textbf{benchmarks} (3 files) --- Data
    ingestion with type inference, pipeline orchestration with caching and
    deterministic replay, built-in profiler, and audit report generation in
    JSON, HTML, and \LaTeX{} formats.
\end{enumerate}

\subsection{Theorem-to-Module Mapping}\label{sec:theorem_map}

Each theorem in Section~\ref{sec:results} is directly realized by one or more
implementation modules:
\begin{itemize}
  \item \textbf{Theorem~1} (locality): \texttt{fragility/bfs.py},
    \texttt{dag/ancestors.py} --- ancestor restriction before fragility BFS.
  \item \textbf{Theorem~2} (NP-hardness): motivates the CDCL approach in
    \texttt{solver/cdcl.py} rather than attempting polynomial-time exact algorithms.
  \item \textbf{Theorem~3} (FPT in treewidth): \texttt{treewidth/dp.py},
    \texttt{treewidth/decomposition.py} --- DP on tree decomposition for exact
    radius computation when moral-graph treewidth $w \leq 5$.
  \item \textbf{Theorem~4} (Cauchy combination): \texttt{ci\_testing/ensemble.py}
    --- the Cauchy combination test for aggregating dependent CI test p-values.
  \item \textbf{Theorem~7} (incremental d-separation):
    \texttt{dag/incremental.py}, \texttt{dag/bitset\_dsep.py} --- $O(p \cdot
    d_{\max})$ identification of changed CI relations under single edge edits.
\end{itemize}

\subsection{Format Support and Conformance}\label{sec:formats}

CausalCert supports seven DAG specification formats, with the two most common
--- GraphViz DOT and DAGitty --- receiving comprehensive parser implementations
with full conformance test suites.

The DOT parser implements a complete lexer producing tokens for identifiers,
strings, numbers, keywords (\texttt{digraph}, \texttt{subgraph}, \texttt{strict}),
and punctuation, followed by a recursive-descent parser that handles subgraphs,
clusters, node/edge attributes, edge chains, and all standard DOT features.
Round-trip fidelity is verified: parsing a DOT string, serializing the resulting
DAG, and re-parsing yields an identical adjacency matrix.

The DAGitty parser handles the full DAGitty specification language including
role annotations (exposure, outcome, adjusted, latent/unobserved), position
coordinates for layout preservation, and edge types (directed, bidirected,
undirected). Bidirected edges are automatically expanded to latent common
cause nodes for DAG-based analysis.

\subsection{Testing and Verification}\label{sec:testing}

The test suite comprises 1{,}364 tests organized into 35 test files. Tests
are categorized as:
\begin{itemize}
  \item \textbf{Unit tests} --- individual function/class testing for each module.
  \item \textbf{Property-based tests} --- Hypothesis-generated inputs testing
    invariants (e.g., tree decomposition validity, edit algebra identities).
  \item \textbf{Conformance tests} --- round-trip parsing for all supported
    formats, cross-format equivalence.
  \item \textbf{Integration tests} --- end-to-end pipeline execution on
    synthetic and published DAGs.
  \item \textbf{Calibration tests} --- Monte Carlo verification that reported
    confidence intervals achieve nominal coverage.
\end{itemize}

\section{Conclusion}\label{sec:conclusion}

We have introduced CausalCert, the first framework for quantitative structural
stress-testing of causal DAGs. Our contributions span three areas:

\textbf{Conceptual:} We formalize the notion of structural robustness via per-edge
fragility scores and the robustness radius, providing the missing structural complement
to confounding robustness measures like E-values.

\textbf{Computational:} We prove a clean complexity picture for the Minimum-Edit DAG
Repair problem: NP-complete in general, fixed-parameter tractable in treewidth, with
an ILP formulation whose LP relaxation has bounded integrality gap.

\textbf{Practical:} The CausalCert pipeline computes fragility scores in polynomial
time and exact robustness radii for DAGs up to $p \approx 20$, covering the vast
majority of published causal DAGs in applied research.

Just as E-values transformed the qualitative debate about unmeasured confounding into
a quantitative one, CausalCert transforms the qualitative debate about DAG correctness
into a quantitative stress test. We believe that reporting structural robustness
alongside confounding robustness should become standard practice in observational
causal inference.

% ══════════════════════════════════════════════════════════════════
% BIBLIOGRAPHY
% ══════════════════════════════════════════════════════════════════
\begin{thebibliography}{99}

\bibitem{pearl2009}
J.~Pearl.
\newblock \emph{Causality: Models, Reasoning, and Inference}.
\newblock Cambridge University Press, 2nd edition, 2009.

\bibitem{spirtes2000}
P.~Spirtes, C.~Glymour, and R.~Scheines.
\newblock \emph{Causation, Prediction, and Search}.
\newblock MIT Press, 2nd edition, 2000.

\bibitem{vanderweele2017}
T.~J.~VanderWeele and P.~Ding.
\newblock Sensitivity analysis in observational research: Introducing the
  E-value.
\newblock \emph{Annals of Internal Medicine}, 167(4):268--274, 2017.

\bibitem{cinelli2020}
C.~Cinelli and C.~Hazlett.
\newblock Making sense of sensitivity: Extending omitted variable bias.
\newblock \emph{Journal of the Royal Statistical Society: Series B},
  82(1):39--67, 2020.

\bibitem{rosenbaum2002}
P.~R.~Rosenbaum.
\newblock \emph{Observational Studies}.
\newblock Springer, 2nd edition, 2002.

\bibitem{chickering2002}
D.~M.~Chickering.
\newblock Optimal structure identification with greedy search.
\newblock \emph{Journal of Machine Learning Research}, 3:507--554, 2002.

\bibitem{chickering1996}
D.~M.~Chickering.
\newblock Learning Bayesian networks is NP-complete.
\newblock In \emph{Learning from Data: AI and Statistics V}, pages 121--130,
  1996.

\bibitem{maathuis2009}
M.~H.~Maathuis, M.~Kalisch, and P.~B\"uhlmann.
\newblock Estimating high-dimensional intervention effects from observational
  data.
\newblock \emph{Annals of Statistics}, 37(6A):3133--3164, 2009.

\bibitem{peters2015}
J.~Peters and P.~B\"uhlmann.
\newblock Structural intervention distance for evaluating causal graphs.
\newblock \emph{Neural Computation}, 27(3):771--799, 2015.

\bibitem{textor2016}
J.~Textor, B.~van~der~Zander, M.~S.~Gilthorpe, M.~Li\'{s}kiewicz, and
  G.~T.~H.~Ellison.
\newblock Robust causal inference using directed acyclic graphs: the R package
  `dagitty'.
\newblock \emph{International Journal of Epidemiology}, 45(6):1887--1894, 2016.

\bibitem{zheng2018}
X.~Zheng, B.~Aragam, P.~Ravikumar, and E.~P.~Xing.
\newblock DAGs with NO TEARS: Continuous optimization for structure learning.
\newblock In \emph{NeurIPS}, pages 9472--9483, 2018.

\bibitem{shimizu2006}
S.~Shimizu, P.~O.~Hoyer, A.~Hyv\"arinen, and A.~Kerminen.
\newblock A linear non-{G}aussian acyclic model for causal discovery.
\newblock \emph{Journal of Machine Learning Research}, 7:2003--2030, 2006.

\bibitem{bello2022}
K.~Bello, B.~Aragam, and P.~Ravikumar.
\newblock DAGMA: Learning DAGs via M-matrices and a log-determinant acyclicity
  characterization.
\newblock In \emph{NeurIPS}, 2022.

\bibitem{sachs2005}
K.~Sachs, O.~Perez, D.~Pe'er, D.~A.~Lauffenburger, and G.~P.~Nolan.
\newblock Causal protein-signaling networks derived from multiparameter
  single-cell data.
\newblock \emph{Science}, 308(5721):523--529, 2005.

\bibitem{lauritzen1988}
S.~L.~Lauritzen and D.~J.~Spiegelhalter.
\newblock Local computations with probabilities on graphical structures and
  their application to expert systems.
\newblock \emph{Journal of the Royal Statistical Society: Series B},
  50(2):157--224, 1988.

\bibitem{neal1996}
D.~A.~Neal and W.~R.~Johnson.
\newblock The role of premarket factors in black-white wage differences.
\newblock \emph{Journal of Political Economy}, 104(5):869--895, 1996.

\bibitem{shen2002}
R.~Shen-Orr, R.~Milo, S.~Mangan, and U.~Alon.
\newblock Network motifs in the transcriptional regulation network of
  \emph{Escherichia coli}.
\newblock \emph{Nature Genetics}, 31(1):64--68, 2002.

\bibitem{morgan2015}
S.~L.~Morgan and C.~Winship.
\newblock \emph{Counterfactuals and Causal Inference}.
\newblock Cambridge University Press, 2nd edition, 2015.

\bibitem{shah2020}
R.~D.~Shah and J.~Peters.
\newblock The hardness of conditional independence testing and the generalised
  covariance measure.
\newblock \emph{Annals of Statistics}, 48(3):1514--1538, 2020.

\bibitem{liu2020}
Y.~Liu and J.~Xie.
\newblock Cauchy combination test: A powerful test with analytic $p$-value
  calculation under arbitrary dependency structures.
\newblock \emph{Journal of the American Statistical Association},
  115(529):393--402, 2020.

\bibitem{meek1995}
C.~Meek.
\newblock Strong completeness and faithfulness in {B}ayesian networks.
\newblock In \emph{UAI}, pages 411--418, 1995.

\bibitem{bodlaender1996}
H.~L.~Bodlaender.
\newblock A linear-time algorithm for finding tree-decompositions of small
  treewidth.
\newblock \emph{SIAM Journal on Computing}, 25(6):1305--1317, 1996.

\bibitem{freuder1990}
E.~C.~Freuder.
\newblock Complexity of $K$-tree structured constraint satisfaction problems.
\newblock In \emph{AAAI}, pages 4--9, 1990.

\bibitem{raghavan1988}
P.~Raghavan and C.~D.~Thompson.
\newblock Randomized rounding: A technique for provably good algorithms and
  algorithmic proofs.
\newblock \emph{Combinatorica}, 7(4):365--374, 1987.

\bibitem{erdos1960}
P.~Erd\H{o}s and A.~R\'enyi.
\newblock On the evolution of random graphs.
\newblock \emph{Publications of the Mathematical Institute of the Hungarian
  Academy of Sciences}, 5:17--61, 1960.

\bibitem{pittel2001}
B.~Pittel, J.~Spencer, and N.~Wormald.
\newblock Sudden emergence of a giant $k$-core in a random graph.
\newblock \emph{Journal of Combinatorial Theory, Series B}, 67(1):111--151,
  1996.

\bibitem{verma1988}
T.~S.~Verma and J.~Pearl.
\newblock Causal networks: Semantics and expressiveness.
\newblock In \emph{UAI}, pages 352--359, 1988.

\bibitem{shachter1998}
R.~D.~Shachter.
\newblock Bayes-Ball: The rational pastime (for determining irrelevance and
  requisite information in belief networks and influence diagrams).
\newblock In \emph{UAI}, pages 480--487, 1998.

\end{thebibliography}

\end{document}
